\documentclass[11pt,letterpaper,reqno]{amsart}
% BEGIN INLINE manuscript-style
\makeatletter
\NeedsTeXFormat{LaTeX2e}[2020/10/01]
% Use with \documentclass[11pt,letterpaper,reqno]{amsart}.
% This file controls presentation only. It supplies no mathematical content.

\RequirePackage[T1]{fontenc}
\RequirePackage{lmodern}
\RequirePackage{microtype}
\RequirePackage[letterpaper,margin=1in,headheight=15pt]{geometry}
\RequirePackage{amsmath,amsthm,amssymb,mathtools}
\RequirePackage{enumitem}
\RequirePackage{perpage}
\RequirePackage{longtable}
\RequirePackage{xcolor}
\definecolor{manuscriptred}{rgb}{0.6,0.1,0.1}
\definecolor{manuscriptblue}{rgb}{0.1,0.3,0.5}
\definecolor{manuscriptmagenta}{HTML}{60445C}
\RequirePackage[breaklinks=true,colorlinks=true,
  linkcolor=manuscriptred,citecolor=manuscriptblue,urlcolor=manuscriptmagenta]{hyperref}
\RequirePackage{cleveref}
\hypersetup{pdfauthor={}}

% Retain the AMS paragraph, title, heading and running-head conventions.
% Equation counters restart each section: 2.i, 2.ii, ... .
% Appendix sections keep their letters: A.i, A.ii, ... .
\numberwithin{equation}{section}
\renewcommand{\theequation}{\thesection.\roman{equation}}
% Prose references include "eq." and omit parentheses around the number.
\renewcommand{\eqref}[1]{\textup{eq.~\ref{#1}}}
\crefname{equation}{eq.}{eqs.}
\Crefname{equation}{eq.}{eqs.}
\creflabelformat{equation}{#2#1#3}
\setcounter{tocdepth}{1}
\renewcommand{\thefigure}{\Roman{figure}}
\renewcommand{\thetable}{\Roman{table}}
\newcommand{\mbneedspace}[1]{\par\begingroup
  \dimen0=\pagegoal\advance\dimen0 by-\pagetotal
  \ifdim\dimen0<#1\relax\newpage\fi\endgroup}
\MakePerPage{footnote}
\renewcommand{\thefootnote}{%
  \ifnum\value{footnote}>9\arabic{footnote}\else\fnsymbol{footnote}\fi}

% Upright statements. Main results are A, B, C; supporting statements
% share a section-local sequence. These end marks indicate typography,
% never proof completeness, validation, or a deferred-proof status.
\theoremstyle{definition}
\newtheorem{ms@theorem}{Theorem}
\renewcommand{\thems@theorem}{\Alph{ms@theorem}}
\newtheorem*{ms@theoremstar}{Theorem}
\newtheorem*{ms@theoremzero}{Theorem 0}
\newtheorem{ms@proposition}{Proposition}[section]
\newtheorem{ms@lemma}[ms@proposition]{Lemma}
\newtheorem{ms@definition}[ms@proposition]{Definition}
\newtheorem{ms@corollary}[ms@proposition]{Corollary}
\newtheorem*{ms@corollarystar}{Corollary}

\newcommand{\ms@majorstart}{%
  \pushQED{\qed}\renewcommand{\qedsymbol}{$\blacksquare\square$}}
\newcommand{\ms@minorstart}{%
  \pushQED{\qed}\renewcommand{\qedsymbol}{$\blacksquare$}}
\newenvironment{theorem}[1][]{\ms@majorstart\ms@theorem[#1]}
  {\popQED\endms@theorem}
\newenvironment{theorem*}[1][]{\ms@majorstart\ms@theoremstar[#1]}
  {\popQED\endms@theoremstar}
\newenvironment{theoremzero}[1][]{\ms@majorstart\ms@theoremzero[#1]}
  {\popQED\endms@theoremzero}
\newenvironment{proposition}[1][]{\ms@minorstart\ms@proposition[#1]}
  {\popQED\endms@proposition}
\newenvironment{lemma}[1][]{\ms@minorstart\ms@lemma[#1]}
  {\popQED\endms@lemma}
\newenvironment{definition}[1][]{\ms@minorstart\ms@definition[#1]}
  {\popQED\endms@definition}
\newenvironment{corollary}[1][]{\ms@majorstart\ms@corollary[#1]}
  {\popQED\endms@corollary}
\newenvironment{corollary*}[1][]{\ms@majorstart\ms@corollarystar[#1]}
  {\popQED\endms@corollarystar}
% The source's lemma/proposition variants are deliberately not carried over:
% choose one marker for each statement type rather than encode proof status.

\theoremstyle{remark}
\newtheorem{remark}[ms@proposition]{Remark}
\newtheorem*{remark*}{Remark}
% The ordinary amsthm proof environment retains its open square.

\crefname{ms@theorem}{theorem}{theorems}
\Crefname{ms@theorem}{Theorem}{Theorems}
\crefname{ms@proposition}{proposition}{propositions}
\Crefname{ms@proposition}{Proposition}{Propositions}
\crefname{ms@lemma}{lemma}{lemmas}
\Crefname{ms@lemma}{Lemma}{Lemmas}
\crefname{ms@definition}{definition}{definitions}
\Crefname{ms@definition}{Definition}{Definitions}
\crefname{ms@corollary}{corollary}{corollaries}
\Crefname{ms@corollary}{Corollary}{Corollaries}
\crefname{remark}{remark}{remarks}
\Crefname{remark}{Remark}{Remarks}

% Bibliography configuration for this manuscript.
\makeatother
% END INLINE manuscript-style

\title[Linear lattice PDE]{Linear lattice PDE with block constraints}
\date{Working draft, 9 November 2026}
\begin{document}
\raggedbottom
\begin{abstract}
We study scalar and gauge-field operators on separated multiscale
lattices with prescribed block averages. The estimates retain the
physical fine and coarse pairings, the actual determining rows and
the boundary convention of each operator. Scalar regularity and
weighted row lifts lead to constrained inverse and reconstruction
operators. Local carrier expansions control changes of determining
region and the admitted covariance matrix functions. Physical
fluctuation gaps use the full specified Lagrangian Hessian, including
the constraint curvature and its multiplier. The coefficient classes
are stated independently of a nonlinear construction of a background.
\end{abstract}
\maketitle
\tableofcontents

\section{Introduction}

Block constraints give an elliptic problem a mass at the scale of
each assigned block. In a gauge theory, a second constraint removes
the gauge directions. These operations interact at a seam between
two scales and at the boundary of the action region. An inverse
estimate must retain the actual rows and adjoints there. Replacing
the boundary problem by a periodic one, or a physical transpose by
an unweighted array transpose, changes the operator.

We develop the linear estimates in four stages. Scalar heat and
singular-integral estimates supply the local regularity. Energetic
row lifts and normalized Gram bounds then give the constrained
quadratic estimates. Exact Schur and current identities construct
the Green and reconstruction operators. Finally, local carrier
expansions preserve the estimates before summation and allow
conditioning and matrix functions on their declared complex domains.

The block maps and their finite geometric identities are the input
from \cite{Averaging}. Scalar covariant bounds originate in
\cite[pp.~573--574]{Bal83}. Flat and mixed-scale constrained
propagators are developed in \cite[Sections D--F]{Bal84a} and
\cite[Sections A--C]{Bal84b}. Regular-background propagators and
local expansions are treated in \cite[Sections C--E]{Bal85p}.
The derivative ranges of these results differ. In particular,
the covariant-Laplacian member of \cite[eq.~(3.42)]{Bal85p}
is excluded in that paper's Theorems 3.12 and 3.13. We state each
operator's output range separately.

The results are uniform in the mesh, depth and admitted volume.
The geometric and coefficient parameters are fixed in the stated
order. We do not construct the nonlinear critical background.
For physical fluctuation statements its complete Hessian
identification is part of the input, including both the inner
constraint multiplier and the next-fibre pivot correction.
The consequences concern the specified forms and do not establish
admission of later nonlinear substitutions or retained histories.


% Construction source: draft/geometry.tex
\section{Geometry, pairings and coefficient classes}\label{mb:geometry}

Fix an odd integer $L\ge3$, an integer $k\ge0$ and a compatible
four-dimensional torus. For the block maps, compatibility
means that every period is an integer multiple of the largest
mesh used and contains at least two sites at that mesh.
A next-lattice assertion includes its next mesh in this
requirement. Period-two axes retain their distinct stored
oriented bonds. Purely scalar assertions need only the
divisibility conditions stated in their own setting.
In units of the top lattice, put
$\eta=L^{-k}$ and $\xi_j=L^{j-k}$. The fine lattice has spacing
$\eta$. A positive bond is stored once; the value on its reverse
is determined by orientation and transport. We use the real
space $\mathfrak g$ of traceless Hermitian $2\times2$ matrices,
identified with $\mathfrak{su}(2)$ by $A\mapsto iA$.
Thus $e^{i\eta A}$ is unitary for real $A\in\mathfrak g$.
The color pairing is $\langle A,B\rangle=\tfrac12\operatorname{Tr}(AB)$,
extended bilinearly to the complexification.

Let $\Omega_0\supset\cdots\supset\Omega_k$ use the nested centred
tilings, with $\Omega_j$ a union of native level-$j$ blocks of
side $\xi_j$, and put $\Omega_{k+1}=\varnothing$. The assigned blocks
$B_y$, called pillars, partition $\Omega_0$ with a half-open
convention. A pillar assigned to level $j$ has side $\xi_y=\xi_j$.
The separation assumption is
\begin{equation}\label{mb:separation}
 \operatorname{dist}_\infty(\Omega_{j+1},\Omega_j^c)
                  >\rho\xi_j,\qquad 0\le j<k.
\end{equation}
Here $\rho$ is fixed sufficiently large. It includes the margins
needed for all stated interior boxes and complete averaging
carriers. A field row at level $j$ joins two level-$j$ sites and
has length $\xi_c=\xi_j$. Interior and outward rows are retained;
an inward bond entering $\Omega_{j+1}$ supplies no row. At the
outer endpoint of an outward row, the own assigned site level
is $j-1$. A level-$j$ endpoint average there includes the
additional transported average of those finer site values.

There are two boundary conventions. The proper problem uses the
original active graph and the action of its inner plaquettes.
The compressed problem extends an unknown by zero and restricts
the original ambient operator, including stencils crossing the
boundary. These forms can differ. Each result below states which
one it uses. On a periodic lattice all adjoints are periodic.
Interior estimates use complete stencils on the specified larger
box. They assert no regularity across an arbitrary rough exterior
boundary.

\subsection{Physical coordinates and adjoints}

Fine site and bond fields have pairing
$\langle u,v\rangle_f=\eta^4\sum_x\langle u(x),v(x)\rangle$.
The sum is over the appropriate sites or stored positive bonds.
Site-top and field-row pairings are
\begin{equation}\label{mb:physical-pairings}
 \langle b,c\rangle_s=\sum_y\xi_y^4\langle b_y,c_y\rangle,
 \qquad
 \langle B,C\rangle_r=\sum_c\xi_c^4\langle B_c,C_c\rangle.
\end{equation}
We use $T^{\mathsf T}$ for the transpose between the declared
bilinear pairings. This transpose is holomorphic in complex
coefficients. At real coefficients it is also the real Hilbert
adjoint. Positivity and operator norms use the corresponding
Hermitian inner products. A plain array transpose does not
include the factors in \eqref{mb:physical-pairings}.

For a connection $U$, the site gradient is
\begin{equation}\label{mb:gradient}
 (D_U\lambda)(b)=\eta^{-1}
       \{\lambda(b_-)-\operatorname{Ad}(U_b)\lambda(b_+)\}.
\end{equation}
The positive site Laplacian is $D_U^{\mathsf T}D_U$. Component
covariant differences of a bond field are written $\nabla_U$.
For a site-based component put
$T_\mu v(x)=\operatorname{Ad}U(x,x+\eta e_\mu)v(x+\eta e_\mu)$ and $\nabla_{\mu,U}v=(T_\mu-I)v/\eta$.
The covariant curl on a positive $\mu\nu$ plaquette is
\begin{equation}\label{mb:curl}
 (\mathrm d_Uv)_{\mu\nu}
    =\eta^{-1}\{(T_\mu-I)v_\nu-(T_\nu-I)v_\mu\}.
\end{equation}
Its transpose uses the declared plaquette set and physical
$\eta^4$ measure. In particular it is different from the
component differences $\nabla_U$.

Write $q$ for the actual transported own-pillar site average,
and $Q$ for the per-length derivative of the straightened field
average. The raw constraint derivative is $\mathsf mQ$, where
$(\mathsf mB)_c=\xi_cB_c$. Thus a raw datum is $b_c=\xi_cB_c$.
Its counting-measure multiplier $\widehat\omega$ corresponds to
the physical per-length multiplier by
\begin{equation}\label{mb:raw-duality}
 \sum_c\langle\widehat\omega_c,b_c\rangle
       =\langle\omega,B\rangle_r,
 \qquad \widehat\omega_c=\xi_c^3\omega_c.
\end{equation}
The analogous conversion is used at every constrained Hessian.
In particular, a coarse-bond current cannot be substituted for
a new-row multiplier without the derivative of the relevant
constraint chart.

At a real connection, the site average is an arithmetic mean
of adjoint transports in each disjoint pillar. Its physical
transpose satisfies $qq^{\mathsf T}=I$. Indeed the transpose
inserts the inverse transport, and the fine density
$\eta^4/\xi_y^4$ sums to one on $B_y$. The same identity holds
for the holomorphic orthogonal transports on their admitted
complex branch. The exact straightened row identity is
\begin{equation}\label{mb:extended-gradient-identity}
                       QD_U=\overline D\,q.
\end{equation}
Here $\overline D$ includes the extended endpoint average just
described. An ordinary difference of two own-top values would
give a different map at an outward row. The block-averaging
identities used in defining these maps are the supplied geometric
input; the inverse and coercivity estimates for them are proved
below.

\subsection{The layer metric and summation}

Put $\xi(x)=\xi_y$ on $B_y$ and extend it by $\eta$ outside
$\Omega_0$. Equivalently, the metric scale is $\xi_j$ in layer
$\Omega_j\setminus\Omega_{j+1}$ and is $\eta$ outside
$\Omega_1$. Define
\begin{equation}\label{mb:metric}
 d(x,z)=\inf_{\gamma:x\to z}\int_\gamma\frac{ds_\infty}{\xi(s)}.
\end{equation}
Paths are taken in the torus or in a lift to its periodic cover.
Changing the assignment of a cell face does not change the
infimum. Distances between assigned cells mean distances between
their fixed centers, with a bounded enlargement when a carrier
is used. The scale outside the active region prevents a path
from acquiring a cheaper metric by leaving that region.

\begin{lemma}[Growth and scale ratios]\label{mb:metric-growth}
For fixed $L$ and $\rho\ge1$, the assigned centers obey
\begin{equation}\label{mb:metric-count}
 \#\{z:d(y,z)\le R\}
       \le C_L(1+R)^5\exp\{8(\log L)R/\rho\},\qquad R\ge0.
\end{equation}
If $y,z$ have levels $j,i$, then
$d(y,z)\ge\rho(|i-j|-1)_+$. Consequently, for any real $p$,
\begin{equation}\label{mb:scale-ratio}
 (\xi_y/\xi_z)^p
       \le L^{|p|}\exp\{|p|(\log L)d(y,z)/\rho\}.
\end{equation}
The constants are independent of depth and compatible volume.
\end{lemma}
\begin{proof}
A path between levels $i$ and $j$ crosses every complete
intervening layer. By \eqref{mb:separation}, each such crossing
has physical length greater than $\rho\xi_r$ within layer $r$.
It therefore costs at least $\rho$ in the metric. The crossing
arcs have disjoint interiors. At most one layer is lost at an
endpoint, proving the level bound and \eqref{mb:scale-ratio}.

A path of metric length at most $R+1$ starting at level $j$
uses no scale greater than $L\xi_jL^{(R+1)/\rho}$. Its physical
length is at most $(R+1)L\xi_jL^{(R+1)/\rho}$. At any reachable
level, the smallest possible mesh is
$\xi_jL^{-1-(R+1)/\rho}$. Counting that level's grid points in
the resulting physical sup-norm ball gives
$C_L(1+R)^4L^{8R/\rho}$. There are at most
$C(1+(R+1)/\rho)$ reachable levels. Summing proves
\eqref{mb:metric-count}. On a torus, choose for each reached
center the lift reached by its path. Distinct periodic centers
give distinct lifted grid points. Counting the whole cover ball
therefore remains an upper bound.
\end{proof}

For species $s,t$ with positive amplitudes $a_s(y),a_t(y)$,
let $k_{st}(y,z)$ be the norm of $P_yKP_z$ between the normalized
cell sup norms $a_t(z)^{-1}\|\cdot\|_\infty$ and
$a_s(y)^{-1}\|\cdot\|_\infty$. Every physical integration weight
is included before this norm is taken. Define
\begin{equation}\label{mb:schur-norm}
 \|K\|_{\delta;s\leftarrow t}
 =\max\left\{\sup_y\sum_z e^{\delta d(y,z)}k_{st}(y,z),
              \sup_z\sum_y e^{\delta d(y,z)}k_{st}(y,z)\right\}.
\end{equation}
Typical amplitudes are $\xi^{-1}$ for a potential,
$\xi^{-3}$ for a force or physical row multiplier, and one
for a site gauge parameter or raw datum.

\begin{lemma}[Weighted composition]\label{mb:weighted-composition}
The norms in \eqref{mb:schur-norm} are submultiplicative for
compatible species. An entry estimate
$k_{st}(y,z)\le C(\xi_y/\xi_z)^p e^{-m d(y,z)}$ implies a
uniform bound in that norm whenever
\begin{equation}\label{mb:rate-reserve}
                    m-\delta>(8+|p|)(\log L)/\rho.
\end{equation}
\end{lemma}
\begin{proof}
Insert the cell partition between two operators. The triangle
inequality gives $k_{su}^{KL}(y,z)\le\sum_w
k_{st}^K(y,w)k_{tu}^L(w,z)$. The triangle inequality for $d$
then proves the row and column bounds for the product.
For the second assertion, sum over unit metric annuli and use
\eqref{mb:metric-count} and \eqref{mb:scale-ratio}. The resulting
series is bounded by a constant times
\[
 \sum_{n\ge0}(n+2)^5
       e^{-[m-\delta-(8+|p|)\log L/\rho]n}.
\]
It converges under \eqref{mb:rate-reserve}. The same argument
centered at $z$ proves the column estimate.
\end{proof}

A cell norm does not remove the Coulomb singularity of a fine
kernel inside a cell. A pointwise kernel estimate will therefore
be stated separately from a cell estimate. We choose geometric
separation with a strict reserve in \eqref{mb:rate-reserve} for
the finitely many scale exponents that occur below.

\subsection{Coefficient hypotheses}

The basic real class consists of unitary connections on the
declared graph whose actual averaging branches are admitted and
whose plaquette curvature satisfies
\begin{equation}\label{mb:curvature-class}
 F_U(p)=(i\eta^2)^{-1}\log U(\partial p),\qquad
        \sup_p\mathsf L(p)^2|F_U(p)|\le\kappa.
\end{equation}
Here $\mathsf L$ is the assigned scale on a complete stencil;
its values on that stencil are comparable within fixed geometric
factors. The logarithm is the branch near the identity. The
curvature ceiling $\kappa$ is small and fixed. A theorem using
stencils crossing an action boundary requires
\eqref{mb:curvature-class} on those stencils as well.

Some results require the stronger regular class. On each stated
complete box of side a fixed multiple of $\xi_y$, it admits a
real gauge $U=e^{i\eta a}$ with
$\xi_y\|a\|_\infty+\xi_y^2\|\nabla a\|_\infty\le\kappa_1$.
When a $C^{1,\alpha}$ potential is needed, it is stated explicitly.
For the strong scalar estimate we require either the explicit
coefficient condition \eqref{mb:eq-local-covariant-strong-coefficients}
or an admitted real Coulomb frame with the hypotheses of
Corollary~\ref{mb:cor-current-holder}. On sufficiently enlarged
complete real cubes, Proposition~\ref{mb:prop-coulomb-frame}
constructs the Coulomb frame from the curvature bound.
The actual Wilson current then gives the additional regularity.
An independent current parameter cannot replace the actual current
in this deduction.

For a complex tube, write $\mathbf U=e^{i\eta A'}U$ about an
admitted real center, and impose the stated bounds on
$\xi_y|A'|$ and $\xi_y^2|\nabla_UA'|$. Independent complex
currents $\mathbf J$ satisfy the separate bound
\[
                 \sup_b\mathsf L(b)^3|\mathbf J(b)|\le\gamma.
\]
Their identification with the derivative of the inner Wilson action
is an additional hypothesis.
An on-shell assertion explicitly requires that identification
on every stencil it uses.

The physical Hessian class specifies the complete Lagrangian
form, including the mean-constraint multiplier and the actual
outer pivot curvature. This is input data independent of the
existence of a nonlinear critical field. The estimates below
prove its linear consequences. They neither construct that field
nor identify a different form by a Ward identity alone.

Unless stated otherwise, constants may depend on fixed $L$,
the positive row-weight bounds and the declared collar geometry.
They do not depend on $\eta$, depth, the shapes or number of
components, or compatible volume. Constants indexed by $p$ or
$\alpha$ may also depend on those fixed exponents. A radius or
constant that depends on a localization length is identified
where it occurs.


% Construction source: outline/heat-green.tex
\section{Scalar kernels and interior regularity}
\label{mb:sec-scalar}

For this section let \(\eta>0\) be the mesh and use measure
\(\eta^4\) at each site and each positively oriented bond. Write
\[
 \partial_\mu u(x)=\frac{u(x+\eta e_\mu)-u(x)}{\eta},
 \qquad
 \partial_\mu^* u(x)=\frac{u(x-\eta e_\mu)-u(x)}{\eta},
 \qquad
 \Delta_\eta=\sum_{\mu=1}^4\partial_\mu^*\partial_\mu.
\]
Here \(D^a\) denotes any product of \(a\) of these differences and
\(d u=(\partial_\mu u)_\mu\). A convolution kernel is a density for
\(\eta^4\sum_x\). In particular, the kernel of the identity is
\(\eta^{-4}1_{\{0\}}\). We use \(|x|_1\) in kernel estimates and the
distance it induces on a periodic lattice. Changing to Euclidean
distance changes only dimension constants.

\subsection{Heat and free Green kernels}

\begin{proposition}[Free kernel bounds]\label{mb:prop-flat-heat}
Let \(p_t^\eta\) be the heat kernel on \(\eta\mathbb Z^4\), and put
\(s=t+\eta^2\). It has mass one and is nonnegative. For every fixed
nonnegative integer \(a\),
\begin{equation}\label{mb:eq-flat-heat}
 |D^a p_t^\eta(x)|\le C_a s^{-(4+a)/2}
             \exp(-c|x|_1/\sqrt s),\qquad t\ge0.
\end{equation}
For \(\mu\ge0\) define \(G_\mu^\eta=\int_0^\infty
e^{-\mu^2t}p_t^\eta\,dt\), and put \(m=\min(\mu,\eta^{-1})\). Then
\begin{equation}\label{mb:eq-flat-green}
 |D^a G_\mu^\eta(x)|\le C_a(|x|_1+\eta)^{-2-a}
                           \exp(-cm|x|_1).
\end{equation}
The convolution by \(G_\mu^\eta\) inverts \(\Delta_\eta+\mu^2\) on
compactly supported tests. For \(\mu>0\), periodization gives the
massive inverse on every compatible rectangular torus.
\end{proposition}

\begin{proof}
The generator is \(\eta^{-2}(\sum_{\mu}(T_{\eta e_\mu}
+T_{-\eta e_\mu})-8)\), where \(T_y f(x)=f(x+y)\).
Its exponential series gives positivity and preservation of constants.
The Fourier density is the integral of
\(e^{ik\cdot x}e^{-t\lambda_\eta(k)}\) over \([-\pi/\eta,\pi/\eta]^4\)
with factor \((2\pi)^{-4}\), where
\(\lambda_\eta(k)=4\eta^{-2}\sum_\mu\sin^2(\eta k_\mu/2)\).
Difference multipliers are \((e^{\pm i\eta k_\mu}-1)/\eta\).

Use the same integral with \(k\) replaced by \(k+i\sigma\), where
\(\sigma_j=c_0\operatorname{sgn}(x_j)(s^{-1/2}+m)\).
Each integrand is entire and \(2\pi/\eta\)-periodic in each variable,
including \(e^{ik\cdot x}\), so opposite sides in the contour shift cancel.
Choose \(c_0\) small so that \(\eta|\sigma_j|\le1\) and
\[
 \operatorname{Re}\lambda_\eta(k+i\sigma)
 \ge \lambda_\eta(k)-C|\sigma|^2,\qquad
 Ct|\sigma|^2\le C' +\mu^2t/4 .
\]
The first inequality follows from the real part of the cosine formula.
The second uses \(t/s\le1\) and \(m\le\mu\). The shifted difference
factors have modulus at most \(C(|k|+s^{-1/2}+m)\).
If \(t\ge\eta^2\), integration against \(e^{-ct|k|^2}\) gives
\(C_a s^{-(4+a)/2}(1+m\sqrt s)^a\). If \(t\le\eta^2\),
the cube volume and \(|k|\le2\pi/\eta\) give the same bound.
The polynomial is bounded by \(C_a e^{\mu^2t/4}\) since \(m\eta\le1\).
Consequently
\[
 e^{-\mu^2t}|D^a p_t^\eta(x)|
 \le C_a s^{-(4+a)/2}
       e^{-c|x|_1/\sqrt s-cm|x|_1-\mu^2t/2}.
\]
Taking \(\mu=0\) proves the heat estimate. To obtain the
Green estimate, integrate the power of \(s\) at \(|x|_1\le\eta\)
and substitute \(v=|x|_1/\sqrt s\) at \(|x|_1>\eta\).
The estimates justify each fixed difference under the integral.
Integrating \(-\partial_t(e^{-\mu^2t}p_t^\eta)\) gives the identity
density. For \(\mu=0\), the large-time boundary term vanishes against
any finitely supported function by the heat supremum bound.
For \(\mu>0\), both periodization and convolution are absolutely
summable. Periodization commutes with the differences and sends the
identity density to the periodic identity density.
\end{proof}

Summing the exponential bound gives, for \(1\le q\le\infty\),
\begin{equation}\label{mb:eq-heat-lq}
 \|D^a p_t^\eta\|_q\le C_{a,q}(t+\eta^2)^{-a/2-2(1-1/q)}.
\end{equation}
Indeed a ball of radius \(r\ge\eta\) has physical mass at most \(Cr^4\).
The same summation outside a ball of radius \(r\) gives an additional
factor \(C_N(1+r/\sqrt{t+\eta^2})^{-N}\) in the \(\ell^1\) estimate,
for any fixed \(N\). In particular, for \(R\ge\eta\), \(b>0\) and
\(0\le a\le3\),
\begin{equation}\label{mb:eq-massive-off-diagonal}
 \sum_{|x|_1\ge bR}\eta^4 |D^aG_{1/R}^\eta(x)|
                    \le C_{a,b}R^{2-a}.
\end{equation}
This follows directly from \eqref{mb:eq-flat-green} and shell
summation, also when \(bR<\eta\). For a periodized kernel, the same
bound holds with periodic distance, since every lift of a point
in the indicated sum lies at distance at least \(bR\).

\subsection{Finite-exponent and H\"older estimates}

\begin{proposition}[Second differences on \(L^p\)]\label{mb:prop-flat-lp}
For \(1<p<\infty\), every mesh \(\eta>0\), and every \(\mu\ge0\),
\begin{equation}\label{mb:eq-flat-lp}
 \|D^2(\Delta_\eta+\mu^2)^{-1}f\|_p\le C_p\|f\|_p.
\end{equation}
This holds on the full lattice and on every compatible rectangular
torus, with the same constant. At \(\mu=0\) the full-lattice
operator is the differentiated Green convolution on finite supports.
On a torus it is the inverse on mean-zero inputs, extended to
vanish on constants.
\end{proposition}

\begin{proof}
The Fourier multiplier has modulus at most one, since the
modulus of each first-difference multiplier is bounded by
\(\lambda_\eta^{1/2}\). This gives the \(L^2\) bound.
The physical kernel \(k=D^2G_\mu^\eta\) satisfies
\[
 |k(z)|\le C(|z|_1+\eta)^{-4},
 \qquad |Dk(z)|\le C(|z|_1+\eta)^{-5}.
\]
If \(y\ne y'\), telescope along a shortest coordinate path from
\(y\) to \(y'\). Writing \(\delta=|y-y'|_1\), every point in the
telescope remains at distance at least \(\tfrac34|x-y|_1\) from
\(x\) when \(|x-y|_1>4\delta\). Summation in shells gives
\begin{equation}\label{mb:eq-kernel-hormander}
 \eta^4\sum_{|x-y|_1>4\delta}
             |k(x-y)-k(x-y')|\le C .
\end{equation}

Use dyadic cubes of side \(2^j\eta\) in the full lattice, down to
singletons. For a finitely supported \(f\) and level \(\lambda>0\),
select the maximal cubes \(C\) with mean \(|f|>\lambda\).
Their total volume is at most \(\lambda^{-1}\|f\|_1\).
Replace \(f\) on each by its mean to obtain \(g\), and let
\(b_C=(f-\operatorname{mean}_C f)1_C\). The parent-to-child
volume ratio is \(16\), hence \(\|g\|_\infty\le16\lambda\)
and \(\|g\|_2^2\le16\lambda\|f\|_1\). Also
\(\sum_C\|b_C\|_1\le2\|f\|_1\).
Outside fixed enlargements of the selected cubes, use
\(\sum_{x\in C} b_C(x)=0\) on each cube. Inserting the kernel at a reference
point of \(C\) and applying \eqref{mb:eq-kernel-hormander} bounds
the \(L^1\) norm of the bad output there by \(C\|f\|_1\).
The \(L^2\) good bound and the volume of the enlargements give
\(\operatorname{vol}\{|Tf|>\lambda\}\le C\lambda^{-1}\|f\|_1\).

For \(1<p<2\), at level \(\lambda\) apply the weak bound to
\(f1_{\{|f|>\lambda\}}\) and the \(L^2\) bound to its complement.
Integration of the distribution function gives inner integrals
\(\int_0^{|f(x)|}\lambda^{p-2}\,d\lambda\) and
\(\int_{|f(x)|}^\infty\lambda^{p-3}\,d\lambda\), respectively,
which yield \(C_p\|f\|_p^p\). Duality gives \(p>2\) because the
adjoint multiplier has the same bounds.

At fixed mesh and positive mass the kernel is absolutely
summable. To pass to periodic data, first truncate the kernel
at a fixed radius. Apply the full-lattice estimate to periodic
data restricted to a growing rectangular union of period cells,
and discard cells within that radius of its boundary. The
proportion discarded tends to zero. The truncated operator
norm differs from the proved bound by at most the absolute
tail of the kernel. Let the radius tend to infinity. On the
finite torus the massive multipliers converge to the stated
mass-zero multiplier, including zero on the constant mode.
\end{proof}

For nested cubes separated by a distance \(bR\), with the outer
cube of radius comparable to \(R\), a sampled smooth cutoff
and \eqref{mb:eq-flat-lp} give
\begin{equation}\label{mb:eq-flat-interior-lp}
 \|D^2u\|_{p,\mathrm{in}}\le C_{p,b}\bigl(
 \|\Delta_\eta u\|_{p,\mathrm{out}}+
 R^{-1}\|Du\|_{p,\mathrm{out}}+
 R^{-2}\|u\|_{p,\mathrm{out}}\bigr).
\end{equation}
The cutoff is one on the inner stencils, and its support, enlarged
by those stencils, is inside the outer cube. The difference product
rule writes \(\Delta_\eta(\chi u)\) as \(\chi\Delta_\eta u\) and
shifted products of \(D\chi,Du\) or \(D^2\chi,u\). Their norms
give the last two terms. All formulas on nested cubes below use
this stencil convention. If a fixed ratio \(R/\eta\) is bounded,
the direct estimate by \(R^{-a}\) times the field norm on the
containing stencils gives the corresponding derivative bound.

For a one-form \(v\), put
\((dv)_{\mu\nu}=\partial_\mu v_\nu-\partial_\nu v_\mu\)
for \(\mu<\nu\). The differences commute and
\(d^*d+dd^*=\Delta_\eta\) on one-forms. Applying the differentiated
inverse to this identity proves
\begin{equation}\label{mb:eq-flat-hodge}
 \|Dv\|_p\le C_p(\|dv\|_p+\|d^*v\|_p).
\end{equation}
The identity follows componentwise by canceling the mixed
differences in the two terms. With a cutoff it gives the
interior estimate with the additional term \(C_{p,b}R^{-1}
\|v\|_{p,\mathrm{out}}\). The zero-frequency constant causes
no term on a torus, as \(D\) kills it.

For later use rescale a physical cube of radius \(R\) to unit
size. At mesh \(0<\eta\le1\), set \(R_1=(\Delta_\eta+1)^{-1}\).
Young's inequality and \eqref{mb:eq-heat-lq} give
\begin{equation}\label{mb:eq-heat-sobolev}
 \|D^a R_1 f\|_q\le C_{a,p,q}\|f\|_p,\qquad
 1\le p\le q\le\infty,\quad a/2+2(1/p-1/q)<1 .
\end{equation}
The time integral converges uniformly at zero by the strict
inequality and at infinity by the mass. In addition, for finite
\(p>4\), \(\alpha=1-4/p\), and \(0<|z|_1\le1\),
\begin{equation}\label{mb:eq-gradient-morrey}
 \|(\tau_z-1)D R_1 f\|_\infty
       \le C_p|z|_1^\alpha\|f\|_p .
\end{equation}
For \(\delta=|z|_1\), telescope the heat difference and use the
smaller of \(C(t+\eta^2)^{-1/2-2/p}\) and
\(C\delta(t+\eta^2)^{-1-2/p}\) in \(L^{p/(p-1)}\).
Integrate the first to \(\delta^2\) and the second thereafter.
Both integrals are \(O(\delta^{1-4/p})\).
Applying this to \(u=R_1u+R_1d^*du\) also gives the first-order
Morrey estimate for \(u\) in terms of its \(L^p\) norm and first
differences. These estimates apply to interior fields by cutoff.

On a cubic torus of physical side \(\ell\), with \(\ell/\eta\) an
integer, write \(G\) for the mean-zero inverse of \(\Delta_\eta\)
and \(\bar u\) for the mean of \(u\). The corresponding gradient
bounds are
\[
 \|DG f\|_\infty\le C_p\ell^{1-4/p}\|f\|_p,\qquad
 \|u-\bar u\|_\infty\le C_p\ell^{1-4/p}\|Du\|_p,\qquad p>4.
\]
To verify the first, scale to side one. For \(t\le1\) the image
sum in \eqref{mb:eq-flat-heat} and its absolute \(\ell^1\) sum give
\(\|Dp_t^{\rm per}\|_{p/(p-1)}\le C(t+\eta^2)^{-1/2-2/p}\).
For \(t\ge1\), all nonzero Fourier modes have
\(\lambda_\eta(k)\ge c|k|^2\) and their sum is \(O(e^{-ct})\).
The second bound follows from \(u-\bar u=Gd^*du\) and the same
kernel estimate. On a rectangular torus this proof is uniform
when the aspect ratios stay in a fixed compact interval.

\begin{proposition}[Discrete Schauder multiplier]
\label{mb:prop-flat-schauder}
For \(0<\alpha<1\), \(0<\eta\le1\), and any two differences,
\[
 \|D^2 R_1 f\|_{C^\alpha}\le C_\alpha\|f\|_{C^\alpha}.
\]
The norm is the supremum norm plus the seminorm on pairs at
distance at most one. The result holds on the full lattice and
on every compatible rectangular torus with its shortest periodic
distance. On fixed nested cubes with the stencil convention of
\eqref{mb:eq-flat-interior-lp},
\begin{equation}\label{mb:eq-flat-interior-schauder}
 \|D^2u\|_{C^\alpha(\mathrm{in})}\le
 C\bigl(\|\Delta_\eta u\|_{C^\alpha(\mathrm{out})}
                       +\|u\|_{C^{1,\alpha}(\mathrm{out})}\bigr).
\end{equation}
\end{proposition}

\begin{proof}
Set \(k=D^2G_1^\eta\). Its total sum is zero since it is absolutely
summable at fixed mesh and each differentiated heat kernel has
zero sum. For \(r\ge\eta\),
\[
 \left|\eta^4\sum_{|z|_1\le r} k(z)\right|\le C .
\]
Indeed split the heat integral at \(t=r^2\). For \(t<r^2\)
use cancellation and bound the sum outside the ball by
\(C_N(t+\eta^2)^{-1}(1+r/\sqrt{t+\eta^2})^{-N}\).
Integrating the bound \(C_Nr^{-N}(t+\eta^2)^{N/2-1}\) gives \(C\).
For \(t\ge r^2\) the ball volume times the supremum bound is
\(Cr^4t^{-3}\), also integrable with integral \(C\).
A ball whose centre is displaced by at most \(r/4\) has the
same bound because the discrepancy lies in an annulus of radius
comparable to \(r\) with bounded absolute kernel mass.

Subtract \(f(x)\) in the convolution for the supremum bound.
Inside radius one sum \(C[f]_\alpha |z|_1^\alpha
(|z|_1+\eta)^{-4}\) by shells. The remaining tail is integrable.
To compare \(x,x'\) at distance \(\delta\le1/16\), subtract
\(f(x)\) from both convolutions. On \(|w-x|_1\le4\delta\),
use the increment centred at each kernel centre. The extra
constant \(f(x')-f(x)\) multiplies a displaced-ball kernel sum,
which is bounded as above. This part is \(O([f]_\alpha\delta^\alpha)\).
On \(4\delta<|w-x|_1\le1\), telescoping bounds the difference of
kernels by \(C\delta |w-x|_1^{-5}\). Shell summation
against the increment gives \(C_\alpha[f]_\alpha\delta^\alpha\).
Beyond radius one the differentiated kernel has an integrable
tail and gives \(C\delta\|f\|_\infty\). The supremum estimate
covers \(1/16<\delta\le1\).

Periodic extension to the full lattice has the same supremum
and no larger seminorm. The convolution of that extension equals
the periodic multiplier by absolute summability. For the local
estimate apply the multiplier to \((\Delta_\eta+1)(\chi u)\).
The cutoff products used in \eqref{mb:eq-flat-interior-lp} have
the stated \(C^\alpha\) bounds. Products crossing the support
boundary are controlled by the same smooth cutoff differences,
with all input stencils still in the outer cube.
\end{proof}

\subsection{Poisson estimates with bounded local data}

For a cube \(B_r\) of radius \(r\) in the maximum norm, field
components are indexed by their bond tails. A source equation
on \(B_{r-\eta}\) uses the full incoming and outgoing stencils
there. Use \(\|\cdot\|_{p,B_r}\) with physical counting measure,
and use \([\cdot]_{\alpha,B_r}\) with physical distances.

\begin{proposition}[Interior Poisson estimates]\label{mb:prop-flat-poisson}
Suppose \(r\ge16\eta\) and
\(\Delta_\eta u=g+d^*F\) on \(B_{r-\eta}\), with \(u,g,F\)
given on \(B_r\). For \(0<\alpha<1\), put
\[
 N=r^{-2}\|u\|_{2,B_r}+r^2\|g\|_{\infty,B_r}
                                      +r\|F\|_{\infty,B_r}.
\]
On \(B_{r/4}\),
\begin{equation}\label{mb:eq-flat-poisson-value}
 \|u\|_\infty+r^\alpha[u]_\alpha\le C_\alpha N.
\end{equation}
If \(F=0\) there is also the bound
\begin{equation}\label{mb:eq-flat-poisson-gradient}
 r\|Du\|_\infty+r^{1+\alpha}[Du]_\alpha\le C_\alpha N.
\end{equation}
If \(F\in C^\alpha(B_r)\), then the right side of
\eqref{mb:eq-flat-poisson-gradient} is instead
\(C_\alpha(N+r^{1+\alpha}[F]_{\alpha,B_r})\).
The constants depend on the exponent and the fixed cube ratios,
uniformly in \(r\) and \(\eta\).
\end{proposition}

\begin{proof}
Use a sampled cutoff \(\chi\) equal to one on \(B_{r/2}\) and
zero outside \(B_{2r/3}\). With \(F_\mu^\chi(x)=
\chi(x+\eta e_\mu)F_\mu(x)\), the product rule gives the exact
full-lattice identity
\[
 \Delta_\eta(\chi u)
   =\chi g+d^*F^\chi
        +\sum_\mu (\partial_\mu\chi)F_\mu+[\Delta_\eta,\chi]u .
\]
The support and all shifted product stencils are in \(B_r\).
Extend the displayed products by zero and apply \(G_0^\eta\).
In the last term move the commutator onto the kernel by
summation by parts. Its adjoint is \(-[\Delta_\eta,\chi]\).
Every nonzero coefficient of that commutator is in a collar
at distance at least \(cr\) from the evaluation stencils.
For \(a=0,1,2\), the commutator applied to \(D^aG_0^\eta(x-\cdot)\)
is bounded there by \(Cr^{-4-a}\). Its terms have one kernel
difference and one cutoff difference or two cutoff differences.
It follows that the last two terms, after \(a\) output differences,
are at most
\[
 C(r^{-2-a}\|u\|_{2,B_r}+r^{1-a}\|F\|_{\infty,B_r}).
\]
The factor \(r^2\) from Cauchy--Schwarz on the collar is included
in the first term. These estimates also control the Hölder
increments of orders zero and one by telescoping an additional
difference along a coordinate path in the evaluation cube.

For the first two terms the relevant kernels are \(G_0^\eta\)
and \(DG_0^\eta\). Their absolute sums over a ball of radius
\(Cr\) are at most \(Cr^2\) and \(Cr\), respectively.
For the difference of two translates of \(DG_0^\eta\) at
separation \(\delta\le r\), split at distance \(2\delta\)
from one centre. The near sum is at most \(C\delta\) by
the first-difference size bound. On the complement telescope
and sum the second-difference bound to obtain
\(C\delta(1+\log(r/\delta))\), increasing the constant for
the fixed support radius. This is at most
\(C_\alpha r^{1-\alpha}\delta^\alpha\). For \(G_0^\eta\)
itself the same telescope gives \(Cr\delta\). Differences
with \(\delta>r\) use the supremum bounds. This proves
\eqref{mb:eq-flat-poisson-value} and
\eqref{mb:eq-flat-poisson-gradient} when \(F=0\).

For H\"older \(F\), the remaining output \(DG_0^\eta*d^*F^\chi\)
uses the same cancellation proof as the massive Schauder
multiplier. For clarity, multiply each massless second kernel
by a sampled smooth cutoff in its difference variable, equal
to one on the evaluation-to-source distances and vanishing
beyond a fixed multiple of \(r\). The truncated-sum bound in
the proof of Proposition~\ref{mb:prop-flat-schauder} used no
mass factor, so it holds for the massless kernel. The total
sum of the cut-off kernel is bounded by this bound and a shell
of bounded absolute kernel mass. After scaling by \(r\), the
near and far split of that proof applies verbatim to the
convolution with \(F^\chi-F^\chi(x)\). The additional term is
the bounded total sum times \(F^\chi(x)\). The cutoff product
has scaled \(C^\alpha\) norm at most
\(C(\|F\|_\infty+r^\alpha[F]_\alpha)\).
This proves the third assertion.
\end{proof}

These estimates also apply to a flat unitary connection on a
lifted torus cube. Trivialize by parallel transport on the
lift and use the componentwise estimate there. Norms are taken
on that lift. If every period is at least a fixed fraction of
the cube radius, a fixed multiplicity bounds the lifted \(L^p\)
norm by the norm on the projected set. To periodize
a massive kernel for such a connection, insert the commuting
unitary holonomies for the deck translations in the image
sum. The absolute tail estimates persist by unitarity. A local
estimate on an embedded cube makes no smallness assumption
on those holonomies.

\begin{lemma}[Scalar subsolution and covariant modulus]
\label{mb:lem-scalar-subsolution}
Let \(w\ge0\) and \(\Delta_\eta w\le H\) on a lifted ball
of radius \(2r\) with full stencils where tested, where \(H\ge0\)
is constant. Then its value at the centre satisfies
\begin{equation}\label{mb:eq-scalar-subsolution}
 w(x)\le C(r^{-2}\|w\|_{2,B_{2r}(x)}+r^2H).
\end{equation}
For \(r/\eta\) bounded the estimate holds by the point-value
bound on every region containing \(x\).
In particular, for a real unitary connection and a site field
\(u\), one may take \(w=|u|\) and
\(H=\|D_U^{\mathsf T}D_Uu\|_\infty\) on that ball.
\end{lemma}

\begin{proof}
When \(r\ge16\eta\), test the inequality against
\(\chi G_0^\eta(x-\cdot)\), where \(\chi\ge0\) is one on the
inner ball and has its cutoff collar inside the outer ball.
Moving the Laplacian gives \(w(x)\) and a commutator supported
on the collar. The positive test has sum at most \(Cr^2\).
The commutator kernel is bounded by \(Cr^{-4}\) by the free
Green and first-difference bounds. Its \(L^2\) norm on the
collar is \(Cr^{-2}\). This proves \eqref{mb:eq-scalar-subsolution}.
If \(r/\eta\) is bounded, use
\(w(x)\le\eta^{-2}\|w\|_{2,B_{2r}(x)}\).

For \(u(x)\ne0\), take the real inner product of each term
\(u(x)-\operatorname{Ad}(U_{xy})u(y)\) in the covariant
Laplacian with \(u(x)/|u(x)|\). Unitarity bounds it below
by \(|u(x)|-|u(y)|\). Sum over nearest neighbours, with
coefficient \(\eta^{-2}\). At a zero of \(u\), the scalar
left-hand side is nonpositive. Thus
\(\Delta_\eta|u|\le |D_U^{\mathsf T}D_Uu|\) at every site
with the full stencil.
\end{proof}


% Construction source: outline/block-gap.tex
\subsection{Scalar block gap and block resolvent}
\label{mb:sec-scalar-block}

Fix a physical block side \(R=m\eta\) with \(m\) a positive
integer. Partition a rectangular torus with periods in \(R\mathbb Z\)
into translates of \(\{0,\eta,\ldots,(m-1)\eta\}^4\).
Write \(q_{\rm ar}\) for the arithmetic site average on this disjoint
partition. This is the flat, one-size version of the site row
\(q\). Coarse scalar products have weight \(R^4\) per block.
In this subsection the star denotes the Hilbert-space adjoint.
All the displayed real operators have the same bilinear transpose.

\begin{proposition}[Mean, gap, and projection]\label{mb:prop-scalar-block-gap}
The adjoint \(q_{\rm ar}^*\) embeds a coarse function as a block-constant
fine function. Thus \(q_{\rm ar}q_{\rm ar}^*=1\) and \(P=q_{\rm ar}^*q_{\rm ar}\) is an
orthogonal projection. If \(m\ge2\), the Neumann difference
\(d_{\rm N}\), using only bonds internal to the individual blocks,
satisfies
\begin{equation}\label{mb:eq-block-poincare}
 \|(1-P)u\|_2^2\le
 \frac{\eta^2}{4\sin^2(\pi/(2m))}\|d_{\rm N}u\|_2^2
                    \le \frac{R^2}{4}\|d_{\rm N}u\|_2^2 .
\end{equation}
For \(m=1\) the left side vanishes. In particular
\(A=\Delta_\eta+R^{-2}P\ge R^{-2}\) on the torus.
Writing \(G=A^{-1}\), the coarse Gram matrix and fine projection
satisfy
\begin{equation}\label{mb:eq-scalar-block-gram}
 q_{\rm ar}G^2q_{\rm ar}^*\ge cR^4,\qquad
 \Pi_{\Delta\ker q_{\rm ar}}
   =1-Gq_{\rm ar}^*(q_{\rm ar}G^2q_{\rm ar}^*)^{-1}q_{\rm ar}G ,
\end{equation}
with a numerical \(c>0\).
\end{proposition}

\begin{proof}
Since \(R^4m^{-4}=\eta^4\), moving the sum in
\(\langle q_{\rm ar} f,u\rangle\) gives exactly the block-constant
embedding. The projection identities follow.
On a one-dimensional path of \(m\) sites, the vectors
\(\cos(\pi k(j+\tfrac12)/m)\), \(0\le k<m\), diagonalize the
Neumann difference form. Substitution in the second difference
gives eigenvalues \(4\eta^{-2}\sin^2(\pi k/(2m))\). At the ends,
reflect the cosine across the adjacent half-site to remove
the missing edge. The eigenvalues are distinct, so the vectors
form a basis. Tensoring in four directions adds the four
eigenvalues and leaves just the constant vector at zero.
The first bound in \eqref{mb:eq-block-poincare} follows.
The second uses \(\sin(\pi/(2m))\ge1/m\).
Adding the interblock bonds only increases the form, and
\(\|u\|_2^2=\|(1-P)u\|_2^2+\|q_{\rm ar}u\|_2^2\), giving the gap.

For the Gram bound choose a right inverse \(T\) for \(q_{\rm ar}\).
For \(m\ge2\), let \(Tu\) on block \(B\) be \(u(B)\) times
\[
 \phi(j)=16\prod_{\nu=1}^4
                     \sin^2\!\bigl(\pi(j_\nu+\tfrac12)/m\bigr).
\]
Each one-dimensional sine square has mean \(1/2\) by summing
the corresponding roots of unity. Hence \(q_{\rm ar}T=1\) and
\(\|T\|\le16\). The second derivative of a sine square on
a continuous interval of length \(R\) is bounded by \(C/R^2\).
At a block face its endpoint sample and the next interior sample
are \(O((\eta/R)^2)\), so the second difference across the face
is bounded by \(CR^{-2}(|u(B)|+|u(B')|)\) for the adjacent block
\(B'\). Summing the pointwise bounds with the physical measure
gives \(\|\Delta_\eta Tu\|_2\le CR^{-2}\|u\|_2\).
For \(m=1\), take \(T=1\) and use \(\|\Delta_\eta\|\le16R^{-2}\).
Thus \(\|ATu\|_2\le CR^{-2}\|u\|_2\) in both cases. Now
\[
 \|u\|_2^2=\langle q_{\rm ar}^*u,Tu\rangle
                  =\langle Gq_{\rm ar}^*u,ATu\rangle
                  \le CR^{-2}\|Gq_{\rm ar}^*u\|_2\|u\|_2 .
\]
This is the first part of \eqref{mb:eq-scalar-block-gram}.
The orthogonal complement of \(\operatorname{ran}(Gq_{\rm ar}^*)\)
is \(\ker(q_{\rm ar}G)\), which equals \(\Delta_\eta\ker q_{\rm ar}\).
Indeed \(A=\Delta_\eta\) on \(\ker q_{\rm ar}\), and conversely for
\(f\in\ker(q_{\rm ar}G)\) put \(u=Gf\). The standard range projection
for the injective map \(Gq_{\rm ar}^*\) now gives the displayed formula.
\end{proof}

In Euclidean coordinate inner products the unitary conversions
are multiplication by \(\eta^2\) on fine sites and bonds and by
\(R^2\) on coarse sites. The conjugated mean has entries \(m^{-2}\)
on its block, with squared row norm one. Keeping entries
\(m^{-4}\) in Euclidean coordinates gives squared row norm
\(m^{-4}\). Then \(q_{\rm ar}^*q_{\rm ar}\) fails to fix constants
when \(m>1\). These unitaries commute with adjoints, orthogonal
projections, and positive or Moore--Penrose inverses by conjugation.
The same conjugation applies to a covariance and its constrained
covariance, since the constraints act between the corresponding
Hilbert spaces. A combinatorial difference is \(\eta d\),
so its form in \eqref{mb:eq-block-poincare} carries the
additional factor \(\eta^{-2}\).

\begin{proposition}[Off-diagonal scalar block inverse]
\label{mb:prop-scalar-block-decay}
Let \(S\) be a set of sites of the torus, \(J_S\) extension
by zero, and \(A_S=J_S^*AJ_S\). This includes the whole torus.
For nonempty site sets \(X,Y\subset S\),
\begin{equation}\label{mb:eq-scalar-block-ct}
 \|1_X A_S^{-1}1_Y\|_{2\to2}
          \le 2R^2e^{-c\,\operatorname{dist}_\eta(X,Y)/R}.
\end{equation}
The same estimate for \(dJ_S A_S^{-1}\), with an output
bond restriction indexed by its tails, has prefactor \(CR\).
For the adjoint estimate the prefactor is also \(CR\), and for
\(dJ_S A_S^{-1}J_S^*d^*\) it is \(C\).
The distance is shortest-path distance in the ambient torus,
with length \(\eta\) per bond. The constants are independent
of \(m,R,S\), and the torus periods.
\end{proposition}

\begin{proof}
Let \(w\) be real with \(|w(x)-w(y)|\le\eta/R\) across each
bond, and write \(M_t=e^{tw}\). For \(|t|\le1\) we prove
\[
 \|A_S^{-1/2}(M_tA_SM_t^{-1}-A_S)A_S^{-1/2}\|
                                      \le C(|t|+t^2).
\]
It suffices to estimate the full-torus form on zero extensions.
The compression commutes with each multiplication, and its
energy is exactly the energy of the zero extension. For a
bond from \(x\) to \(y\), put \(\delta=w(x)-w(y)\) and
\(a=e^{t\delta}-1,\ b=e^{-t\delta}-1\).
The Laplacian conjugation error has form contribution
\(-\eta^2(a\overline{u_x}v_y+b\overline{u_y}v_x)\).
Split the parenthesis as
\[
 a(\overline{u_x}-\overline{u_y})v_y
 +b\overline{u_y}(v_x-v_y)+(a+b)\overline{u_y}v_y .
\]
Here \(|a|+|b|\le C|t|\eta/R\) and
\(|a+b|\le Ct^2\eta^2/R^2\). Summing and applying
Cauchy--Schwarz bounds the error by
\[
 C|t|R^{-1}(\|du\|_2\|v\|_2+\|u\|_2\|dv\|_2)
                          +Ct^2R^{-2}\|u\|_2\|v\|_2 .
\]
For the block term, \(P\) has coordinate entries \(m^{-4}\)
on each block, and the oscillation of \(w\) on one block is
at most four. The absolute row and column sums of the conjugation
error for \(R^{-2}P\) are therefore at most \(C|t|R^{-2}\).
The gap and the gradient energy bound give the asserted relative
estimate. This uses the full block mean in the compression,
including at blocks cut by the boundary.

Choose a numerical \(t_0>0\) so the relative error has norm
at most \(1/2\) for \(|t|\le t_0\). Factoring the conjugated
operator between \(A_S^{1/2}\) on both sides gives the inverse
by a Neumann series, with middle factor of norm at most two.
This proves \(\|M_tA_S^{-1}M_t^{-1}\|\le2R^2\).
On bonds let the multiplier use the tail value of \(w\).
Then \(M_t d M_t^{-1}=d+\mathcal E_t\) with
\(\|\mathcal E_t\|_{2\to2}\le C|t|/R\), by the same
nearest-neighbour bound on \(w\). Thus the conjugated gradient
times \(J_S A_S^{-1/2}\) has norm at most \(C\), using the
energy of the zero extension. The input divergence follows
by taking the adjoint of this bound with parameter \(-t\).
Combining these factors with the middle inverse proves the
gradient prefactors.

Finally take \(w(x)=\operatorname{dist}_\eta(x,Y)/R\), using
the tails for a bond input. The weight is zero on \(Y\)
and at least \(\operatorname{dist}_\eta(X,Y)/R\) on \(X\).
Undoing the conjugation proves all the estimates.
\end{proof}

The Dirichlet operator here compresses \(P\) itself. Averaging
over only the surviving sites would change the operator.
The off-diagonal statements are norms on weighted Hilbert
spaces. A pointwise estimate for a bounded cell source uses
the separate local argument in Lemma~\ref{mb:lem-scalar-subsolution}.


% Construction source: outline/covariant.tex
\subsection{Local covariant scalar equations}
\label{mb:sec-local-covariant}

On a lifted cube of radius comparable to \(R\ge\eta\), choose a
frame in which
\(U(x,x+\eta e_\mu)=\exp(i\eta A_\mu(x))\).
The frame may be real, with \(A_\mu\) traceless Hermitian, or complex,
with \(A_\mu\) traceless complex. The coefficient conditions
used for finite-\(p\) regularity are
\begin{equation}\label{mb:eq-local-covariant-coefficients}
 R\|A\|_\infty\le K,\qquad
 R^{2-4/q}\|DA\|_q\le K_q\quad(1<q<\infty).
\end{equation}
All differences in these conditions are on the prescribed outer
stencils. Only finitely many \(q\) are used for each conclusion.
For the second-difference Hölder estimate we use the stronger
condition, in the same frame,
\begin{equation}\label{mb:eq-local-covariant-strong-coefficients}
 R\|A\|_\infty+R^2\|DA\|_\infty
                    +R^{2+\alpha}[DA]_\alpha\le K_\alpha .
\end{equation}
On a fixed cube this also controls the scaled \(C^\alpha\)
norm of \(A\), by telescoping coordinate paths.

For example, suppose the prescribed frame satisfies \(d^*A=0\),
\(R\|A\|_\infty\le C\kappa\), and each plaquette logarithm satisfies
\(\|(i\eta^2)^{-1}\log U(\partial p)\|\le\kappa/R^2\).
For small \(\kappa\), multiplying the four exponentials gives
\[
 (i\eta^2)^{-1}\log U(\partial p)=(dA)(p)+O(\|A\|_\infty^2).
\]
The linear term is the oriented sum of the link exponents,
and the analytic remainder before division by \(\eta^2\)
is quadratic in those exponents. The interior Hodge estimate
\eqref{mb:eq-flat-hodge}, applied on a slightly smaller cube,
gives \(K_q\le C_q\kappa\). This deduction uses only the
admitted frame, its Coulomb condition and the curvature bound.
For a complex logarithmic perturbation, the corresponding
norms in \eqref{mb:eq-local-covariant-coefficients} add to
those of the real logarithm if the perturbation is specified
in the same logarithmic coordinates. No unitary estimate
will be used for such a complex coefficient.

With \(D_Uu(x,\mu)=\eta^{-1}
(u(x)-\operatorname{Ad}U_\mu(x)u(x+\eta e_\mu))\),
the bilinear transpose satisfies
\(\Delta_U=D_U^{\mathsf T}D_U\). Its site formula is
\begin{equation}\label{mb:eq-covariant-site-laplacian}
 (\Delta_Uu)(x)=\eta^{-2}\sum_\mu
 \bigl(2u(x)-\operatorname{Ad}U_\mu(x)u(x+\eta e_\mu)
       -\operatorname{Ad}U_\mu(x-\eta e_\mu)^{-1}
                                    u(x-\eta e_\mu)\bigr).
\end{equation}
Indeed \((\operatorname{Ad}g)^{\mathsf T}=\operatorname{Ad}g^{-1}\)
for the invariant bilinear form. This identity remains valid
at complex \(g\).

\begin{lemma}[Exact lower-order terms]\label{mb:lem-covariant-lower}
On the rescaled unit cube, with mesh \(h=\eta/R\) and potential
\(RA\), the difference \(\Delta_U-\Delta_h\) is the sum of
first differences with bounded coefficients and a multiplication
term \(b\) with bounded \(L^q\) norms. Its bounds are uniform
in \(h\), with constants depending only on \(K,K_q\).
In the class \eqref{mb:eq-local-covariant-strong-coefficients}
these coefficients are uniformly bounded in \(C^\alpha\).
The difference \(D_U^{\mathsf T}+d^*\) is a zeroth-order sum
with bounded coefficients and one-step shifts.
\end{lemma}

\begin{proof}
In unit coordinates, denote the rescaled potential again by
\(A\), and put \(E_\pm(a)=(e^{\pm ih\operatorname{ad}a}-1)/h\).
In \eqref{mb:eq-covariant-site-laplacian}, replace each translated
\(u\) by its central value plus its signed \(h\)-difference.
The first-difference coefficients are \(-E_+(A_\mu(x))\)
and \(-E_-(A_\mu(x-he_\mu))\) for the forward and the adjoint
backward difference. The multiplication coefficient is
\[
 b(x)=-\sum_\mu h^{-1}
       \bigl(E_+(A_\mu(x))+E_-(A_\mu(x-he_\mu))\bigr).
\]
Since \(E_+(a)+E_-(a)=O_K(h|a|^2)\), this is a sum of
bounded quadratic terms and backward divided differences
of \(E_+(A)\). The latter are bounded by \(C_K|DA|\)
by the integral differential of the matrix exponential.
In \(C^\alpha\) the same identity bounds those terms by
products of a smooth function of the two potential values
and one difference of \(A\). There is no second difference
of \(A\) in \(b\). Taking the transpose of the corresponding
first-order formula for \(D_U\) proves the last assertion.
\end{proof}

\begin{lemma}[A left change of the connection]
\label{mb:lem-site-connection-change}
Write \(U_E(x,\mu)=e^{i\eta E_\mu(x)}U(x,\mu)\).
For each direction let \(T_\mu u(x)=\operatorname{Ad}U_\mu(x)
u(x+\eta e_\mu)\), and put
\(\widetilde E_\mu(x)=\operatorname{Ad}U_\mu(x-\eta e_\mu)^{-1}
E_\mu(x-\eta e_\mu)\).
If \(\eta(\|E\|_\infty+\|\widetilde E\|_\infty)\) is bounded,
the difference \(\Delta_{U_E}-\Delta_U\) has no second
differences of its argument. The first-difference coefficients
relative to the forward and backward transports of \(U\) have
size \(C(|E|+|\widetilde E|)\). Its multiplication coefficients
are bounded by
\[
 C\sum_\mu\bigl(\eta^{-1}|E_\mu-\widetilde E_\mu|
                           +|E_\mu|^2+|\widetilde E_\mu|^2\bigr).
\]
This is a bilinear identity and norm bound in real or complex
frames, on the stencils where both operators are defined.
\end{lemma}

\begin{proof}
The shifted operator for \(U_E\) is \(e^{i\eta\operatorname{ad}E_\mu}T_\mu\),
and its inverse is \(e^{-i\eta\operatorname{ad}\widetilde E_\mu}T_\mu^{-1}\).
Thus the difference of the two positive Laplacians in direction
\(\mu\) is \(-\eta^{-2}(M_+T_\mu+M_-T_\mu^{-1})\), where
\(M_+=e^{i\eta\operatorname{ad}E_\mu}-1\) and
\(M_-=e^{-i\eta\operatorname{ad}\widetilde E_\mu}-1\).
Write \(T_\mu^{\pm1}u=u+(T_\mu^{\pm1}-1)u\).
The latter terms are a single transported difference times
\(M_\pm/\eta\). In the zeroth term \(M_++M_-\), expand both
exponentials to linear order with a bounded quadratic remainder.
This gives the asserted coefficient and its precise
\(\eta^{-1}(E_\mu-\widetilde E_\mu)\) derivative.
\end{proof}

In the same frame, a relative left perturbation also preserves
the stated logarithmic coefficient classes. In unit coordinates
on a common small exponential chart,
\[
 (ih)^{-1}\log(e^{ihE}e^{ihA})=A+E+h\,\mathcal B_h(A,E).
\]
The last function is analytic and has uniform derivatives of
each fixed order on smaller balls. Indeed the logarithm of
the product minus its linear term has a zero of order two
in \(h\). Division by \(h^2\) extends it at zero. The
remainder vanishes when either of the two potentials is zero.
Consequently its \(A\)-derivative is \(O(|E|)\), and its
\(E\)-derivative is \(O(|A|)\). A discrete first difference of
this identity uses the integral of these derivatives along
the segment joining two neighbouring input pairs. This gives
the finite-\(p\) first-difference bounds in terms of
\(\|DE\|_p+\|E\|_\infty\|DA\|_p\), with the zeroth bound
\(C\|E\|_\infty\). Products and the same integral formula
give the \(C^{1,\alpha}\) bound on the logarithm in a
\(C^{1,\alpha}\) tube. A covariant difference of \(E\)
in the fixed frame differs from \(DE\) by an \(O(AE)\)
term, with the same scaled \(C^\alpha\) estimate in
the stronger class.

\begin{proposition}[Local site bootstraps]\label{mb:prop-local-scalar}
Work on nested lifted cubes after scaling their fixed physical
radius \(R\) to one, and write \(h=\eta/R\). The norms in the
statement are in these rescaled coordinates, including the
first differences and the physical measure \(h^4\).
Assume \eqref{mb:eq-local-covariant-coefficients} and an
outer energy bound
\[
 \|u\|_{2,\mathrm{out}}+\|D_Uu\|_{2,\mathrm{out}}\le B .
\]
One may add a zeroth-order operator \(M\) to \(\Delta_U\),
provided that every cutoff step has a reserved outer box
containing the full input carrier of \(M\) at all source sites
in the cutoff support, and
\(\|Mu\|_{q,\mathrm{supp}\,\chi}\le K'_q\|u\|_{q,\mathrm{out}}\)
on these pairs of stocks for the finite \(q\) used below.

If \((\Delta_U+M)u=g\), \(4<p<\infty\), and
\(0<\alpha\le1-4/p\), then on the innermost cube
\begin{equation}\label{mb:eq-local-covariant-point}
 \|u\|_\infty+\|Du\|_\infty+[Du]_\alpha
                              \le C(B+\|g\|_{p,\mathrm{out}}).
\end{equation}
If \((\Delta_U+M)u=g_0+D_U^{\mathsf T}F\) and \(2<p<\infty\),
then
\begin{equation}\label{mb:eq-local-covariant-div}
 \|u\|_p+\|Du\|_p
                  \le C(B+\|g_0\|_{p,\mathrm{out}}
                                +\|F\|_{p,\mathrm{out}}).
\end{equation}
For \(p>4\) the left side can also include \(\|u\|_\infty\).
The constants depend on the exponents, coefficient bounds,
and the fixed cutoff and carrier margins. They are independent
of \(h\) and of the outside domain. After reversing the dilation,
\(B=R^{-2}\|u\|_{2,\mathrm{out}}
+R^{-1}\|D_Uu\|_{2,\mathrm{out}}\).
The two source norms acquire factors \(R^{2-4/p}\) for \(g,g_0\)
and \(R^{1-4/p}\) for \(F\). The left side uses scaled derivatives,
for example \(R^{1+\alpha}[Du]_\alpha\) in
\eqref{mb:eq-local-covariant-point} and
\(R^{-4/p}\|u\|_p+R^{1-4/p}\|Du\|_p\) in
\eqref{mb:eq-local-covariant-div}.
\end{proposition}

\begin{proof}
The comparison \(D_U=-d+O_K(1)\) first bounds the flat
outer \(H^1\) norm. A cutoff and the identity
\(v=R_1v+R_1d^*dv\) give \(u\in L^3\) by
\eqref{mb:eq-heat-sobolev}, since
\(\tfrac12+2(\tfrac12-\tfrac13)<1\).
In the first equation, the new flat zeroth term \(bu\)
is in \(L^2\) by \(b\in L^6\). First-order terms are in \(L^2\).
For the cutoff equation \((\Delta_h+1)v=s\), the cutoff
commutator uses only \(Du,u\) on the preceding box.
With \(s\in L^2\) we obtain \(D^2u\in L^2\) from
\eqref{mb:eq-flat-lp}, \(Du\in L^3\) and \(u\in L^{12}\)
from \eqref{mb:eq-heat-sobolev}.
At the next cutoff, \(b\in L^4\) and \(u\in L^{12}\)
give \(bu\in L^3\). The first-order terms are now in
\(L^3\), and \(M\) reads \(u\in L^3\) on its reserved
box. An \(L^3\) source yields bounded \(u\) and
\(Du\in L^q\) for every \(q<12\).
Take \(q=\min(p,8)\) in this step. If \(p>8\), one more
\(L^8\) step gives bounded \(Du\).
Thus the final cutoff equation has source in \(L^p\),
and \eqref{mb:eq-gradient-morrey} proves
\eqref{mb:eq-local-covariant-point}.

For a divergence source, each cutoff equation is the sum
of a flat divergence of a bounded cutoff times \(F\) and
an ordinary source. The latter contains the previous
commutator terms and only zeroth-order products with \(F\),
by Lemma~\ref{mb:lem-covariant-lower} and the product rule
used for \eqref{mb:eq-flat-poisson-value}.
At each step the divergence contribution to the gradient
is bounded in \(L^p\) by \eqref{mb:eq-flat-lp}.
The first ordinary source is in \(L^2\) as before, and
its gradient is in \(L^3\). If \(p\le3\) this gives
\eqref{mb:eq-local-covariant-div} already, since the field
is in \(L^3\) from the energy step.
If \(p>3\), \(R_1d^*\) maps \(L^p\) to \(L^{12}\)
by \eqref{mb:eq-heat-sobolev} (or first by restriction
to a lower exponent if \(p>12\)).
Thus the field is in \(L^{12}\), including its divergence
contribution, and the next ordinary source is in \(L^3\).
It gives a bounded field and gradient in every \(L^q\)
with \(q<12\), sufficient for \(q=\min(p,8)\).
If \(p>8\), its next \(L^8\) source gives bounded
gradient, while the divergence contribution remains
in \(L^p\). For \(p>4\) the divergence contribution
to the field is also bounded because
\(R_1d^*:L^p\to L^\infty\).
This proves the claims without estimating a difference of \(F\).
Each invocation reads the preceding box only, with
the stated margins. For bounded \(R/\eta\) the desired
suprema, differences and seminorms follow directly
from the energy on the containing stencils.
\end{proof}

If \(g_0\) is bounded and \(F\in C^\alpha\), \(0<\alpha<1\),
the conclusion \eqref{mb:eq-local-covariant-point} holds
for the divergence equation with the right side replaced by
\(C(B+\|g_0\|_\infty+\|F\|_{C^\alpha})\) in unit coordinates.
To see this, use \eqref{mb:eq-local-covariant-div} for
some \(p>4/(1-\alpha)\) on an intermediate box. In the final
cutoff equation the ordinary source, including \(bu\) and
the first-order terms, is in \(L^p\). Its contribution is
controlled by \eqref{mb:eq-gradient-morrey}. For the
divergence part use the multiplier in
Proposition~\ref{mb:prop-flat-schauder} on the cut-off \(F\).
The same argument with \(F=0\) gives the corresponding
bounded ordinary forcing estimate.

There is also a bounded-data version without an a priori
energy bound. For \(\Delta_Uu=g+D_U^{\mathsf T}F\), an outer
supremum bound on \(u\) gives an interior energy bound by
\(C(\|u\|_\infty+\|g\|_2+\|F\|_2)\) in unit coordinates.
Indeed write the equation in the flat form of
Lemma~\ref{mb:lem-covariant-lower} and test against \(\chi^2u\)
in the Hermitian pairing. For the flat term, on a bond from
\(x\) to \(y\), the real part of
\((\overline{u_y}-\overline{u_x})(\chi_y^2u_y-\chi_x^2u_x)\)
equals
\[
 |\chi_yu_y-\chi_xu_x|^2
                     -(\chi_y-\chi_x)^2\operatorname{Re}(\overline{u_y}u_x).
\]
This gives the squared gradient of \(\chi u\) up to
\(C\|u\|_\infty^2\).
The first-order and divergence terms are bounded by a
small multiple of that gradient square plus
\(C(\|u\|_\infty^2+\|F\|_2^2)\), by the product rule
and Cauchy--Schwarz. For a shifted first difference,
choose the cutoff on the same bond endpoint as the source.
The difference from \(d(\chi u)\) is a bounded cutoff derivative
times \(u\). The zeroth term is bounded by
\(\|b\|_1\|u\|_\infty^2\), and the ordinary forcing by
Cauchy--Schwarz. This proves the energy bound for real or
complex coefficients. It extends to an added \(M\) with
the stated \(L^2\) carrier bound by testing its output on
the complete outer box. After one more cutoff, the bounded
force case of \eqref{mb:eq-local-covariant-point} yields
the gradient and every fixed H\"older exponent below one.

\begin{corollary}[Second covariant differences]
\label{mb:cor-local-covariant-schauder}
In the class \eqref{mb:eq-local-covariant-strong-coefficients},
if \(\Delta_Uu=g\), then on fixed nested unit cubes
\begin{equation}\label{mb:eq-local-covariant-schauder}
 \|D_U^2u\|_{C^\alpha,\mathrm{in}}\le
 C\bigl(\|g\|_{C^\alpha,\mathrm{out}}
                      +\|u\|_{C^{1,\alpha},\mathrm{out}}\bigr).
\end{equation}
Here a covariant second difference is a composition of
forward or backward transported differences on site-based
components. The source seminorm includes pairs across adjacent
pillar faces whenever the cube crosses those faces. For example,
if the source jumps at a face, the first-difference estimate
from bounded forcing still applies, while the source seminorm
in \eqref{mb:eq-local-covariant-schauder} then includes the jump.
In physical coordinates, set
\(\|f\|_{C^\alpha_R}=\|f\|_\infty+R^\alpha[f]_\alpha\).
Then \eqref{mb:eq-local-covariant-schauder} reads
\[
 R^2\|D_U^2u\|_{C^\alpha_R,\mathrm{in}}\le C\bigl(
 R^2\|g\|_{C^\alpha_R,\mathrm{out}}+\|u\|_{\infty,\mathrm{out}}
                         +R\|Du\|_{C^\alpha_R,\mathrm{out}}\bigr).
\]
\end{corollary}

\begin{proof}
Lemma~\ref{mb:lem-covariant-lower} bounds the lower-order
terms in \(C^\alpha\), so \eqref{mb:eq-flat-interior-schauder}
applies to the flat equation. Composing two covariant
differences and subtracting the corresponding ordinary
ones leaves products of \(A,DA,u,Du\) with bounded
smooth coefficients, by the same divided-difference identity.
Their \(C^\alpha\) norms have the stated bound.
\end{proof}

\begin{corollary}[Checking the strong coefficient condition in a Coulomb frame]
\label{mb:cor-current-holder}
Suppose a real Coulomb frame \(U=e^{i\eta A}\) is already given on
the outer stock, with \(d^*A=0\),
\(R\|A\|_\infty\le C_0\kappa\), and
\[
 R^2\|F\|_\infty\le\kappa,\qquad
 R^3\|J_{\rm left}(U)\|_\infty\le j .
\]
Here \(F(p)=(i\eta^2)^{-1}\log U(\partial p)\) on the small
branch, and \(J_{\rm left}\) represents the variation
\(U\mapsto e^{it\eta B}U\) of the Wilson action
\(\sum_p(1-\tfrac12\operatorname{tr}U(\partial p))\)
in the physical link pairing. Use complete current stencils
throughout this stock. For \(\kappa\) below a fixed ceiling,
on a smaller stock and for every \(0<\alpha<1\),
\begin{equation}\label{mb:eq-current-holder}
 R^2\|DA\|_\infty+R^{2+\alpha}[DA]_\alpha
                       \le C_\alpha(\kappa+j),\qquad
 R^3\|\Delta_\eta A\|_\infty\le C(j+\kappa^2).
\end{equation}
This supplies \eqref{mb:eq-local-covariant-strong-coefficients}
when a Coulomb frame and the actual current are specified.
\end{corollary}

\begin{proof}
Scale to \(R=1\), with mesh \(h=\eta/R\), so the potential,
curvature and current are multiplied by \(R,R^2,R^3\),
respectively. The current scale also follows directly by
dilating its defining variation and physical pairing.
By the interior flat Hodge estimate, \(\|DA\|_p\le C_p\kappa\)
on an intermediate stock. For the pointwise equation let
\(S_h(F)=(P-P^{-1})/(2ih^2)\) for \(P=e^{ih^2F}\).
The exact Wilson current is the covariant curl transpose
of \(S_h(F)\). In coordinates this reads
\[
 J_{{\rm left},\mu}(x)
   =h^{-1}\sum_{\nu\ne\mu}
       \bigl(S_{h,\mu\nu}(x)-
         \operatorname{Ad}U_\nu(x-he_\nu)^{-1}
                    S_{h,\mu\nu}(x-he_\nu)\bigr).
\]
Indeed varying a positive link gives the plaquette insertion
with coefficient \(ih\). Its reversed occurrence gives the
negative transported insertion at the other plaquette base.
Cyclicity of trace brings both to the varied link tail.
In each trace, the traceless insertion pairs with
\((P-P^{-1})/2\), since \(P+P^{-1}\) is scalar.
Conjugation by \(P\) leaves \(P-P^{-1}\) fixed.
Dividing the resulting \(h^4\) contribution by the physical
bond weight gives the displayed formula.

Replacing the transport in that formula by the identity
costs \(C\|A\|_\infty\|F\|_\infty\).
The sine remainder satisfies
\(|S_h(F)-F|\le Ch^4|F|^3\), and a first difference of
this remainder costs at most \(Ch^3\|F\|_\infty^3\).
In \(F=dA+E(A)\), the analytic curvature remainder
satisfies \(|DE(A)|\le C\|A\|_\infty |DA|\) on a fixed
stencil enlargement. This follows by subtracting its
values at neighbouring plaquettes and differentiating the
analytic quadratic remainder with respect to its four link
arguments. Since \(d^*A=0\) and \(d^*dA=\Delta_h A\)
on full Coulomb stencils, we obtain
\[
 |\Delta_h A-J_{\rm left}(U)|
              \le C\bigl(\kappa\max_{\rm stencil}|DA|
                                   +\kappa^2+h^3\kappa^3\bigr).
\]
Thus \(\Delta_h A\) is bounded in any fixed \(L^p\) on the
intermediate stock. Apply \(D R_1\) to the equation for a
cutoff of \(A\), taking \(p>4\). Its source consists of the
cut-off \(\Delta_h A\) and bounded cutoff products of \(DA,A\).
Equation~\ref{mb:eq-heat-sobolev} gives
\(\|DA\|_\infty\le C(\kappa+j)\) on a smaller stock.
Insert this bound into the pointwise equation to obtain
\(\|\Delta_h A\|_\infty\le C(j+\kappa^2)\).
The bounded ordinary forcing estimate
\eqref{mb:eq-flat-poisson-gradient} gives the stated H\"older
bound on the innermost stock. For bounded \(R/\eta\), direct
finite-stencil bounds from the potential supremum give
the first-difference and H\"older bounds. The same pointwise
identity then gives the claimed Laplacian bound.
Undo the dilation to obtain \eqref{mb:eq-current-holder}.
\end{proof}

For real coefficients the corresponding transported Hölder
seminorm compares values in a common frame by parallel transport
on a fixed shortest coordinate path inside the cube.
In the chosen frame, such a transport over distance \(s\)
differs from the identity by at most \(C_Ks/R\), by
telescoping its factors. The same bound holds in the declared
complex class because a product along such a path, and its
inverse, have norm bounded by \(e^{C_Ks/R}\).
For first differences use the first-order Morrey estimate
on \(A\) with some \(q>4/(1-\alpha)\) and the expression
of \(D_U\) in terms of \(d\) and \(E_+(A)\).
This gives the transported first-difference bounds from
\eqref{mb:eq-local-covariant-point}. For second differences
use \eqref{mb:eq-local-covariant-schauder} and its supremum
bound. Real site gauges preserve these transported norms.

These local results use full stencils inside a lifted stock.
Norms on a periodic lift can be estimated by the corresponding
projected norms when the lift has bounded multiplicity, as
specified after Proposition~\ref{mb:prop-flat-poisson}.
A Neumann Laplacian on a graph retains only its
internal bond terms in its quadratic form, whereas a
Dirichlet compression retains the full quadratic form of
the zero extension. For a real unitary connection on either
graph, the heat semigroup is dominated in norm by the scalar
heat semigroup with that same boundary convention. To check
this, split the generator into its diagonal degree term and
its transported adjacency. The time-ordered expansion with
respect to the diagonal term sums over paths in the graph.
Each path has a nonnegative scalar weight times a product
of unitaries. Taking the norm gives the scalar expansion
with the same diagonal term, including the killing term
for zero extension. At an unbuffered boundary with bounded
\(R/\eta\), the point-value and present-stencil difference
estimates suffice. The interior bootstraps require their
stated collars. In particular a complex tube uses the
coefficient estimates rather than this unitary domination.


% Construction source: assembly/coulomb-frame.tex
\mbneedspace{10\baselineskip}
\subsection{Obtaining the local coefficient frame}

\begin{proposition}[A small real Coulomb frame]\label{mb:prop-coulomb-frame}
Let \(R/\eta\) be a positive integer. Suppose a real connection
is given on a complete lifted cube of radius \(8R\), with
\(\|F_U\|_\infty\le\kappa R^{-2}\) there in the logarithm
branch near one. For sufficiently small fixed \(\kappa\),
it admits a real gauge on the cube of radius \(6R\) in
which \(U=e^{i\eta A}\),
\[
          d^{\mathsf T}A=0,\qquad R\|A\|_\infty\le C\kappa .
\]
The divergence equation holds on complete stencils there.
The constant is independent of \(\eta/R\) and exterior data.
If the actual Wilson current has \(R^3\|j_U\|_\infty\le j\)
on the required complete inner stock, the frame also
satisfies \eqref{mb:eq-current-holder} on a fixed smaller
cube.
\end{proposition}

\begin{proof}
Rescale to \(R=1\) and denote the fine spacing by \(h\).
In a rooted comb gauge on the original cube, each link
holonomy is a product of at most \(C/h\) conjugated
plaquettes. Unitarity and the defect bound give
\(\|A\|_\infty\le C\kappa\) for the per-length link logarithm.
Expanding each plaquette product gives
\(F_U=dA+E(A)\), where
\(|E(A)(p)|\le C\sum_{b\subset\partial p}|A(b)|^2\);
the raw coefficients \(hA\) are in a fixed small
logarithm neighbourhood. Consequently \(\|dA\|_\infty\le C\kappa\).

Extend \(A\) to a periodic auxiliary lattice of side \(32\)
by multiplying by a fixed smooth cutoff, equal to one
on the cube of radius seven and zero before radius eight.
Use zero on any link not inner to the original cube.
This leaves every complete stencil at radius six unchanged,
including the neighbouring tails in its divergence.
For fine enough \(h\), each
plaquette meeting its nonzero support lies in the original
cube, and the difference product rule gives
\(\|d(\chi A)\|_\infty\le C(\|dA\|_\infty+\|A\|_\infty)\).
For the finitely coarse \(h\) the direct four-link bound
gives the same estimate because \(h\) has a fixed
positive lower bound. The extended connection therefore
has the same small bounds on its potential and physical
plaquette curvature up to fixed constants.

On the auxiliary torus use a mean-zero Hermitian site
field \(\lambda\) and gauge \(e^{i\lambda(x)}\).
For its left logarithm write
\[
 A^\lambda_\mu(x)=
    (ih)^{-1}\log(e^{i\lambda(x)}
       e^{ih A_\mu(x)}e^{-i\lambda(x+he_\mu)})
       = A_\mu(x)-d_\mu\lambda(x)+N_\mu(A,\lambda)(x).
\]
To estimate \(N\), let \(l=\lambda(x)\), \(r=hA_\mu(x)\)
and \(s=h d_\mu\lambda(x)\). The analytic logarithm of
\(e^{il}e^{ir}e^{-i(l+s)}\) vanishes at \(r=s=0\), and its
linear part at the origin is \(i(r-s)\).
After removing that linear part, the remainder is
\(O((|l|+|r|+|s|)(|r|+|s|))\); its \(l\)-derivative is
\(O(|r|+|s|)\) and its \(s\)-derivative is
\(O(|l|+|r|+|s|)\). These bounds follow by integrating
the bounded second derivatives on a fixed smaller
matrix ball. Since \(|s|\le2\|\lambda\|_\infty\),
division by \(h\) leaves a factor \(A\) or
\(d\lambda\) in the remainder estimate. With
\(\|\lambda\|_*=\|\lambda\|_\infty+\|d\lambda\|_p\),
for fixed \(p>8\), this proves
\[
 \begin{split}
 \|N(A,\lambda)\|_p &\le C(\kappa+\|\lambda\|_*)^2,\\
 \|N(A,\lambda)-N(A,\widetilde\lambda)\|_p
  &\le C(\kappa+\|\lambda\|_*+\|\widetilde\lambda\|_*)
                             \|\lambda-\widetilde\lambda\|_* .
 \end{split}
\]

The map \(\Delta_h^{-1}d^{\mathsf T}\) into this mean-zero
norm is bounded uniformly in \(h\). Its first difference
bound is Proposition~\ref{mb:prop-flat-lp}. The periodic
Poincar\'e inequality at side \(32\) bounds its \(L^p\)
norm by its first difference norm. This Poincar\'e
inequality follows by averaging coordinate paths between
pairs of sites. The first-order Sobolev estimate
\(\|\lambda\|_\infty\le C(\|\lambda\|_p+\|d\lambda\|_p)\)
follows from the periodic mean-zero estimate after
\eqref{mb:eq-gradient-morrey}, at the fixed side \(32\).
The zero mode is absent because \(\lambda\) has mean zero.
Thus
\[
 \lambda=\Delta_h^{-1}d^{\mathsf T}
                              (A+N(A,\lambda))
\]
is a contraction on \(\|\lambda\|_*\le C_0\kappa\):
choose \(C_0\) larger than twice the linear bound
for \(A\), then choose \(\kappa\) so the Lipschitz
constant is at most \(1/2\) and the quadratic term
fits in the remaining half. Its real fixed point
satisfies \(d^{\mathsf T}A^\lambda=0\) and
\(\|A^\lambda\|_p\le C\kappa\).
The raw logarithms \(hA^\lambda\) stay in the fixed
small branch by \(\|\lambda\|_\infty=O(\kappa)\).

Curvature remains \(O(\kappa)\) under the real gauge.
The same plaquette expansion now gives
\[
 \|dA^\lambda\|_{p/2}\le C(\kappa+\|A^\lambda\|_p^2).
\]
Use the torus Hodge bound \eqref{mb:eq-flat-hodge}
with zero divergence, and the first-order Sobolev
bound at \(p/2>4\), to obtain
\(\|A^\lambda\|_\infty\le C\kappa\).
Restrict the composition of the two gauges to the
original inner cube. Its links agree with the
auxiliary connection there. The actual current on
complete inner stencils transforms by the unitary
colour rotations. The last assertion follows from
Corollary~\ref{mb:cor-current-holder} and rescaling.
\end{proof}


% Construction source: assembly/flat-contour.tex
\section{The flat classical form and the contour gauge}\label{mb:sec-flat-contour}

We first use raw logarithms and counting measures. The fine
mesh in the flat comparison is one, the intermediate grid has
spacing \(n=L^k\), and its curl \(\mathrm d_n\) uses unscaled
differences. In the subsequent one-step contour problem we measure
length in units of that intermediate spacing. The flat action
Hessian is \(\|\mathrm da\|_2^2\) for
\(\langle X,Y\rangle=\tfrac12\operatorname{tr}(XY)\). All
flat scalar calculations apply to the three colours separately.

For a positively oriented \(\mu\)-link let
\[
 (S_n a)_\mu(y)=n^{-4}\sum_{x\in B_n(y)}
                          \sum_{t=0}^{n-1}a_\mu(x+t e_\mu).
\]
Write \(P_n\) for the raw derivative of the plain \(k\)-fold
average at the identity, \(r_n\) for restriction to its centres,
and \(q_n\) for the untransported block mean. The flat rows of
block averaging give
\begin{equation}\label{mb:eq-flat-plain}
 P_n a=S_n a+\mathrm d_n\sigma_n(a),\qquad
 P_n\mathrm df=\mathrm d_n r_nf,\qquad
 S_n\mathrm df=\mathrm d_n q_nf
\end{equation}
for a linear site map \(\sigma_n\). Indeed the first step differs
from \(S_L\) by the difference of the two averaged contour sums.
At each further step the centred block digits and the segment
digits partition their longer counterparts. Thus the straight
rows compose, and the gradient difference remains a gradient
under further straight rows. The two identities on gradients
follow by telescoping the segment and the contours, respectively.

\begin{proposition}[The flat effective action]\label{mb:prop-flat-effective}
On a compatible periodic lattice, the raw quadratic value
\[
 E_n(a)=\min_{P_nv=a}\|\mathrm dv\|_2^2
\]
is also the minimum subject to \(S_nv=a\). It satisfies
\begin{equation}\label{mb:eq-flat-effective}
            \|\mathrm d_n a\|_2^2\le E_n(a)
                                   \le 256\|\mathrm d_n a\|_2^2 .
\end{equation}
The corresponding translation-invariant matrix has a Fourier
symbol holomorphic and bounded on a fixed complex strip,
independent of \(k\) and the periods. Its kernel and the periodic
image sums have uniform exponential decay in intermediate units.
The same bounds hold for the real flat unitary connections in
their transported coordinates.
\end{proposition}

\begin{proof}
Both \(q_n\) and \(r_n\) are onto. If \(P_nv=a\), add a fine
gradient with \(q_nf=\sigma_n(v)\). If \(S_nv=a\), add a fine
gradient with \(r_nf=-\sigma_n(v)\). By \eqref{mb:eq-flat-plain}
these changes leave the curl unchanged and give the stated
constraints in the other convention.

The coarse curl of \(S_nv\) averages fine plaquette sums over
\(n\)-by-\(n\) sheets. Its nonnegative weights have row sum
\(n^2\) and counting column sum \(n^{-2}\). For the column sum,
first sum over all coarse origins and block sites, which gives
each fine translate weight \(n^{-4}\), and then sum the \(n^2\)
sheet positions. Weighted Cauchy--Schwarz gives the lower bound
in \eqref{mb:eq-flat-effective}.

For the upper bound use the nodal hat
\(\varphi(t)=(1-|t|/n)_+\) and interval
\(\beta(t)=1_{\{0,\ldots,n-1\}}(t)\). Define \(I_j\) on coarse
\(j\)-cochains by the product interpolation with factor \(n^{-j}\),
using \(\beta\) in the \(j\) oriented coordinate directions and
\(\varphi\) in the others. The products are centred at the
cochain's starting site, and are periodized on a torus.
The identity
\(\varphi(t+1)-\varphi(t)=n^{-1}[\beta(t+n)-\beta(t)]\)
gives \(\mathrm d I_1=I_2\mathrm d_n\). The row and column
sums of \(I_2\) are \(n^{-2}\) and \(n^2\). Hence
\(\|I_2\|_{2\to2}\le1\).

In each coordinate the convolution \(S_n I_1\) has weights
\[
 a_0=\frac34+\frac1{4n^2},\qquad
 a_{\pm1}=\frac{1-n^{-2}}8 .
\]
In a transverse coordinate the two off-centre coefficients
are the sums of the linear tails of the hat. In the segment
coordinate they count the overshoots of a translated segment,
one triangular sum \(1+\cdots+(n-1)/2\) on each side, divided
by \(n^2\). Thus \(S_n I_1\) is the same four-factor convolution
on every link component. Each factor has Fourier multiplier
at least \(a_0-2a_1\ge1/2\). Its inverse commutes with
\(\mathrm d_n\), so \(I_1(S_n I_1)^{-1}a\) is a feasible
straight lift with curl norm at most \(16\|\mathrm d_na\|_2\).
This proves the upper bound and nonemptiness. Finite-dimensional
orthogonal projection onto the image of the curl gives attainment.

Equivalently the minimizing curvature is the least-norm exact
two-cochain with the prescribed sheet-average curl. To verify the
converse, lift any such exact cochain as \(\mathrm dv\). The
difference \(S_nv-a\) is closed. On the periodic lattice its
nonzero Fourier modes are coarse gradients, and its constant mode
is a constant link field. These can be removed by a fine gradient
and a constant fine link field, respectively. Both operations
preserve the fine curl. In a twisted unitary character sector the
same argument applies, with a constant mode only in the invariant
sector. The sheet averages compose by the same block and segment
digit partition. Their kernels on exact fine cochains are therefore
nested, and the least-norm curvatures are the orthogonal projections
onto the decreasing orthogonal complements of those kernels.

To obtain the symbol, let \(p\) be the intermediate momentum
and \(\zeta_m=(p+2\pi m)/n\), with
\(m\in\{-(n-1)/2,\ldots,(n-1)/2\}^4\). Put
\[
 \begin{aligned}
 b_{\nu,m}(p)&=\frac{\sin((p_\nu+2\pi m_\nu)/2)}
                   {n\sin((p_\nu+2\pi m_\nu)/(2n))},\\
 f_{\mu,m}&=\Bigl(\prod_\nu b_{\nu,m}^2\Bigr)b_{\mu,m}^2,\qquad
 s_m=4n^2\sum_\nu\sin^2((p_\nu+2\pi m_\nu)/(2n)).
 \end{aligned}
\]
The removable sinc ratios use their analytic values.
With unitary counting Fourier transforms the straight row
acts on an alias by the multiplier
\(n^{-1}e^{i(n-1)\zeta_{m,\mu}/2}
(\prod_\nu b_{\nu,m})b_{\mu,m}\).
The factor \(n^{-1}\) is \(n\) from the segment times \(n^{-2}\)
from the ratio of the square roots of the two site counts.
Adding the squared fine divergence to the curl gives the
scalar Feynman symbol \(s_m/n^2\) on each alias. Its pushed
quadratic covariance is therefore diagonal with entries
\[
 D_\mu(p)=\sum_m f_{\mu,m}(p)/s_m(p)
 \quad(p\text{ real},\ p\ne0).
\]
The squared factors are the continuation of the product of the
two oppositely signed filters. The starting-link phase cancels
in that product.

Let \(t^\pm_\mu=e^{\pm ip_\mu}-1\), as columns. Removing a fine
gradient makes its divergence zero and changes the straight
datum only by a coarse gradient. Conversely a fine gradient can
restore any such change because \(q_n\) is onto. Minimization
of the Feynman value over \(a+t^+z\) gives the effective symbol
\begin{equation}\label{mb:eq-flat-symbol}
 \mathcal E_n
  = D^{-1}
    -\frac{D^{-1}t^+(t^-)^{\mathsf T}D^{-1}}
                 {(t^-)^{\mathsf T}D^{-1}t^+},
 \qquad D=\operatorname{diag}(D_\mu).
\end{equation}
For a real nonzero \(p\) the central alias is positive in each
component, and the denominator is strictly positive.
The constant mode has zero action by its constant fine lift.

On \(|\operatorname{Re}p_\nu|<\pi+\delta\),
\(|\operatorname{Im}p_\nu|<\delta\) with \(\delta>0\) fixed
sufficiently small, the sine denominator gives
\(\left|b_{\nu,m}\right|\le C(1+|m_\nu|)^{-1}\).
For \(m\ne0\), \(\operatorname{Re}s_m\ge c(1+|m|^2)\):
the real sine arguments lie in a fixed subinterval of
\((- \pi,\pi)\), and the imaginary squares cost \(O(\delta^2)\)
after multiplying by \(n^2\). Hence the sum \(R_\mu\) over
noncentral aliases is holomorphic and uniformly bounded by
\(C\sum_{m\in\mathbb Z^4}\prod_\nu(1+|m_\nu|)^{-2}\).
Write \(s=s_0\), \(f_\mu=f_{\mu,0}\),
\(B=\prod_\nu b_{\nu,0}^2\), and
\(H_\mu=(f_\mu+s R_\mu)^{-1}\).
Near zero \(s=O(|p|^2)\) and \(f_\mu=1+O(|p|^2)\), with
constants independent of \(n\). The exact identity
\(\sum_\mu t_\mu^-t_\mu^+/f_\mu=s/B\) gives
\[
 \sum_\mu t_\mu^-t_\mu^+H_\mu=s J,\qquad
 J=B^{-1}-\sum_\mu t_\mu^-t_\mu^+R_\mu H_\mu/f_\mu .
\]
On a fixed neighbourhood of zero \(J=1+O(|p|^2)\) is
nonzero. Substituting \(D^{-1}=sH\), with
\(H=\operatorname{diag}(H_\mu)\), cancels both powers of \(s\)
in \eqref{mb:eq-flat-symbol}. The extension is
\[
       \mathcal E_n=sH-J^{-1}Ht^+(t^-)^{\mathsf T}H .
\]
This removes the complex zero cone of \(s\), and gives the
stated value at zero. Away from that neighbourhood, the
central real \(b_{\nu,0}\) are bounded above and away from
zero, and \(s\) and \(\sum_\mu|t_\mu^+|^2\) are bounded
below. Thus \(D\) and the denominator in
\eqref{mb:eq-flat-symbol} are bounded away from zero on the
remaining real compact set. The uniform alias estimates
and Cauchy's bound on a smaller complex neighbourhood keep
these lower bounds on a fixed strip there. The expressions
agree on overlap by the displayed cancellation. Reindexing
aliases proves periodicity. Shifting the Fourier contour in
each direction gives an exponential kernel bound; sampling
the symbol on a finite torus gives its absolutely convergent
image sum.

For a real flat connection, take a parallel frame on the
cover. The commuting unitary adjoint holonomies decompose
into characters of the deck group. The symbol is unchanged
in each character sector and is sampled at shifted real
momenta. The images acquire unit-modulus factors, which
preserve the same bound. The interpolation and gradient
corrections also commute with these deck actions.
\end{proof}

For a real reference \(V\) on a lattice of spacing one, write
\(P_V=D q_V(0)\), where
\[
 q_V(a)=i^{-1}\log(\mathcal M(e^{ia}V)\mathcal M(V)^{-1}).
\]
On a retained family of complete \(L\)-blocks include all
faces incident to those blocks, with their double blocks and
all internal contour paths. Extend a field by zero on the
remaining links if only retained coordinates are varied.
Let \(b_V a\) be the derivative of the relative contour
logarithm
\[
 i^{-1}\log[T_{e^{ia}V}(y,x)T_V(y,x)^{-1}]
\]
at zero, with one row
for each \(x\ne y\) in a retained block. At the identity it
is the arithmetic contour path sum. The quadratic form
\(\|b_V a\|_2^2\) is the relative trace-contour Hessian.
When the contours of \(V\) are one it is the absolute
trace-contour Hessian.

\begin{lemma}[The outer curl and contour inequality]\label{mb:lem-outer-geometry}
For \(L\) fixed and sufficiently small plaquette defect on the
complete retained stocks,
\begin{equation}\label{mb:eq-outer-geometry}
 \|a\|_2^2\le C_L\bigl(\|\mathrm d_Va\|_{2,\mathcal P}^2
                           +\|b_Va\|_2^2+\|P_Va\|_2^2\bigr).
\end{equation}
Here the supported links are those incident to retained blocks,
the plaquette stock \(\mathcal P\) contains the internal plaquettes
of the retained blocks and the face plaquettes comparing
neighbouring crossing links, and \(P_V\) includes each incident
face row. If the field vanishes on a spanning tree inside each
retained block, the contour term can be omitted.
The constants are independent of the retained family and volume.
\end{lemma}

\begin{proof}
At a flat reference use a parallel frame on each rectangular
one- or two-block stock. Inside a retained block, a path from
its centre to the head of an internal link and a path to its
tail followed by that link differ by a sum of elementary
plaquette boundaries. To see this, replace successive path
steps by coordinate paths to the centre, interchanging adjacent
different coordinate steps by their plaquette. The paths have
bounded length depending on \(L\); summing these elementary
differences and averaging over the prescribed contours gives a
bound on the internal link by two contour sums and a finite sum
of curls. A tree path can be used instead when its link values
vanish. Internal links of an unretained neighbour have zero
values.

On a face, a plaquette with one transverse step expresses the
difference of adjacent crossing values using its curl and its
two internal links. The transverse grid is connected and has
\(L^3\) vertices. Thus all crossing values differ from a common
one by terms already bounded. Every straight segment in the
face row has one crossing and otherwise internal links.
Its plain row is their arithmetic average plus the endpoint
contour correction. Its total coefficient on the common
crossing value is one. The row therefore bounds that common
value. Squaring these bounds uses only constants depending
on \(L\), and each link or plaquette is used a bounded number
of times.

For a real curvature-small field, put each such stock in its
comb frame. The Wilson curl, the contour derivative and the
plain row differ in that frame from the flat matrices by
\(O_L(\max|\partial V-1|)\) in the local coefficient norms.
This follows by differentiating the admitted finite contour
and logarithm formulas on their fixed one-step balls.
The change of frame is a unitary rotation at each input or
output tail, and it preserves zero tree values. The squared
errors summed over the complete stocks are at most
\(C_L\max|\partial V-1|^2\|a\|_2^2\) by their overlap.
Decrease the defect ceiling to absorb them. If a stock wraps,
make the estimates on its rectangular cover with the
given projected branch tests and descend by deck invariance.
\end{proof}

\begin{corollary}[Actual contour precisions]\label{mb:cor-contour-gap}
Let \(C_V\) be the derivative of the actual pivot graph and
\(h_V\) its pivot-supported right inverse of \(P_V\).
Use the ambient Wilson action whose plaquette set contains
every internal and face-comparison plaquette in the stock $\mathcal P$ of Lemma~\ref{mb:lem-outer-geometry}, with the varied link field
extended by zero. Any selection of free coordinates is
compressed after forming the chart precision.
At a real reference with sufficiently small defect, the one-step
matrix
\begin{equation}\label{mb:eq-one-step-contour}
 C_V^{\mathsf T}H_W(V)C_V+(b_VC_V)^{\mathsf T}(b_VC_V)
  -\left\langle h_V^{\mathsf T}J,\,
                         D^2q_V(0)[C_V\,\cdot,C_V\,\cdot]\right\rangle
\end{equation}
has spectrum in \([c_L,C_L]\) for real independent raw currents
\(J\) with sufficiently small supremum norm. For the actual
Wilson current at an axial reference it is the Hessian of the
Wilson and absolute contour actions in the exact fibre chart.
At a constrained-critical reference this is the physical
Gaussian Hessian. Its inverse is
holomorphic on a fixed smaller complex coefficient
neighbourhood and has an exponential kernel there.

At the flat field the corresponding next-step precision
after \(k\) classical averaging levels is
\(C_1^{\mathsf T}\mathcal E_n C_1+
 (b_1C_1)^{\mathsf T}(b_1C_1)\), with \(n=L^k\) and the
intermediate lattice spacing taken as one. Its spectral bounds
and the exponential kernel bounds for this precision and its
inverse are uniform in \(k\) and compatible volume.
The flat upper spectral bound may be chosen as $CL^8$,
with a numerical $C$ independent of $L$ and $k$.
\end{corollary}

\begin{proof}
The pivot graph has free restriction equal to the identity and
\(P_V C_V=0\). Its second derivative lies on the pivots and
equals \(-h_VD^2q_V(0)[C_V\,\cdot,C_V\,\cdot]\) by two
differentiations of the constraint. The chain rule gives the
current term in \eqref{mb:eq-one-step-contour}, including its
sign. The absolute contour term has the relative Hessian at
the axial reference, and its current is zero there.

The finite one-step pivot, mean and logarithm bounds give
\(\|C_Vu\|_2\le C_L\|u\|_2\) and operator norm at most
\(C_L\|J\|_\infty\) for the current term. The Wilson second
variation is the transported curl square up to
\(C_L\max|\partial V-1|\|C_Vu\|_2^2\), by expanding the
four insertion factors in a plaquette trace. Apply
\eqref{mb:eq-outer-geometry} and use
\(\|C_Vu\|_2\ge\|u\|_2\) to absorb the two errors. The
same finite coefficient bounds give the upper estimate.
The entries are holomorphic on admitted complex tubes.
A smaller coefficient increment changes the matrix in norm
by less than half the gap, since each row and column reads a
fixed number of raw coordinates. Inverting by a Neumann
series proves uniform analyticity. Conjugation by
\(e^{s\,d(\cdot,b)}\) changes the matrix in norm by \(O_L(|s|)\)
because it has fixed range, so the conjugated inverse stays
bounded for a fixed small \(|s|\). Taking a single column
gives the exponential bound.

For the flat classical form use \eqref{mb:eq-flat-effective}
and \eqref{mb:eq-outer-geometry}; its curl-null fields
include the constant modes, which are controlled by the
actual outer rows and contour form. The flat \(\mathcal E_n\)
has uniformly summable exponential kernels by
Proposition~\ref{mb:prop-flat-effective}. Composing with the
fixed local \(C_1\) and \(b_1\) preserves such kernels on the
block lattice. The same conjugation estimate applies using
the exponential first moment in place of finite range.
For the explicit flat upper bound, the straight raw row
$S_L$ has absolute row sum $L$ and column sum $L^{-3}$.
The internal contour mean has row sum $O(L)$ and column
sum at most one. Thus $P_L=S_L-\mathrm d_L\sigma$
has norm at most $CL^{1/2}$. The pivot right inverse has
norm $L^3$, giving $\|C_1\|\le CL^4$.
The contour derivative has row sum $O(L)$ and column sum
at most $L^4$, and it does not read the crossing pivots.
Hence $\|b_1C_1\|^2\le CL^5$. Combine these counts with
$\mathcal E_n\le256\mathrm d_n^{\mathsf T}\mathrm d_n$
and the bounded unit-grid curl to obtain $CL^8$.
Periodic image bounds and the flat character decomposition
preserve both estimates on tori.
\end{proof}


% Construction source: assembly/rows-restoration.tex
\section{Site averages, restoration and determining row lifts}
\label{mb:sec-rows-and-restoration}

Write \(\Xi_s\) for multiplication by \(\xi_y\) on site tops.
For a fine bond or plaquette the scale \(\mathsf L\) is the
largest assigned corner scale. We use the field energy space \(X\)
with norm
\begin{equation}\label{mb:eq-field-energy}
 \|v\|_X^2=\|\nabla_Uv^0\|_2^2+\|v\|_{\rm m}^2,
 \qquad \|v\|_{\rm m}=\|\mathsf L^{-1}v\|_2 .
\end{equation}
Here \(v^0\) is zero off the active bonds. The comparison norm
counts its boundary jumps. If a stencil has two nonzero bond
values, the connecting shift is inside the active graph. With
only one nonzero value its norm is independent of the arbitrary
unitary exterior connection. The form of the Wilson action
still uses only inner plaquettes. Let \(Y\) be the field-row
space with norm
\begin{equation}\label{mb:eq-row-energy}
 \|b\|_Y^2=\|a_f^{1/2}b\|_{2,r}^2,\qquad
               a_f(c)=d_c\xi_c^{-2},\quad 0<d_*\le d_c\le d^*.
\end{equation}
Thus \(Y\) is uniformly equivalent to the norm
\(\|b\|_{Y_0}^2=\sum_c \xi_c^2|b_c|^2\). The dual norm for
the physical row pairing is
\(\|\omega\|_{Y^*}=\|a_f^{-1/2}\omega\|_{2,r}\).

\subsection{Weighted site estimates}

Choose \(\delta>0\) small and fixed, then impose \(\rho\delta\ge1\).
On the whole torus set
\begin{equation}\label{mb:eq-slow-length}
 \ell(x)=\inf_z\{\xi(z)+\delta\,\operatorname{dist}_\infty(x,z)\}.
\end{equation}
It is \(\delta\)-Lipschitz and
\(\xi(x)/L\le \ell(x)\le\xi(x)\). In fact if \(x\) has level \(j\),
a \(z\) at level at least \(j-1\) has \(\xi(z)\ge\xi_j/L\);
a \(z\) at a smaller level is at distance \(>\rho\xi_{j-1}\)
by separation. On the bare layer \(\ell=\eta\). The values at
neighbouring sites and on their fine bond are comparable, and
\(|D(\ell^{-1})|\le C\delta\ell^{-2}\) for a fine difference.
In norms containing a bond scale, \(\ell\) is evaluated at its tail.
The metric defined using \(\ell\) instead of \(\xi\) lies between
\(d\) and \(Ld\).

Set \(\Delta_U=D_U^{\mathsf T}D_U\) on the true graph and
\(a_s(y)=c_y\xi_y^{-2}\), where \(0<c_*\le c_y\le c^*\).
Let \(A_s=\Delta_U+q^{\mathsf T}a_s q\).
For a real connection in \eqref{mb:curvature-class},
\begin{equation}\label{mb:eq-site-poincare}
 \begin{split}
 \|\ell^{-1}u\|_2^2
       &\le C(\|D_Uu\|_2^2+\|\Xi_s^{-1}qu\|_{2,s}^2),\\
 \|u\|_2^2&\le C(\|\ell D_Uu\|_2^2+\|qu\|_{2,s}^2).
 \end{split}
\end{equation}
For each nonbare pillar, use a comb frame on its complete
surrounding box. Its potential is \(O_L(\kappa/\xi_y)\).
The \(q\)-transports are \(I+O_L(\kappa)\) in this frame
by the summable block transport comparison. The ordinary
block Poincar\'e inequality follows by writing a difference
as a coordinate path and summing over both endpoints.
The covariant gradient and the actual mean change its right
side by at most \(C_L\kappa^2\xi_y^{-2}\|u\|_{2,B_y}^2\)
after enlarging its constant. Absorb this term by taking
\(\kappa\) small. The bare pillar has the identity mean.
Sum the block inequalities to prove the first line of
\eqref{mb:eq-site-poincare}, or multiply each by \(\xi_y^2\)
before summing to prove the second. Only internal pillar bonds
are used. In particular \(A_s\) is positive definite.

\begin{lemma}[Compact lifts of the actual site average]\label{mb:lem-site-lift}
Under the preceding real curvature condition there is a
right inverse \(E_s\) of \(q\) with seeds supported in their
pillars. For \(r=-1,0,1\),
\begin{equation}\label{mb:eq-site-lift-H1}
 \|\ell^{-r}E_s b\|_2\le C\|\Xi_s^{-r}b\|_{2,s},\qquad
 \|\ell^{-r}D_U E_s b\|_2\le C\|\Xi_s^{-r-1}b\|_{2,s}.
\end{equation}
It can be chosen with the additional bound
\begin{equation}\label{mb:eq-site-lift-H2}
 \|E_s b\|_{V_2}\le C\|\Xi_s^{-2}b\|_{2,s},\qquad
 \|u\|_{V_2}^2=\|\Delta_Uu\|_2^2+\|\ell^{-1}D_Uu\|_2^2+
                                            \|\ell^{-2}u\|_2^2 .
\end{equation}
The constants are independent of mesh, depth and volume.
\end{lemma}

\begin{proof}
On a complete nonbare box of side \(C_L\xi_y\) containing
the pillar, choose its comb frame. For a fixed colour \(b_y\)
let \(u\) minimize the covariant gradient energy with boundary
value \(b_y\) in that frame. The Dirichlet form is positive
on zero boundary fields by paths to the boundary, and its
equation is \(\Delta_Uu=0\) at interior sites. Comparison
with the constant field and the small potential give
\(\|D_Uu\|_2\le C_L\kappa\xi_y|b_y|\). For \(u-b_y\) the
zero-boundary Poincar\'e inequality has the bound
\(\|u-b_y\|_2\le C\xi_y\|D_U(u-b_y)\|_2\). The ordinary
inequality follows by summing along paths parallel to one
axis to the boundary. In the comb frame its replacement by
\(D_U\) costs \(C_L\kappa\|u-b_y\|_2\), which is absorbed.
Thus \(\|u-b_y\|_2\le C_L\kappa\xi_y^2|b_y|\).

Multiply \(u\) by a fixed smooth bump sampled at scale \(\xi_y\),
supported strictly inside \(B_y\) and with positive mean.
For sufficiently fine ratios \(\eta/\xi_y\), its two-step collar
lies in the harmonic region. Its first two differences are
bounded by \(C\xi_y^{-1}\) and \(C\xi_y^{-2}\). The product rule
therefore gives \(\|\chi u\|_2\le C\xi_y^2|b_y|\) and
\(\|D_U(\chi u)\|_2\le C\xi_y|b_y|\). Moreover
\[
 \|\Delta_U(\chi u)\|_2
       =\|[\Delta_U,\chi]u\|_2
       \le C(\xi_y^{-1}\|D_Uu\|_2+\xi_y^{-2}\|u\|_2)
       \le C|b_y|.
\]
For the commutator, expand the two opposite covariant shifts
of \(\Delta_U\). Terms with a shift difference of \(u\) cost
a first difference of \(\chi\), and the terms proportional to
\(u\) have its second difference. This uses no derivative of
the connection. The transported average of \(\chi u\) is
\((m_y I+O_L(\kappa))b_y\), where \(m_y\) is the ordinary
mean of \(\chi\), uniformly bounded away from zero. The error
uses the \(L^2\) bound on \(u-b_y\) and the actual transports
on the pillar. Normalize the seed by the inverse of this
matrix and undo the frame.

For the finitely coarse ratios, use a single site seed,
normalized by the inverse of its exact site-row coefficient.
Its amplitude is bounded by \(C|b_y|\) because the number of
sites is then bounded. Its \(i\) differences cost at most
\(C\eta^{-i}\) in a fixed collar, which is \(O(\xi_y^{-i})\)
in these cases. This includes the bare pillar. Distinct
pillars are disjoint and their fixed-stencil collars have
bounded overlap and comparable scales. Summing proves
\eqref{mb:eq-site-lift-H1} and \eqref{mb:eq-site-lift-H2}.
If a box does not embed, construct on a large periodic cover
and average the linear lift over simultaneous deck shifts
of source and output. The right inverse equation and the
norm estimates are preserved and descend with the common
sheet-count factor.
\end{proof}

\begin{proposition}[Normalized site Gram matrix and restoration]
\label{mb:prop-site-gram}
Under the preceding real hypotheses put
\begin{equation}\label{mb:eq-site-projector}
 G_s=A_s^{-1},\quad F=G_s q^{\mathsf T}\Xi_s^{-2},\quad
 M=F^{\mathsf T}F,\quad P=FM^{-1}F^{\mathsf T},\quad R=I-P .
\end{equation}
There are uniform constants \(m_*,m^*>0\) with
\(m_*I\le M\le m^*I\). The normalized columns \(F,\ell D_UF\)
and \(M^{-1}\) are bounded after conjugation by
\(e^{s\phi}\) for \(|s|\le s_0\), where \(\phi\) is
one-Lipschitz for the metric of \(\ell\) and its top samples
are used on top spaces. In particular their normalized
cell-\(L^2\) entries decay exponentially.
The fine blocks of \(G_s\) have the bound
\(C\xi_y\xi_z e^{-s_0 d(y,z)}\) from \(B_z\) to \(B_y\);
those of \(D_UG_s\) have bound \(C\xi_z e^{-s_0 d(y,z)}\).

Define
\begin{equation}\label{mb:eq-site-restoration-definition}
                  \Lambda=G_sRD_U^{\mathsf T},\qquad S=I-D_U\Lambda .
\end{equation}
Then
\begin{equation}\label{mb:eq-restoration-properties}
 \begin{split}
 \operatorname{ran}R&=\Delta_U(\ker q),\qquad
 qG_sR=0,\qquad\Delta_U G_s R=R,\\
 RD_U^{\mathsf T}S&=0,\qquad QS=Q,\qquad S^2=S,\\
 \|\Lambda v\|_{V_2}+\|D_U\Lambda v\|_X+\|Sv\|_X
                         &\le C\|v\|_X .
 \end{split}
\end{equation}
On a fixed smaller weighted \(C^1\) complex tube these formulas
continue holomorphically, with the same bounds up to fixed
constants. \(R\) there is the bilinear projection given by
\eqref{mb:eq-site-projector}.
\end{proposition}

\begin{proof}
Let \(u=Fb\) and
\(\mathcal E_+(u)=\|\ell D_Uu\|_2^2+\|qu\|_{2,s}^2\).
Test \(A_su=q^{\mathsf T}\Xi_s^{-2}b\) against
\(\ell^2u\). On each real bond the exact real-part identity
\[
 \operatorname{Re}\langle D_Uu,D_U(\ell^2u)\rangle
    = |D_U(\ell u)|^2
       -|D\ell|^2\operatorname{Re}
           \langle u_-,\operatorname{Ad}U_b u_+\rangle
\]
bounds the summed gradient term below by
\(\tfrac12\|\ell D_Uu\|_2^2-C\delta^2\|u\|_2^2\).
The weighted difference of the actual mean and its constant
weight on a pillar satisfies
\[
 \|\Xi_s^{-2}(q(\ell^2u)-\ell(y)^2 qu)\|_{2,s}
                     \le C_L\delta\|u\|_2 .
\]
Indeed the absolute variation of \(\ell^2/\xi_y^2\) on
a pillar is \(O_L(\delta)\); the transported average is a
contraction on the physical \(L^2\) spaces. Since
\(\ell(y)\ge\xi_y/L\), the site mass term is bounded below
by \(c\|qu\|_{2,s}^2-C_L\delta^2\|u\|_2^2\) by Young's
inequality. Also \(\|\Xi_s^{-2}q(\ell^2u)\|_{2,s}\le\|u\|_2\).
Choose \(\delta\) small enough to absorb the errors by
\eqref{mb:eq-site-poincare}. The result is
\begin{equation}\label{mb:eq-gram-test}
       c\mathcal E_+(Fb)\le \|b\|_{2,s}\|Fb\|_2 .
\end{equation}
For the lower bound test the same equation against
\(E_s\Xi_s^2b\). Its exact \(q\)-value is \(\Xi_s^2b\).
The \(r=1\) lift bound controls the gradient term by
\(C\|\ell D_UFb\|_2\|b\|_{2,s}\); the mass term is the
pairing of \(qFb\) with \(c_y b\). Thus
\(\|b\|_{2,s}^2\le C\mathcal E_+(Fb)^{1/2}\|b\|_{2,s}\).
Together with \eqref{mb:eq-gram-test} this proves the lower
bound, and \eqref{mb:eq-site-poincare} proves the upper bound.

Conjugate the weighted site equation by \(e^{s\phi}\).
On a fine bond the change of a gradient is at most
\(C|s|\ell^{-1}\), since a single bond has metric length
\(O(\eta/\ell)\). A site mean and its transpose change by
\(C|s|\) in normalized form norms because its carrier has
bounded metric diameter. The preceding weighted energy test
is therefore changed by \(O(|s|)\) times its energy. It
bounds the conjugated \(F\) and \(\ell D_UF\) and gives an
\(O(|s|)\) difference from the unweighted operators by the
resolvent identity. This applies to the opposite weight
and hence to the conjugated transpose. Consequently the
conjugated \(M\) differs in norm by \(O(|s|)\) and its
inverse is bounded for a fixed sufficiently small \(s_0\).
For \(G_s\) use the first line of \eqref{mb:eq-site-poincare}
to test the conjugated equation in the ordinary energy norm
\(\|D_Uu\|_2+\|\ell^{-1}u\|_2\). Its right hand side \(f\)
has dual bound \(C\|\ell f\|_2\). Taking \(\phi\) to be
distance from the source cell and restricting the outputs
gives the two stated blocks, including their length factors.

The range identities follow directly from
\(F^{\mathsf T}A_s=\Xi_s^{-2}q\). If \(q\lambda=0\) then
\(A_s\lambda=\Delta_U\lambda\) belongs to \(\ker F^{\mathsf T}\).
Conversely \(F^{\mathsf T}f=0\) gives \(qG_sf=0\) and
\(\Delta_U G_sf=f\). Thus \(\ker F^{\mathsf T}=\operatorname{ran}R\).

We need two energy derivatives on \(\ker q\). Set
\(w=\ell^{-1}\), \(f=\Delta_U\lambda\) and \(q\lambda=0\).
The first block inequality applied to \(w\lambda\) has mean
error at most \(C_L\delta\|\ell^{-2}\lambda\|_2\). This
follows by subtracting \(w(y)q\lambda\) in each pillar.
It follows that
\(\|\ell^{-2}\lambda\|_2^2\le C\|D_U(w\lambda)\|_2^2\)
after absorbing the error. Testing the equation against
\(w^2\lambda\) and using the same bond identity gives
\[
 \|D_U(w\lambda)\|_2^2
  \le\|f\|_2\|\ell^{-2}\lambda\|_2+
                   C\delta^2\|\ell^{-2}\lambda\|_2^2 .
\]
A second absorption and the product rule prove
\(\|\ell^{-2}\lambda\|_2+\|\ell^{-1}D_U\lambda\|_2
 \le C\|f\|_2\).

For the remaining derivative, extend \(\lambda\) by zero
outside \(\Omega_0\). It vanishes on the bare layer because
its actual mean there is its site value. The separation
therefore keeps every nonzero two-step stencil of this
extension inside the active region. For the unitary
covariant shifts, expanding the squared curl and divergence
on the full torus yields the component-gradient square
up to terms containing a commutator of two shifts.
The diagonal terms agree by unitarity. In the mixed terms
interchanging the two shifted directions leaves a factor
\(\operatorname{Ad} U(\partial p)-I\) divided by \(\eta^2\).
The scale bound in \eqref{mb:curvature-class} bounds the
sum of these terms by \(C_L\kappa\|v\|_{\rm m}^2\).
The curl of \(D_U\lambda\) is itself such a commutator
applied to \(\lambda\), with norm at most
\(C_L\kappa\|\ell^{-2}\lambda\|_2\).
Applying this calculation to \(v=D_U\lambda\) gives
\begin{equation}\label{mb:eq-zero-site-H2}
 \|\lambda\|_{V_2}+\|\nabla_U D_U\lambda\|_2
                         \le C\|\Delta_U\lambda\|_2 .
\end{equation}
This calculation uses no exterior curvature, since all
terms with nonzero fields or shifted fields lie in the
bare separated collar or its interior.

Now \(\lambda=\Lambda v\) is in \(\ker q\) and
\(\Delta_U\lambda=RD_U^{\mathsf T}v\). The actual divergence
is bounded by \(C\|v\|_X\): it is the restriction of the
divergence of the zero extension, with exterior bond
values zero. Thus \eqref{mb:eq-zero-site-H2} proves the
norm bounds. The range identities imply
\(RD_U^{\mathsf T}S=0\), and \(S\) is the identity on that
kernel. The exact relation \eqref{mb:extended-gradient-identity}
gives \(QD_U\Lambda=0\), hence \(QS=Q\).

For a complex shift of size \(\varepsilon\) the gradient
coefficient changes by \(C\varepsilon\ell^{-1}\) and
the actual normalized site row by \(C\varepsilon\).
The latter follows by the holomorphic transported-mean
comparison and the geometric sum of raw radii over its
complete backward blocks. The \(A_s\) energy form,
the plus-weight tested form, and their conjugated forms
therefore have \(O(\varepsilon)\) perturbations in their
declared energy norms. A relative Neumann series gives
\(G_s,F\), and a second one gives \(M^{-1}\) and \(R\).
In particular \(\|R_{\mathbf U}-R_U\|_{2\to2}\le C\varepsilon\).
The range identities are algebraic. The covariant
Laplacian difference obeys
\[
 \|(\Delta_{\mathbf U}-\Delta_U)\lambda\|_2
 \le C\varepsilon\bigl(\|\ell^{-1}D_U\lambda\|_2+
                             \|\ell^{-2}\lambda\|_2\bigr).
\]
To obtain it, collect opposite shifts before estimating:
one connection increment multiplies a first difference
of \(\lambda\), and the derivative of the increment
multiplies \(\lambda\). The weighted \(C^1\) bound controls
these by \(\varepsilon/\ell\) and \(\varepsilon/\ell^2\),
including the product of two increments. On the bare
collar the direct finite-stencil bound has those same
powers because \(\ell=\eta\). For \(q_{\mathbf U}\lambda=0\),
subtract \(E_s q_U\lambda\) to use
\eqref{mb:eq-zero-site-H2}. Its \(V_2\) norm is at most
\(C\|\Xi_s^{-2}(q_U-q_{\mathbf U})\lambda\|_{2,s}
 \le C\varepsilon\|\ell^{-2}\lambda\|_2\)
by \eqref{mb:eq-site-lift-H2}. The subtracted lift's full
gradient-derivative norm is bounded by its \(V_2\) norm
up to \(C\): the preceding curl-divergence calculation
applies to interior seeds and the finitely coarse collars
are controlled by their atomic norm. Absorbing these errors
proves \eqref{mb:eq-zero-site-H2} on the moving kernel,
and then the claimed restoration bounds.
\end{proof}

\subsection{Energetic lifts of the field rows}

\begin{lemma}[A right inverse of the determining field rows]\label{mb:lem-field-lift}
For the separated no-inward determining geometry and the
curvature-small real class, the actual per-length \(Q\) has
a right inverse \(E_Q:Y\to X\) with norm at most \(C\).
There is a right inverse \(E_L=SE_Q\) into
\(X_L=\ker(RD_U^{\mathsf T})\), with the same uniform bound.
The lifts and the row inf-sup estimate extend to a fixed
smaller complex tube about a real reference. On this tube
the \(X\) norm may be taken at the real centre.
\end{lemma}

\begin{proof}
At the flat field consider a level \(j\ge1\) pillar with
\(n=L^j=2m+1\) sites per side. A seed for its \(\mu\)-direction
has only that component nonzero. Inside the pillar, at fine
index \(u\), take the product
\(\phi(u_\mu)\prod_{\nu\ne\mu}\chi(u_\nu)\) times its colour,
using only bonds with both endpoints in the pillar.
Here \(\chi\) is a nonnegative even bump sampled from a fixed
smooth profile of \(u/n\) supported in \((-1/4,1/4)\).
For sufficiently large \(n\), set
\(\beta=n^{-1}\sum_u\chi(u)\) and
\(s=n^{-1}\sum_u(u/n)^2\chi(u)\). Both are uniformly positive.
The straight per-length row weights are \(n^{-5}\)
times their segment multiplicities. The two rows incident
to this pillar have longitudinal counts \(u+m+1\) and
\(m-u\) at the positive link with tail \(u\). Their
response functionals are
\[
 \ell_+(\phi)=\beta^3n^{-2}\sum_u(u+m+1)\phi(u),\qquad
 \ell_-(\phi)=\beta^3n^{-2}\sum_u(m-u)\phi(u).
\]
The profiles
\[
 \phi_+(u)=\beta^{-4}\chi(u)+
       \frac m n\frac{\beta^{-3}}s\frac u n\chi(u),\qquad
 \phi_-(u)=\beta^{-4}\chi(u)-
       \frac{m+1}n\frac{\beta^{-3}}s\frac u n\chi(u)
\]
have exact responses \((1,0)\) and \((0,1)\), by evenness
and the two moments. Their first differences are \(O(n^{-1})\).
For the remaining odd \(3\le n\) below the fixed cutoff,
use transverse centre atoms and longitudinal atoms at
\(u=-1,0\). The two response columns form
\(n^{-5}\bigl(\begin{smallmatrix}m&m+1\\m+1&m\end{smallmatrix}\bigr)\),
of determinant \(-n^{-9}\). Its inverse gives the required
seeds with constants depending only on the cutoff.

Assign each row to a level-\(j\) endpoint, using the
interior endpoint for an outward row, and choose the seed
with the response on that side. The full regular grid at
that level sees exactly the one assigned row. At the bare
level the seed is the individual stored fine link.
At most eight seeds are assigned to one pillar. The support
volume \(\xi_j^4\) and one inverse length from each energy
factor give
\[
                   \|E^0b\|_X^2\le C\sum_c\xi_c^2|b_c|^2
\]
for the sum \(E^0\) of the seeds. This includes the bare
boundary jump, whose scale is \(\eta\).

A level-\(j\) seed is invisible to lower rows because their
carriers miss \(\Omega_j\). A row of level \(j+2\) or higher
has carrier in \(\Omega_{j+1}\), so it also misses the seed.
The only possible leakage is into the next level.
The flat straight rows on full grids satisfy
\begin{equation}\label{mb:eq-flat-row-composition}
                    Q_{j+1}^{0,\rm full}=Q_L^0 Q_j^{0,\rm full}.
\end{equation}
The site and segment digits partition their larger ranges,
so the weights multiply to \((Ln)^{-5}\).
The one-step \(Q_L^0\) has row sum one and counting column
sum \(L^{-4}\), the two longitudinal counts adding to \(L\).
Its counting \(L^2\) operator norm is at most \(L^{-2}\).
The \(Y_0\) norms add a factor \(L\) across levels. Therefore
\(Q^0 E^0=I+T\) with only adjacent-level upper leakage and
\(\|T\|_{Y_0\to Y_0}\le L^{-1}\). No depth factor enters:
the level blocks of \(T\) have distinct source and target
levels. Unit transports in local flat frames preserve these
bounds, also with nontrivial flat holonomies.

At a real curvature-small reference, put the same seed
profiles in their local comb frames, rotating the colour
to its own row tail frame. Use a common real frame on a
complete row carrier when computing its action. The gauge
transition between two such frames differs from its value
at the pillar centre by \(C_L\kappa\), by telescoping the
bond equation along a path of length \(O_L(\xi_j)\);
both potentials are \(O_L(\kappa/\xi_j)\).
The per-length straight row differs from its local flat row
by at most \(C_L\kappa n^{-4}\) per fine input, by summing
the relative one-step distortions
\(O_L(\kappa L^{i-j})\) over its backward scales. This is the
straight-row comparison on the complete carriers. The same
normalization bounds row and column sums. Each seed affects
only a bounded number of rows at its own and next level,
and each such row meets at most \(C_L\) seeds.
The entrywise errors thus give a \(C_L\kappa\) perturbation
in \(Y_0\). The leading next-level response is bounded by
the scalar \(Q_L^0\) weights, since the fixed colour
rotations have norm one. Hence
\(\|QE^0-I\|_{Y_0\to Y_0}\le L^{-1}+C_L\kappa<1\).
The seed energy estimate is unchanged up to a fixed constant:
the connection term in its first difference costs at most
\(C_L\kappa/\xi_j\).

The formula \(E_Q=E^0(QE^0)^{-1}\) now proves the result
on real fields. For a complex shift, keep the real seeds
fixed. On a smaller tube the analytic row coefficients
change by \(C_L\varepsilon n^{-4}\) per input on the same
carriers. This follows by Cauchy's formula applied to the
bounded normalized coefficient on the fixed physical
perturbation ball. The contraction reserve therefore
persists and the lift is holomorphic. The identities
\eqref{mb:eq-restoration-properties} give \(QE_L=I\) with
horizontal range. Testing a row covector on \(E_L b\),
where \(b\) attains its dual norm, proves the inf-sup
bound \(C^{-1}\) in \(X_L\). Short periods are handled
on a fixed large cover by averaging the seed lift over
simultaneous deck translations before solving \(QE^0\).
All the norms and the right inverse equation descend.
\end{proof}


% Construction source: assembly/site-pointwise.tex
\subsection{Site Green bounds for a bounded cell source}

\begin{proposition}[Local norms of the site Green operator]
\label{mb:prop-site-pointwise}
Use the true graph, actual site average and real curvature
hypotheses of Proposition~\ref{mb:prop-site-gram}. If \(f\)
is supported on one pillar \(B_z\) and \(u=G_s f\), then
\begin{equation}\label{mb:eq-site-pointwise}
 \xi_y^{-2}\|u\|_{\infty,B_y}
       +\|\Delta_Uu\|_{\infty,B_y}
               \le C e^{-m d(y,z)}\|f\|_\infty .
\end{equation}
In the regular local coefficient class, for \(0<\alpha<1\)
the left side may also include
\(\xi_y^{-1}\|D_Uu\|_{\infty,B_y}
 +\xi_y^{\alpha-1}[D_Uu]_{\alpha,y}^{U}\).
The first-difference stencils and transported H\"older pairs
use fixed comparable-scale enlargements of \(B_y\), inside
the complete local stock. At a bare boundary use the present
stencils. With this local coefficient class the estimates
continue to the admitted weighted \(C^1\) complex tube,
where the norms use the fixed real frames. The rate is
positive with a fixed reserve for the scale ratios used below.
\end{proposition}

\begin{proof}
The actual mass term \(q^{\mathsf T}a_s q u\) at a site of
\(B_w\) has norm at most \(C\xi_w^{-4}\|u\|_{2,B_w}\):
the transported average costs \(\xi_w^{-2}\) by
Cauchy--Schwarz, and \(a_s\) supplies \(\xi_w^{-2}\).
Apply the real scalar modulus estimate of
Lemma~\ref{mb:lem-scalar-subsolution} on an interior ball
of radius comparable to \(\xi_y\), with its neighbouring
whole pillars retained for this mass. The scales there
are comparable by separation. Its bound is
\(C\xi_y^{-2}\sum_w\|u\|_{2,B_w}
  +C\xi_y^2\|f\|_{\infty,\rm loc}\) over a fixed enlarged stock.
At a bare boundary the point-value bound directly costs
\(C\eta^{-2}\|u\|_{2,\rm loc}\).

The two exponential energy blocks of
Proposition~\ref{mb:prop-site-gram}, and
\(\|f\|_{2,B_z}\le \xi_z^2\|f\|_\infty\), give a bound
\(C\xi_y^{-1}\xi_z^3 e^{-s_0 d(y,z)}\|f\|_\infty\)
for the first term. By \eqref{mb:scale-ratio}, the
factor \((\xi_z/\xi_y)^3\) is absorbed at a smaller
positive rate, using the prescribed separation.
A local force contributes only at bounded metric distance
from \(z\). The equation then bounds \(\Delta_Uu\)
from the local \(L^\infty\) bound on \(u\), proving
\eqref{mb:eq-site-pointwise}.

For the differentiated assertion, apply
Proposition~\ref{mb:prop-local-scalar} in the declared
local frames at radius comparable to \(\xi_y\).
The normalized input energy
\(\xi_y^{-2}\|u\|_{2,\rm loc}+\xi_y^{-1}\|D_Uu\|_{2,\rm loc}\)
has exactly the preceding bound. The mass term is
zeroth order there, with \(L^p\) norm bounded by
\(C\xi_y^{-2}\|u\|_{p,\rm outer}\) on its complete
carrier. For any fixed \(\alpha<1\) choose
\(p>4/(1-\alpha)\). The result bounds the scaled first
differences and their transported H\"older quotient.
In the finitely coarse bare collars present-stencil
differences and H\"older quotients follow from the
point-value bound and \(\xi_y=\eta\).
For complex coefficients use the continued energy blocks
and the complex local bootstrap in the same real frames.
The mass and scale estimates are unchanged up to constants
because the complex transport condition numbers are
uniformly bounded on their complete carriers.
The equation again supplies the Laplacian bound.

In particular, for the normalized site column \(F_z\)
of \(F\) in a physical orthonormal top colour basis,
the right side in \(A_s F_z=q^{\mathsf T}\Xi_s^{-2}e_z\)
has supremum norm \(C\xi_z^{-4}\) and is supported in
\(B_z\). Here \(e_z\) is a unit top vector with value
of magnitude \(\xi_z^{-2}\). Thus
\[
 \|F_z\|_{\infty,B_y}+\xi_y^\alpha[F_z]_{\alpha,y}^{U}
                      \le C_\alpha \xi_y^{-2}e^{-m' d(y,z)}
\]
in the regular class. This follows from the derivative
bound and the bounded metric diameter of the cell,
and spends an additional \(4\log L/\rho\) to replace
\(\xi_y^2\xi_z^{-4}\) by \(\xi_y^{-2}\).
The estimate is for a column with its physical
normalization, rather than for a point-mass fine source.
\end{proof}


% Construction source: outline/site-divergence.tex
\begin{corollary}[Site Green operator on a cell flux]
\label{mb:cor-site-divergence}
Use the regular local coefficient class and the actual
graph of Proposition~\ref{mb:prop-site-pointwise}. Let
\(F\) be a bounded field supported on the bonds incident
to a pillar \(B_z\). Put
\(u=G_sD_U^{\mathsf T}F\), with the divergence in the
physical variational pairing. For \(2<p<\infty\) and
\(0<\alpha<1\), one has
\begin{equation}\label{mb:eq-site-divergence-point}
 \xi_y^{-1}\|u\|_{\infty,B_y}
 +\xi_y^{\alpha-1}[u]_{\alpha,y}^{U}
 +\xi_y^{-4/p}\|D_Uu\|_{p,B_y}
       \le C_{p,\alpha}e^{-m d(y,z)}\|F\|_\infty .
\end{equation}
If \(F\) also has the scaled \(C^\alpha\) bound on the
complete lifted source stock, then
\begin{equation}\label{mb:eq-site-divergence-gradient}
 \|D_Uu\|_{\infty,B_y}+\xi_y^\alpha[D_Uu]_{\alpha,y}^{U}
 \le C_\alpha e^{-m d(y,z)}
               \bigl(\|F\|_\infty+\xi_z^\alpha[F]_{\alpha,z}^{U}\bigr).
\end{equation}
Here the source stock includes the zero values of the
supported flux on the collars used by any cutoff meeting
it. In particular its seminorm includes increments to
those zero values. Output seminorms and bare stencils use
the same conventions as Proposition~\ref{mb:prop-site-pointwise}.
Both statements hold on the same type of smaller complex
coefficient tube, with a fixed positive reserve in the rate.
\end{corollary}

\begin{proof}
The weighted site energy estimate in the proof of
Proposition~\ref{mb:prop-site-gram} applies with the dual
source \(D_U^{\mathsf T}F\). Indeed, testing the divergence
against a site field uses one difference of the test.
Under exponential conjugation its additional multiplication
coefficient is \(O(|s|\ell^{-1})\), by the difference
product rule across each bond. This is also controlled
by the site energy. Taking the distance weight from the
source stock gives, on every fixed enlarged output pillar,
\[
 \xi_y^{-1}\|u\|_{2,\mathrm{loc}}
       +\|D_Uu\|_{2,\mathrm{loc}}
   \le C e^{-s d(y,z)} \xi_z^2\|F\|_\infty .
\]
The physical volume of the complete flux stock accounts
for \(\xi_z^2\). Thus the rescaled energy input in
Proposition~\ref{mb:prop-local-scalar} is bounded by
\(C\xi_y(\xi_z/\xi_y)^2 e^{-s d(y,z)}\|F\|_\infty\).
The strict scale reserve absorbs the ratio with loss at most
\(2\log L/\rho\) in the rate.

On a complete stock the equation is
\((\Delta_U+q^{\mathsf T}a_s q)u=D_U^{\mathsf T}F\).
The mass is an admitted zeroth-order term for that
proposition on each cutoff support, with its complete
pillar inputs. The rescaled \(L^p\) flux norm on those
stocks is at most \(C\xi_y\|F\|_\infty\), and it vanishes
unless the stock is at bounded metric distance from \(z\).
Apply the divergence estimate for the requested \(p\).
For the value and its H\"older bound take a finite
exponent \(p_1>4/(1-\alpha)\) and use the first-order Morrey
estimate. This proves \eqref{mb:eq-site-divergence-point}.

For \eqref{mb:eq-site-divergence-gradient}, use the
H\"older-flux version of the same bootstrap, with the
mass term still estimated as an ordinary finite-\(p\)
source. When the local flux is nonzero, the two scales
are comparable and the rescaled \(C^\alpha\) flux norm
is at most \(C\xi_y(\|F\|_\infty+\xi_z^\alpha[F]_{\alpha,z}^{U})\).
Undoing the output scaling yields the bound. At a bare
boundary the corresponding atomic energy estimates cost
\(\eta^{-2}\) for a point value and \(\eta^{-\alpha}\)
for a H\"older quotient. The displayed physical weights
cancel these factors. The complex proof uses the continued
weighted site energy estimate and the bilinear local
divergence equation. Its local estimates use the fixed
real comparison frames.
\end{proof}


% Construction source: assembly/layered-gap.tex
\section{Coercivity on separated determining regions}
\label{mb:sec-layered-gap}

We use the covariant curl \eqref{mb:curl} on the inner plaquettes,
with physical measure \(\eta^4\).

Let \(H_W(U)\) be the Wilson Hessian in the physical left
exponential coordinates. For the actual per-length nonlinear
straightened map \(\mathcal Q_U\), put
\begin{equation}\label{mb:eq-constraint-form}
 \begin{split}
 B_\omega(v,w)&=\sum_c\xi_c^4
    \langle\omega(c),D^2\mathcal Q_{U,c}(0)[v,w]\rangle,\\
 \gamma_\omega&=\sup_c\xi_c^3|\omega(c)|.
 \end{split}
\end{equation}
The fixed-input-leg moment of the actual averaging derivative
is \(C_L\xi_c(\eta/\xi_c)^4\). Pairing this moment with
\(\xi_c^4|\omega(c)|\) leaves at most
\(C_L\gamma_\omega\eta^4\xi_c^{-2}\) on each fine input.
The carriers have bounded overlap and their scales are
comparable to the fine scales they read. By symmetry the same
bound holds on either leg. Weighted Schur then gives
\begin{equation}\label{mb:eq-constraint-small}
             |B_\omega(v,w)|\le C_L\gamma_\omega
                                         \|v\|_{\rm m}\|w\|_{\rm m}.
\end{equation}
The argument applies on the admitted complex chart. The
estimate is a coefficient estimate for all input vectors,
even though it was obtained on a small analytic ball.

For real curvature-small \(U\) the Wilson Hessian satisfies
\begin{equation}\label{mb:eq-wilson-form}
 \begin{aligned}
 |\langle v,H_W(U)w\rangle_f-\langle\mathrm d_Uv,\mathrm d_Uw\rangle_f|
                                    &\le C_L\kappa\|v\|_X\|w\|_X,\\
 \langle v,H_W(U)v\rangle_f&\ge
    (1-C_L\kappa)\|\mathrm d_Uv\|_2^2-C_L\kappa\|v\|_{\rm m}^2 .
 \end{aligned}
\end{equation}
Indeed varying a plaquette inserts four \(\eta v\)'s transported
by its successive prefixes. The trace Hessian with the final
plaquette matrix replaced by \(I\) is the pairing of their
sums, since the trace kills the commutator parts. In the
difference each term contains the final plaquette defect
\(O(\kappa\eta^2\mathsf L^{-2})\) and two factors \(\eta\).
Their product supplies the physical measure and the mass weight.
The prefix curl and \(\mathrm d_U\) differ only by a coefficient
\(\eta^{-1}(I-\operatorname{Ad}U(\partial p))\) on two
of the link values, as follows by multiplying the four prefixes
and using the reversed-link rule. Its norm is at most
\(C_L\kappa\|v\|_{\rm m}\) because \(\eta\le\mathsf L\).
Cauchy--Schwarz and the bounded number of plaquettes at a
link give \eqref{mb:eq-wilson-form}. In particular no
boundary-crossing action term has entered this estimate.

\mbneedspace{10\baselineskip}
\subsection{The flat decomposition}

\begin{lemma}[A regular decomposition at a real flat connection]
\label{mb:lem-flat-decomposition}
On the true active graph at a real flat connection \(V\),
with the separated no-inward geometry and the actual site and
field rows, every field has a decomposition
\[
 v=D_V\lambda+r,\qquad q\lambda=0,\qquad
             \|r\|_X\le C(\|\mathrm d_Vv\|_2+\|Qv\|_Y).
\]
The constant is independent of topology, flat holonomies,
depth and compatible volume.
\end{lemma}

\begin{proof}
First we relate the outward average to differences of
individual pillar values. Two pillars sharing a fine face
have equal or adjacent levels. For an equal-level pair the
coarse difference is its interior row. For an unequal pair
the larger scale is \(s\). In a parallel frame on their
two-scale stock, the larger value \(a\) and the \(L^4\)
smaller values \(b_w\) in the whole outer block enter the
outward row as \(s^{-1}(a-\bar b)\), with \(\bar b\) their
arithmetic mean. All smaller rows internal to that block
are retained. For each face neighbour \(w\), write
\(a-b_w=(a-\bar b)+(\bar b-b_w)\).
The finite \(L\)-grid Poincar\'e inequality estimates the
second difference by the internal smaller differences.
Conversely choose one face neighbour \(w_0\) and write
\(a-\bar b=(a-b_{w_0})+L^{-4}\sum_w(b_{w_0}-b_w)\).
The last sum uses the same smaller rows along coordinate
paths of length at most \(4L\). Thus a route through a
bounded number of comparable pillars is controlled by
its actual interior and outward differences, together with
the internal rows of each full outer block. The length
factors on that route differ by bounded powers of \(L\).
This proof on lifts retains all faces, even if a periodic
pair of pillars has more than one shared face.

Next consider a convex fine box \(D\) of side lengths
comparable to \(s\). In a parallel frame project the field
orthogonally onto gradients of site functions on that box.
The remainder \(r_D\) is orthogonal to gradients, including
the boundary site equations, and has curl equal to the
given field's curl. In one dimension with \(N+1\) vertices,
the scalar path Neumann matrix has eigenvalues
\(4\eta^{-2}\sin^2(\pi m/[2(N+1)])\), \(0\le m\le N\).
This follows by substituting the cosines with half-site
offset. The edge matrix has exactly its nonzero spectrum,
since the incidence map is onto edges. The cochain Hodge
matrix on a box is the sum of those tensor factors, using
the edge matrix in the component direction and the scalar
matrices in the others. This identity follows by expanding
the curl and divergence squares, whose mixed terms cancel.
Each component has its nonzero edge frequency. Its
quadratic form controls the interior component differences
and its lowest eigenvalue is at least \(cs^{-2}\).
Therefore
\begin{equation}\label{mb:eq-local-hodge}
       s^{-1}\|r_D\|_{2,D}+\|\nabla_Vr_D\|_{2,D/2}
                            \le C\|\mathrm d_Vv\|_{2,D}.
\end{equation}
The gradient potential can be given zero average on a
specified central pillar of side comparable to \(s\) by
subtracting a parallel constant.

Choose a fixed large \(K=K(L)\) and take \(\rho\) larger
than all the required \(C_LK\) collars. A box \(D_y\) of
radius \(K\xi_y\) about a pillar top can meet at most two
adjacent assigned levels. Otherwise it would cross an
intervening layer of width greater than its radius, by
separation. Such boxes and fixed enlargements have bounded
overlap, by counting tops of comparable scales. Smooth
sampled bumps supported in the doubled pillar cubes and
equal to one on their pillars, divided by their sum, give
a partition \(\sum_y\chi_y=1\), with first two difference
bounds \(C_L\xi_y^{-1}\) and \(C_L\xi_y^{-2}\). Indeed the
sum is between one and the fixed overlap bound, and each
term contributing to a stencil has comparable scale.

At the true boundary, first write \(v=v_b+w\), where
\(v_b\) equals \(v\) within \(H\eta\) of the exterior and
vanishes beyond \(2H\eta\), with \(H>10K\) fixed.
Take \(\rho>4H+C_L\). All its nonzero bonds are bare
determining rows, where \(Q_0v=v\), and no higher carrier
meets that support. The cutoff difference costs at most
\(C\eta^{-1}\) there. Thus
\[
 \|v_b\|_X+\|\mathrm d_Vw\|_2+\|Qw\|_Y
                           \le C(\|\mathrm d_Vv\|_2+\|Qv\|_Y).
\]
A nonbare box is inside the active region by separation.
A bare box meeting the exterior lies entirely in the
zero collar of \(w\), where its decomposition is set to
zero. Hence every nonzero local decomposition uses only
the flat connection on an inner rectangular box.

On \(D_y\) write \(w=D_V\lambda_y+r_y\) with
\(q_y\lambda_y=0\) using \eqref{mb:eq-local-hodge}.
If two cutoff supports or their two-step collars intersect,
their scales are comparable and their centres are within
\(C_L\sigma\), where \(\sigma=\max(\xi_y,\xi_z)\).
For \(K\) chosen large enough their half-boxes intersect
in a neighbourhood containing the central pillars and a
coordinate route between them. Only \(C_L\) pillars of
comparable sizes meet this route. Enlarge it to include the
complete outward stocks used in the first paragraph;
this fits into the intersection by the choice of \(K\).
On this route \(\overline D q\lambda_z=Q(w-r_z)\) by
\eqref{mb:extended-gradient-identity}. Its actual site value
at \(z\) is zero. The route estimate gives
\[
  |q_y\lambda_z|\le C e_{yz},\qquad
  e_{yz}^2=\sum_{c\ {\rm in\ route}}\xi_c^2|Qw(c)|^2+
                                  \|\mathrm d_Vw\|_{2,D_y\cup D_z}^2 .
\]
For the remainder contribution, the per-length row has
physical \(L^2\) functional norm at most \(C\xi_c^{-2}\),
by the product of its coefficient sum and its maximum
coefficient divided by \(\eta^4\). Equation
\eqref{mb:eq-local-hodge} bounds \(\xi_c|Qr_z(c)|\)
by \(C\|\mathrm d_Vw\|_{2,D_z}\). This proves the stated
route estimate. If a box meets the exterior, all boxes whose
cutoffs overlap its cutoff are bare and remain inside the
zero collar of \(w\), by \(H>10K\) and separation. Their
potentials and remainders are zero. Every other pair
neighbourhood used in the estimate is inner.
Poincar\'e with the \(y\)-pillar mean on
this comparable-volume intersection now gives
\[
 \|\lambda_y-\lambda_z\|_2\le C\sigma^2 e_{yz},\qquad
 \|D_V(\lambda_y-\lambda_z)\|_2\le C\sigma e_{yz}.
\]
The gradient bound comes from \(r_z-r_y\). For the first
bound the mean costs \(\sigma^2|q_y\lambda_z|\) in physical
\(L^2\). Bounded overlap gives
\(\sum e_{yz}^2\le C(\|\mathrm d_Vw\|_2^2+\|Qw\|_Y^2)\).

Put \(\lambda_*=\sum_y\chi_y\lambda_y\).
In \(w-D_V\lambda_*\) the derivatives of the cutoff multiply
only differences \(\lambda_y-\lambda_z\), because their
sum is zero. One more difference produces at most
\(\sigma^{-2}\|\lambda_y-\lambda_z\|_2+
 \sigma^{-1}\|D_V(\lambda_y-\lambda_z)\|_2\)
for each overlap. The remaining terms cost the two terms
in \eqref{mb:eq-local-hodge}. The same estimates control
the mass norm. Thus the \(X\) bound for the glued remainder
follows by bounded overlap. Since \(q_y\lambda_y=0\), the
glued mean \(q_y\lambda_*\) only reads those same potential
differences. The physical \(L^2\) norm of its averaging
functional is \(\xi_y^{-2}\), which cancels the factor
\(\sigma^2\). Consequently
\[
          \sum_y|q_y\lambda_*|^2
                       \le C(\|\mathrm d_Vw\|_2^2+\|Qw\|_Y^2).
\]
At a flat field the compact scalar bump lifts of the site
averages obey
\(\|D_V E_s b\|_X^2\le C\sum_y|b_y|^2\):
the first weighted and second ordinary differences have
pointwise size \(C\xi_y^{-2}|b_y|\) on volume
\(O(\xi_y^4)\). Bare atomic seeds give the same estimate.
Subtract \(E_s q\lambda_*\) from the potential and add its
gradient and \(v_b\) to the remainder. This gives the
decomposition.

For short periods apply the construction on a compatible
cover on which these finitely sized boxes embed. Deck
averaging the potential preserves its zero site means and
the norm bound for its remainder. The input and resulting
remainder norms have the same sheet-count factor. This
uses no global choice of flat frame.
\end{proof}

\subsection{The real and complex form gaps}

\begin{theorem}[Uniform layered form gaps]\label{mb:thm-layered-gap}
On the separated aligned geometry of \eqref{mb:separation},
use the real curvature class \eqref{mb:curvature-class}
with \(\kappa\) sufficiently small and the positive weights
of \eqref{mb:eq-row-energy} and \eqref{mb:eq-site-projector}.
The actual forms
\begin{equation}\label{mb:eq-layered-forms}
 \begin{split}
 K_0&=H_W(U)+D_URD_U^{\mathsf T}+Q^{\mathsf T}a_fQ,\\
 K_F&=H_W(U)+D_UD_U^{\mathsf T}+Q^{\mathsf T}a_fQ
 \end{split}
\end{equation}
satisfy \(c\|v\|_X^2\le\langle v,K_i v\rangle_f\le C\|v\|_X^2\)
for \(i=0,F\). The same is true after subtracting the form
\(B_\omega\) from \(K_0\) when
\(\gamma_\omega\) is sufficiently small. The constants are
uniform in mesh, depth, shapes and compatible volume.
The corresponding normalized energy inverses are holomorphic
and uniformly bounded on a fixed smaller weighted \(C^1\)
complex tube, with the actual constraints and projections
continued there.
\end{theorem}

\begin{proof}
At a real flat reference \(V\), use
Lemma~\ref{mb:lem-flat-decomposition}. Its potential satisfies
\[
 \Delta_V\lambda=R D_V^{\mathsf T}(v-r)
\]
by the range identity \eqref{mb:eq-restoration-properties}.
Equation~\ref{mb:eq-zero-site-H2} yields
\(\|D_V\lambda\|_X\le C(\|R D_V^{\mathsf T}v\|_2+\|r\|_X)\).
Together with the remainder bound this proves the \(K_0\)
lower bound at \(V\). The upper bound follows from the
local differences and the physical Schur bound for \(Q\).

Suppose \(U\) and a real flat reference \(V\) on the same
geometry are in a common gauge with
\(U_b=e^{i\eta A_b}V_b\) and
\(\varepsilon=\sup_b\mathsf L(b)|A_b|\) small.
The local comparison of the actual rows and the shifted
gradient coefficients gives
\[
 \begin{split}
 \|\ell(D_U-D_V)u\|_2 &\le C\varepsilon\|u\|_2,\\
 \|(q_U-q_V)u\|_{2,s}
 +\|\Xi_s^{-2}(q_U-q_V)(\ell^2u)\|_{2,s}
                                      &\le C\varepsilon\|u\|_2,\\
 \|(Q_U-Q_V)v\|_Y&\le C\varepsilon\|v\|_{\rm m}.
 \end{split}
\]
For the row bounds pass to a flat frame on each complete
carrier and use the averaging coefficient estimates with
the shift as parameter. An entry of a per-length row
changes by \(C_L\varepsilon(\eta/\xi_c)^4\), since its
relative perturbation at each intermediate scale is
\(O_L(\varepsilon L^{i-j})\) for a level-\(j\) row.
The geometric sum and the bounded carrier overlap yield
the stated physical norms. The same transported average
comparison proves the \(q\) bounds.

The site column difference solves
\[
 A_{s,U}(F_U-F_V)b
  =(q_U^{\mathsf T}-q_V^{\mathsf T})\Xi_s^{-2}b
                    -(A_{s,U}-A_{s,V})F_Vb .
\]
Test it with \(\ell^2(F_U-F_V)b\).
The left side is the plus-weight energy in the proof of
Proposition~\ref{mb:prop-site-gram}. In the gradient
difference one \(D\) is replaced by an \(O(\varepsilon
\ell^{-1})\) coefficient; the remaining derivative and
the weight are controlled by the plus-weight energies
using the product rule. The mean difference uses the
two site-row bounds just proved. It follows that
\(F_U-F_V\) and its \(\ell\)-weighted first derivative
are \(O(\varepsilon)\) in their physical operator norms.
The Gram inverse identity then gives
\(\|R_U-R_V\|_{2\to2}\le C\varepsilon\).
Also
\(\|(R_U D_U^{\mathsf T}-R_VD_V^{\mathsf T})v\|_2
 \le C\varepsilon\|v\|_X\),
since the adjoint coefficient difference is bounded by
\(C\varepsilon\|v\|_{\rm m}\).
Finally \(\mathrm d_U-\mathrm d_V\) obeys the same mass
bound. Equations~\ref{mb:eq-wilson-form} and
\eqref{mb:eq-layered-forms} give a form difference
\(C(\kappa+\varepsilon)\|v\|_X\|w\|_X\). This proves the
gap in the common gauge by absorption in the flat form.

To remove the global small-potential condition, localize on
the original graph. We give the two estimates needed for
localization. Write \(P=I-R\). The mass-normalized matrix
\[
 K_s^{\rm sm}=\ell D_U P D_U^{\mathsf T}\ell
     =(\ell D_UF)M^{-1}(\ell D_UF)^{\mathsf T}
\]
has exponential cell-\(L^2\) entries by
Proposition~\ref{mb:prop-site-gram}. By taking \(\rho\)
large with fixed reserve in \eqref{mb:rate-reserve},
these entries have summable polynomial moments on the
layer metric, in particular their second moment.

Choose a fixed large radius \(R_*>1\), then increase
\(\rho\) so that balls of radius \(C_L R_*\xi_y\) meet
at most two adjacent scales. There is a squared partition
\(\sum_i\chi_i^2=1\) with \(\chi_i\) supported on balls of
radius \(2R_*\xi_i\), bounded overlap, and
\[
 \sum_i|D\chi_i(x)|^2\le C R_*^{-2}\ell(x)^{-2},\qquad
 \sum_i|\chi_i(x)-\chi_i(z)|^2\le
                       C\min\{1,d_\ell(x,z)^2/R_*^2\}.
\]
To construct it, select centres in decreasing scale
with disjoint balls of radius \(R_*\xi_i/(20L)\)
until maximal. Their enlarged balls cover by
comparability of scales where they meet. Packing
gives bounded overlap. Use smooth bumps equal to one
on the covering balls and divide by the square root of
their squared sum. Include bare exterior cells for this
construction if the active region is proper. Integrating
the vector of cutoff derivatives along a metric path
proves the second bound. Within each cell the sum of
the squared suprema of the changes from its centre
is \(O(R_*^{-2})\), by the derivative bound and the
bounded number of overlapping supports.

The localization error of \(K_0\) is bounded by
\(CR_*^{-2}\|v\|_{\rm m}^2\). In fact for a symmetric
matrix \(K\) its error is the double cutoff
\(-\tfrac12\sum_i[\chi_i,[\chi_i,K]]\).
For \(K_s^{\rm sm}\), on a cell pair expand each cutoff
difference into its centre difference and the two
within-cell differences. Each term of the double
commutator is left and right multiplication by the
corresponding scalars or cutoff deviations. By
Cauchy--Schwarz in \(i\), its block operator norm is
bounded by \(CR_*^{-2}(1+d_\ell(y,z))^2\) times the
entry norm. Summation uses the second moment, with
\(\ell^{-1}v\) as input to the normalized smoothing
matrix. The remaining fine-range Wilson and divergence
matrices have coefficients \(O(\eta^{-2})\) in physical
array form. Their squared cutoff differences cancel
that denominator, since \(|D\chi_i|\) costs
\(C R_*^{-1}\ell^{-1}\).
The row mass has uniformly bounded normalized coefficient
Schur sums on its bounded metric carriers, so the same
double-cutoff estimate applies. The local shift formula
for \(\|v\|_X^2\) gives the same localization error.
A zero-extension jump is exact since it has only one
nonzero value.

We also need that changing a connection outside a large
patch changes its compressed projector exponentially little.
Suppose two admitted real fields agree on the cube of
radius \(24R_*\xi_i\) in a common gauge, with comparable
scales there, and have the same determining geometry.
For a cutoff \(\theta\) with gradient in an intermediate
annulus, the exact identity
\[
 \theta G_{s,1}-G_{s,2}\theta
           =G_{s,2}(A_{s,2}\theta-\theta A_{s,1})G_{s,1}
\]
has middle factor supported in the annulus with the
site energy-to-dual-energy bound. Its terms are the
gradient commutator and the two actual site-row cutoff
differences, since the coefficients agree wherever
\(\theta\) or its required neighbours are nonzero.
The exponential conjugated energy estimates of
Proposition~\ref{mb:prop-site-gram} make the relevant
tails \(O(e^{-cR_*})\) in that norm. Here a path leaving
the patch from its core must pay length at least
\(cR_*\) before leaving its comparable-scale cube.
The fine inverse carries its two factors of \(\xi_i\);
for \(F\) the source \(\Xi_s^{-2}\) cancels them.
For top input on a middle cube of radius \(8R_*\xi_i\)
this gives an \(O(e^{-cR_*})\) bound on the difference of
the two columns in the interior. Their outputs outside
a larger middle cube have the same small norm by the
column tail estimate. The two Gram matrices compressed
to the middle top cube therefore differ in norm by
\(O(e^{-cR_*})\).

Both compressed Gram matrices retain the uniform gap.
The difference between the compressed inverse and the
compression of the global inverse, applied to input on
the inner cube, is the product of the compressed inverse,
the off-diagonal Gram block, and the exterior tail of
the global inverse. This follows by restricting
\(MM^{-1}=I\) to the middle cube. The exterior tail is
\(O(e^{-cR_*})\) by the weighted inverse bound. Thus
inserting \(F M^{-1}F^{\mathsf T}\), with its analogous
column and transpose tails, gives
\begin{equation}\label{mb:eq-projector-shield}
 \|1_{\rm core}(R_1-R_2)1_{\rm core}\|_{2\to2}
                                                    \le C e^{-cR_*}
\end{equation}
on the cube of radius \(3R_*\xi_i\). These estimates
are on supported vectors and require no sum of errors
over the number of fine sites.

For an interior localization patch choose a comb frame
on the cube of radius \(32R_*\xi_i\), then extend its
potential by a scalar cutoff on that cube, equal to one
on the agreement cube. Keep the original graph and all
its assigned averages. The comb estimate gives potential
size \(C_LR_*\kappa/\xi_i\) there. The cutoff derivative
is \(O((R_*\xi_i)^{-1})\). Expanding the plaquette
product shows that the new curvature is at most
\(C_L(\kappa+R_*^2\kappa^2)/\xi_i^2\): the linear
plaquette difference comes from the original curvature
and the cutoff derivative, and the quadratic product
costs \(C_LR_*^2\kappa^2/\xi_i^2\).
Thus the comparison field is globally in a small
potential gauge with the identity on the same graph.
For \(v_i=\chi_i v\), the inner Wilson and actual row
forms agree exactly with this comparison field, by their
complete stocks. The divergence is supported in the core,
so \eqref{mb:eq-projector-shield} changes its gauge form
by at most \(Ce^{-cR_*}\|v_i\|_X^2\).
We used orthogonality of \(R\) to write the form as
\(\langle D_U^{\mathsf T}v_i,R D_U^{\mathsf T}v_i\rangle\).

If a bare patch's enlarged comparison cube would meet
the exterior, its entire support is in the bare
boundary collar by the larger separation. On that
support \(Qv_i=v_i\), higher carriers miss it, and the
row mass is at least \(d_*\eta^{-2}\|v_i\|_2^2\).
The possible negative Wilson remainder is
\(C\kappa\eta^{-2}\|v_i\|_2^2\) by
\eqref{mb:eq-wilson-form}. The direct bound
\(\|\nabla_U v_i^0\|_2\le C\eta^{-1}\|v_i\|_2\)
therefore gives the required gap on these patches.
Choose \(R_*\) so the two localization errors are
small relative to the fixed flat gap, then choose the
separation for the patches and the rate reserve.
Finally reduce \(\kappa\) to admit all comparison
fields in the common-gauge gap. Summing the local
bounds proves the \(K_0\) lower bound on the real
curvature class. The same argument on a compatible
cover descends to short periods, since the site
and Gram inverses preserve the deck-invariant space.

The upper bounds for both forms follow from their
local differences, the row bound and \(\|R\|\le1\).
Also \(K_F-K_0=D_UPD_U^{\mathsf T}\) is nonnegative
at a real reference. Equation~\ref{mb:eq-constraint-small}
gives the asserted multiplier perturbation. On the
complex tube, the actual \(R,D,Q\) comparison from
Proposition~\ref{mb:prop-site-gram} and the analytic
averaging bounds give an \(O(\varepsilon)\) form
difference in \(X\). For the Wilson coefficients, the
\(C^1\) weighted shift changes each plaquette by
\(O(\eta^2\varepsilon\mathsf L^{-2})\) and each prefix
transport by \(O(\eta\varepsilon\mathsf L^{-1})\).
The calculation for \eqref{mb:eq-wilson-form} gives
the same \(X\)-form perturbation. A relative Neumann
series about the real positive form proves bounded
holomorphic inverses. The complex radii are chosen
after the real uniform gap.
\end{proof}

\begin{corollary}[The constrained Lagrangian saddle]\label{mb:cor-physical-saddle}
Under Theorem~\ref{mb:thm-layered-gap}, with
\(\gamma_\omega\) small, the Lagrangian Hessian
\(H_L=H_W-B_\omega\) is coercive on
\(X_L\cap\ker Q\). Given \(b\in Y\) and a force \(f\) in
the dual of \(X_L\), the equations on \(X_L\)
\[
                 H_L v-Q^{\mathsf T}\mu=f,\qquad Qv=b
\]
have a unique solution with
\(\|v\|_X+\|\mu\|_{Y^*}\le C(\|f\|_{X_L^*}+\|b\|_Y)\).
The conclusion concerns the specified inner action
and actual row curvature at the given pair.
\end{corollary}

\begin{proof}
On the simultaneous kernel the two positive penalties
in \(K_0\) vanish exactly, so the form gap gives the
coercivity. Write \(v=E_L b+z\) using
Lemma~\ref{mb:lem-field-lift}, solve the coercive
equation on \(z\in X_L\cap\ker Q\), and define
\(\langle\mu,b'\rangle_r=(H_Lv-f)(E_Lb')\).
Every vector of \(X_L\) is the sum of its \(E_LQ\)
part and a vector of the simultaneous kernel. This
gives the full saddle and the estimate. Testing the
difference of two solutions on the kernel and then
on the lift proves uniqueness.
\end{proof}


% Construction source: assembly/compressed-gap.tex
\subsection{Ambient compression and the flat tree bound}

Let the real connection be specified also on the complete
ambient plaquettes touching the active graph. Compress its
Wilson Hessian by zero extension of the field variables.
The site operator may use its corresponding Dirichlet
compression, or the original graph Laplacian. For the
Dirichlet convention include its boundary diagonal in
the site form. We retain the same complete determining
rows and their physical pairings. Assume the scaled
curvature bound also on the crossing plaquettes.

The alternative single-comb averages use the full differential
of the relative average in \cite[eqs.~(89)--(92), (124)]{Bal85a},
with its matching centred site average
\[
 q_{\rm comb}\lambda(y)
     =L^{-4}\sum_{x\in B_y}
             \operatorname{Ad}U(\Gamma_{yx})\lambda(x),
\]
where $\Gamma_{yx}$ is the chosen comb path from $y$ to $x$.
Here is the field average's one-step definition. For a coarse bond $c$, let $\Gamma_{c,x}$ run
from its initial centre along the comb to $x$, along the
translated long segment to $x'$, and back along the terminal
comb. Write $W(c)$ for the central long segment. The relative
mean is
\[
 \overline W(c)=
 \exp\!\left(L^{-4}\sum_x
       \log\{W(\Gamma_{c,x})W(c)^{-1}\}\right)W(c).
\]
All loop logarithms use their admitted small branches.
For a fine tangent $a$, first solve the tree equations
$(a-D\psi)|_{\rm comb}=0$ with $q_{\rm comb}\psi=0$.
Tree integration determines $\psi$ up to a transported
constant on each block, which the last mean condition fixes.
Let $\overline W_t$ be this mean evaluated on
$(e^{i\eta ta}U)^{e^{-it\psi}}$, and put
\[
 Q_{\rm comb}a(c)=(iL\eta)^{-1}
    \left.\frac{d}{dt}\log\{
          \overline W_t(c)\overline W_0(c)^{-1}\}\right|_{t=0}.
\]
This full differential includes all differentiated transport
and endpoint factors. Iterate these actual maps
with the complete outward endpoint mean. These are the
single-comb maps used in the compressed background operator
of \cite[pp.~392--394, eqs.~(3.13)--(3.19)]{Bal85p}.

\begin{proposition}[The compressed zero-current estimates]
\label{mb:prop-compressed-gap}
Under the preceding separated geometry and coefficient bounds,
the analogues of \(A_s,M,K_F,K_0\) for the ambient compression
have the same uniform normalized real gaps and energy
inverse bounds. The row lift and zero-average restoration
remain uniformly bounded. These assertions also hold for
the full centred single-comb differential and its matching
site rows just defined. Its precise point, Laplacian,
divergence-source and H\"older output ranges are proved in
Corollary~\ref{mb:cor-compressed-regularity} below, after
the local source and field norms have been introduced.
\end{proposition}

\begin{proof}
The Dirichlet site form differs from the original one by a
nonnegative diagonal supported in the bare boundary collar.
Each diagonal coefficient is at most \(8\eta^{-2}\): it
counts the missing outside neighbours in the gradient form
of the zero extension. The real weighted energy tests and
the compact lift proof for the site Gram remain valid
with this term. It is nonnegative in the tested form
and bounded by the bare mass norm, also when paired with
the compact lift. It vanishes on \(\ker q\), because \(q\)
is the identity on bare sites. Hence
\(\Delta_{\rm Dir}(\ker q)=\Delta_U(\ker q)\);
the real orthogonal \(R\) is unchanged. The divergence of
an active field is the restriction of the ambient
divergence of its zero extension, which vanishes at all
exterior sites.

A crossing plaquette contributes its transported curl
square and a remainder as in \eqref{mb:eq-wilson-form},
using zero field values on its inactive links.
Every nonzero field value in these extra terms is in
the bare collar by separation. The principal squares
are nonnegative. The sum of the remainder absolute
values is bounded by \(C_L\kappa\eta^{-2}\) times the
squared \(L^2\) norm on that collar, whereas the sum of
the principal squares is bounded above by the same mass
bound without \(\kappa\). Absorbing the small remainder
in Theorem~\ref{mb:thm-layered-gap} proves the compressed
field gap and the upper bound. This argument retains
the crossing terms; it does not subtract them as a
small operator.

The single-comb rows agree at a flat field after
centring the blocks. In the same local real comb frames,
the summable one-step comparison gives \(O_L(\kappa)\)
differences in both physical row and column norms, with
the same $(\eta/\xi_c)^4$ coefficient density for field
rows. These are the full finite-path analytic row estimates;
see \cite[eqs.~(139)--(147), pp.~39--40; Proposition~7, p.~43]{Bal85a}.
For the iteration, write a per-length one-step derivative as
$S_i+E_i$ in a common real comb frame. Its absolute row
sum is $1+O_L(\delta_i)$, its column sum is
$L^{-4}(1+O_L(\delta_i))$, and each entry of $E_i$ is
$O_L(\delta_i L^{-4})$. The relative potential at that
step gives $\delta_i\le C_L\kappa L^{i-j}$ for a row
at final level $j$. Telescope the product, use the later
row sums and earlier column sums on either side of each
error entry, and sum this geometric sequence. This proves
both normalized error sums and the point density
$C_L\kappa(\eta/\xi_c)^4$, uniformly in the depth.
The common-frame site resolvent and Gram comparison
used in Theorem~\ref{mb:thm-layered-gap} applies to this
change of \(q\) at fixed connection, also in the presence
of the same bare diagonal. Thus \(R\) and the field
forms change by \(O_L(\kappa)\) in the corresponding
normalized norms. The seed lift proof applies with these complete row
coefficients. To check the exact gradient identity, put
$a=D\lambda$ in the defining tree correction. Then
$\lambda-\psi$ is parallel on every comb, and its value
at the centre is $q_{\rm comb}\lambda$ because
$q_{\rm comb}\psi=0$. Coarse endpoint covariance of the
relative mean therefore gives
\[
 Q_{\rm comb}D\lambda=
       \overline D_{\rm comb}q_{\rm comb}\lambda.
\]
At an outward row the complete further endpoint average
is included in $\overline D_{\rm comb}$. Iteration preserves
this identity and gives the corresponding restoration identities.

The conjugated energy estimates use the same gradient,
actual row and finite-range commutators as before;
the extra diagonal commutes with the weights and the
crossing Hessian uses only bare stencils. Away from the
bare boundary the complete local elliptic equations have
the same stencils and the same coefficients as in the
interior bootstrap. At the boundary the bare energy
column gives the point and derivative bounds by the
atomic estimates with \(\xi_y=\eta\). Applying the
interior site and zero-current component estimates with
these energy columns gives the stated point and
first-difference bounds. The site operator does not
acquire a fine Dirichlet boundary in a nonbare layer.
\end{proof}

For clarity the flat strict tree case has a lower bound
with an explicit dependence on \(L\). Take in each
retained block the tree whose paths from its centre
follow a single fixed coordinate order. If the field
vanishes on those trees and outside the retained stock,
then for its flat straight one-step constraint,
\begin{equation}\label{mb:eq-flat-tree-quantitative}
       S_La=0\quad\Longrightarrow\quad
                   \|a\|_2^2\le C L^5\|\mathrm da\|_2^2 .
\end{equation}
Here the raw \(S_L\) differs by the factor \(L\) from the
per-length straight row, which does not affect its kernel.
For an internal bond the difference of the two tree
paths and that bond bounds its value by a sum of at most
\(C L\) plaquette curls. Insert the bond step into the
fixed coordinate order and move it past the later
coordinate segments, forming strips of width one.
Each strip has no repeated plaquettes and there are at
most four strips. Cauchy--Schwarz and at most \(4L^4\)
internal bonds give the bound \(CL^5\) on a block.

On a face normal to \(\mu\), use the central crossing
as reference, with value \(\beta\). The loop consisting
of a tree path to a crossing link, that link, the return
tree path, and the central segment has a filling by at
most \(C L^2\) plaquettes in the double block. To see the
area and multiplicity, split the transverse coordinates
into those before and after \(\mu\) in the order. Move
the unit crossing step past the later transverse
segments, using width-one strips. The earlier transverse
segments then bound rectangles of width \(L\) along
the central segment. This uses at most six strips or
rectangles, each without repeated plaquettes.
Consequently each crossing value differs from \(\beta\)
by at most \(CL\|\mathrm da\|_{2,\rm double\ block}\).
The same bound holds for a complete translated
longitudinal segment: extrude its initial tree path
in direction \(\mu\) by \(L\); each transverse segment
sweeps one such rectangle, while the longitudinal
segment has zero swept area. The tree values and the
internal parts of the central segment vanish. The
zero straight average of the translated segments
bounds \(|\beta|\) by the same error. Squaring and
summing over \(L^3\) crossings gives \(CL^5\).
Summing over blocks and faces uses fixed overlap.
The proof also applies on rectangular lifts and
descends to compatible tori.

At a flat reference the mass-removed Schur form on
supported top data, in either proper or compressed
convention, dominates their coarse curl square with
coefficient one. Indeed the minimizing lift has exact
row value, its action curl plus the gauge term are
the mass-removed form, and the flat flux row is a
physical contraction. In the compressed case the
additional flat plaquette squares are nonnegative.
Combining this observation with \eqref{mb:eq-flat-tree-quantitative}
gives the strict tree gap \(c L^{-5}\) on top-supported
data in the kernel of their outer straight constraint.


% Construction source: assembly/curvature-slices.tex
\subsection{The deterministic curvature slices}

Let \(\mathcal R_j\) be the level-\(j\) sheet average on
exact fine two-cochains, including \(\mathcal R_0=I\).
In counting measure let \(W_j\) be the orthogonal complement
of \(\ker\mathcal R_j\) inside the exact cochains and \(P_j\)
its orthogonal projection. The sheet composition gives
\(W_{j+1}\subset W_j\). Thus \(P_jP_k=P_kP_j=P_k\) for
\(j\le k\), and every exact cochain \(F\) satisfies
\begin{equation}\label{mb:eq-flat-slice-sum}
 F=P_KF+\sum_{j=1}^K (P_{j-1}-P_j)F .
\end{equation}
The terms are mutually orthogonal because they belong
to \(W_K\) and \(W_{j-1}\cap W_j^\perp\), respectively.
These are identities of deterministic finite operators.
The linear background current also has its row form.
Since \(P_jF\) is in the range of the transpose of
\(\mathcal R_j\) restricted to exact cochains, it has
the form \(P_{\rm ex}\mathcal R_j^{\mathsf T}\phi\).
The identity \(\mathrm d_j S_{L^j}=\mathcal R_j\mathrm d\)
then gives
\[
 \mathrm d^{\mathsf T}P_jF
       =S_{L^j}^{\mathsf T}\mathrm d_j^{\mathsf T}\phi .
\]
Here \(P_{\rm ex}\) is the projection onto exact
cochains, and \(\mathrm d^{\mathsf T}P_{\rm ex}
=\mathrm d^{\mathsf T}\). Thus the minimizing
background has a row multiplier even away from
the zero datum.

\begin{proposition}[Kernel of a flat curvature slice]
\label{mb:prop-flat-slice-kernel}
For \(1\le j\le K\) at the flat periodic reference, the counting-measure
kernel of \(P_{j-1}-P_j\) obeys
\[
  |(P_{j-1}-P_j)(p,q)|
       \le C_L n^{-4}e^{-c_L d_{j-1}(p,q)},\qquad n=L^{j-1},
\]
where the fine mesh is \(1/n\) in level-\((j-1)\) units
and \(d_{j-1}\) is periodic distance in those units.
The bound is uniform in \(j\) and compatible volume.
\end{proposition}

\begin{proof}
Let \(H_n a\) be the unique least-norm fine curvature with
sheet-average equal to the curl of the intermediate datum
\(a\). Then \(H_n^{\mathsf T}H_n=\mathcal E_n\) in counting
coordinates, by Proposition~\ref{mb:prop-flat-effective}.
We first show that \(H_n\) has a kernel bounded by
\(Cn^{-2}e^{-c d_{j-1}}\).

Use the notation of \eqref{mb:eq-flat-symbol}. Define
\(B_m=\prod_\nu b_{\nu,m}\) and the diagonal matrix
\[
 (\mathcal D_m)_{\mu\mu}
     =e^{-i(n-1)\zeta_{m,\mu}/2}B_m b_{\mu,m}.
\]
Write
\[
 t_{m,\mu}^{\rm f}=e^{i\zeta_{m,\mu}}-1.
\]
The Feynman minimizer of the gauge-adjusted datum has
alias \(n\,\mathcal D_m\mathcal E_n a/s_m\).
Indeed its fine covariance multiplier is \(n^2/s_m\),
the transpose filter is \(\mathcal D_m/n\), and the
minimizing coarse force is \(\mathcal E_n a\) by
\eqref{mb:eq-flat-symbol}. Correcting the datum by a
fine gradient leaves its curl unchanged. At a fine
offset \(r\) in the intermediate block the curvature
symbol is therefore
\begin{equation}\label{mb:eq-curvature-lift-symbol}
 \frac1n\sum_m e^{i\zeta_m\cdot r}
       \frac{t_m^{\rm f}\wedge(\mathcal D_m\mathcal E_n a)}{s_m}.
\end{equation}
The additional factor \(n^{-2}\) converting fine
alias coefficients to the point value has already
been included.

For the noncentral aliases, the uniform strip bounds
from Proposition~\ref{mb:prop-flat-effective} give
\(|t_m^{\rm f}|\le C(1+|m|)/n\) and
\(|s_m|\ge c(1+|m|^2)\).
The phase factors are bounded on that strip because
\(|r_\nu|\le n/2\). Thus the noncentral sum is bounded
by \(Cn^{-2}\|a\|\) times
\[
 \sum_{m\in\mathbb Z^4}(1+|m|)^{-1}
                              \prod_\nu(1+|m_\nu|)^{-1}.
\]
This converges by dyadic shells: on the shell of
radius \(2^l\) the first factor costs \(2^{-l}\), and
the four truncated harmonic sums give at most
\(C(1+l)^4\).

For the central alias use
\(\mathcal E_n=sH-J^{-1}Ht^+(t^-)^{\mathsf T}H\)
from the origin cancellation proof, where \(s=s_0\).
The term \(\mathcal D_0 Ht^+\) has leading part
\(\mathcal D_0\operatorname{diag}(f_{\mu,0}^{-1})t^+
  =nB_0^{-1}t_0^{\rm f}\).
This follows by cancelling \(\sin(p_\mu/2)\) against
\(b_{\mu,0}\), including its phase.
Its wedge with \(t_0^{\rm f}\) is zero. The remaining
part is divisible by \(s\), since
\(H-\operatorname{diag}(f_{\mu,0}^{-1})\) is divisible
by \(s\). Hence the quotient in
\eqref{mb:eq-curvature-lift-symbol} is holomorphic
and \(O(n^{-1})\) there. Away from zero, \(s\)
is bounded away from zero on a smaller strip and
the same bound follows directly. The sum is periodic
by alias reindexing. The shifted contour and the
periodic image sum prove the bound for \(H_n\).

The slice subspace \(W_{j-1}\cap W_j^\perp\) is the
range of \(H_n\) on intermediate link data \(a\) with
\(S_La=0\) and zero link values on the prescribed
internal coordinate trees. Indeed any element of this
slice is \(H_n a\) and has
\(\mathrm d_L S_L a=0\) by sheet composition. A
gradient on the intermediate grid and a constant link
field can remove this closed straight datum without
changing \(\mathrm da\). Next subtract the site
potential obtained by integrating along each internal
tree. Choose its additive block constants so its
block average is zero. Its gradient then leaves
\(S_La=0\) and makes all tree links vanish.
The reverse inclusion follows from the straight
flux identity.

Let \(C_{\rm tr}\) parametrize this hard-tree and
straight-constraint subspace by its remaining
nonpivot coordinates. The flat crossing pivot
coefficient is \(L^{-3}\) and it is the only pivot
in its row, so \(C_{\rm tr}\) is bounded and local
with constants depending only on \(L\).
Its free restriction is the identity. By
\eqref{mb:eq-flat-tree-quantitative} and
\eqref{mb:eq-flat-effective}, the matrix
\(A_{\rm tr}=C_{\rm tr}^{\mathsf T}\mathcal E_n C_{\rm tr}\)
has a uniform positive gap. Its inverse has an
exponential kernel by the same weighted conjugation
argument as in Corollary~\ref{mb:cor-contour-gap}.
The orthogonal projection onto the slice is exactly
\[
 H_n C_{\rm tr} A_{\rm tr}^{-1}
                         C_{\rm tr}^{\mathsf T}H_n^{\mathsf T}.
\]
Both exterior factors give \(n^{-2}\), and their
product is summed only over the fixed-width
intermediate block cells. Multiplying the summable
exponential kernels at a smaller rate therefore
gives \(C_L n^{-4}e^{-c_Ld_{j-1}}\).
\end{proof}


% Construction source: assembly/moving-saddle.tex
\subsection{The full complex saddle on fixed data spaces}

Fix a real reference \(U_0\) of Theorem~\ref{mb:thm-layered-gap}
and use its energy norm \(X\). Let \(Z_0=\operatorname{ran}R_0\)
with its physical site \(L^2\) pairing. At a complex perturbation
\(U_A=e^{i\eta A}U_0\) in the admitted weighted \(C^1\) tube,
define
\[
 \mathcal C_A v=(Q_Av,\Gamma_Av),\qquad
 \Gamma_A=R_0 R_A D_A^{\mathsf T}:X\longrightarrow Z_0 .
\]
Only the target is fixed by \(R_0\). The site and field
operators are the actual continued ones.

\begin{proposition}[The moving-slice saddle]\label{mb:prop-moving-saddle}
For sufficiently small fixed tube radius and
\(\gamma_\omega=\sup_c\xi_c^3|\omega_c|\), let
\(H_{A,\omega}=H_W(U_A)-C_\omega(U_A)\), with \(C_\omega\)
the operator of the form \eqref{mb:eq-constraint-form}.
The independent multiplier \(\omega\) may be complex.
For \(f\in X^*\), \(b\in Y\) and \(z\in Z_0\), the system
\[
 \begin{aligned}
 H_{A,\omega}v-Q_A^{\mathsf T}\mu-D_A R_A R_0\nu &= f,\\
 Q_Av &= b,\qquad \Gamma_Av=z
 \end{aligned}
\]
has a unique solution in \(X\times Y^*\times Z_0\),
with norm at most \(C(\|f\|_{X^*}+\|b\|_Y+\|z\|_2)\).
The inverse on this fixed product is holomorphic.
At \(z=0\), the field and physical row multiplier are
exactly the solution of the constrained saddle on the
actual moving slice \(\ker(R_A D_A^{\mathsf T})\).
\end{proposition}

\begin{proof}
At the real centre the combined constraint has exact lift
\[
        (b,z)\longmapsto E_{L,0}b+D_0G_{s,0}z .
\]
The first term has zero projected divergence. The second
has zero \(Q_0\)-value by \eqref{mb:extended-gradient-identity}
and \(q_0G_{s,0}R_0=0\); its projected divergence is \(z\).
The restoration and field lift bounds give the uniform
energy norm. On its simultaneous kernel the real Wilson
form is coercive by Theorem~\ref{mb:thm-layered-gap},
because both positive penalties vanish. The full reference
saddle is therefore invertible: lift the two data, solve
on the coercive kernel, and evaluate the residual
functional on the same lift to recover the two multipliers.

The row, gradient and normalized projector estimates
give \(\|\mathcal C_A-\mathcal C_0\|\le C\|A\|_{\rm tube}\)
in the fixed spaces. Bilinear symmetry of \(R_A,R_0\)
gives the displayed transpose of \(\Gamma_A\).
The Wilson comparison and the second-leg moment give
\[
 \|H_{A,\omega}-H_W(U_0)\|_{X\to X^*}
       \le C(\|A\|_{\rm tube}+\gamma_\omega).
\]
For the Wilson term the prefix variation costs
\(O(\eta\|A\|_{\rm tube}/\mathsf L)\) and the curvature
variation costs \(O(\eta^2\|A\|_{\rm tube}/\mathsf L^2)\).
The first uses the undifferentiated potential, the second
its weighted first difference. Their insertion into the
actual plaquette form gives the stated energy comparison.
A relative Neumann series about the reference saddle
gives the estimate and holomorphy.

For \(w\in\operatorname{ran}R_A\), \(R_0w=0\) implies
\(w=(R_A-R_0)w=0\) by the small norm difference.
Hence \(\ker\Gamma_A=\ker(R_A D_A^{\mathsf T})\).
Restriction of the first equation to this kernel
eliminates \(\nu\). Conversely \(\Gamma_A\) is onto,
since \(\mathcal C_A\) has a right inverse obtained by
perturbing the displayed lift. Thus a residual functional
vanishing on its kernel belongs to its transpose range,
with a unique gauge multiplier. If \(f\) was supplied only
on the slice, any norm-preserving extension to \(X\) gives
the same field and physical multiplier. The fixed-space
holomorphy claim uses fixed source extensions.
\end{proof}


% Construction source: audit/reconstruction.tex
\section{Projected Green operators and coarse response}
\label{mb:sec-reconstruction}

All inverses in this section act on the prescribed finite fine-field or site space. We use the physical bilinear transpose \(T\), and use Hermitian norms on their complexifications. Write \(\Xi_s\) and \(\Xi_f\) for multiplication by the assigned lengths on the site and field rows, respectively. At an outward field row this is the row length, even when the other endpoint has a smaller site length. Let \(a_s\) and \(a_f\) be the positive diagonal row masses, comparable to \(\Xi_s^{-2}\) and \(\Xi_f^{-2}\). Put
\[
 X_y(u)=\|\ell^{-1}u\|_{2,B_y}+\|\nabla_U u\|_{2,B_y}.
\]
The row space remains $Y$ with its physical dual $Y^*$ from \eqref{mb:eq-row-energy}. Here \(B_y\) is the fine cell of an assigned top, with a fixed derivative collar when a difference is measured, and \(\ell\) is the comparison length of \eqref{mb:eq-slow-length}. The norms \(X_y\) sum in squares, up to bounded overlap, to the field energy norm \(X\). At the bare physical boundary use the atomic difference convention for \(X\). We work with the separated domains, complete rows and admitted real backgrounds of \cref{mb:thm-layered-gap}. Their complex tubes use the same spaces and fixed reference norms.

For a site function \(u\) and a site source \(h\), define
\[
 \|u\|_{\mathcal Y_s}=\|u\|_2+\|\ell Du\|_2,\qquad
 \|u\|_{\mathcal Z_s}=\|\ell^{-2}u\|_2+\|\ell^{-1}Du\|_2+
               \|\nabla_U Du\|_2,\qquad
 \|h\|_{\mathcal Y_s^\flat}
     =\sup_{v\ne0}\frac{|\langle h,\ell^2 v\rangle|}
                            {\|v\|_{\mathcal Y_s}} .
\]
The local versions use the complete cell collars and bare atomic stencils.

\begin{lemma}[Balanced site bounds and positive elimination]
\label{mb:lem-site-woodbury}
Keep \(A_s,G_s,F,M,P,R,\Lambda,S\) of \cref{mb:prop-site-gram}. In the distance weights \(w=e^{s\psi}\) of that proposition, \(G_s\) satisfies
\(\|wG_sw^{-1}h\|_{\mathcal Y_s}\le C\|h\|_{\mathcal Y_s^\flat}\)
for arbitrary site sources. For the inner Wilson Hessian \(H_W\) put
\[
 K_F=H_W+DD^T+Q^Ta_fQ,\quad G_F=K_F^{-1},\qquad
 K_0=K_F-DPD^T .
\]
Then \(\Theta=M-(DF)^TG_FDF\) has a uniform positive upper and lower bound and
\begin{equation}\label{mb:eq-zero-current-woodbury}
 G_0=K_0^{-1}=G_F+G_FDF\Theta^{-1}(DF)^TG_F .
\end{equation}
The identities and inverse bounds extend holomorphically to a common smaller complex tube.
\end{lemma}

\begin{proof}
For an arbitrary source \(h\), test \(A_su=h\) against \(\ell^2u\). The left side is the balanced test of \eqref{mb:eq-gram-test}: its leading terms are comparable to \(\|\ell Du\|_2^2+\|qu\|_{\rm top}^2\). The differences of the weight cost \(C_L\operatorname{Lip}(\ell)\|u\|_{\mathcal Y_s}^2\) by the product and complete-row estimates there. The pillar Poincar\'e bound absorbs them. On the right the pairing is bounded by \(\|h\|_{\mathcal Y_s^\flat}\|u\|_{\mathcal Y_s}\) by definition, which proves the balanced inverse bound. Conjugation by \(w\) adds a relative error \(C_L|s|\) from first differences and the \(O_L(|s|)\) complete-row oscillations. This proves the weighted bound on arbitrary sources.

The layered gap gives \(K_0\ge\delta K_F>0\) with a fixed \(\delta>0\). In real orthonormal physical coordinates use \(C=G_F^{1/2}DFM^{-1/2}\). Conjugating the preceding inequality gives \(I-CC^T\ge\delta I\). The singular values of \(C\) then give \(I-C^TC\ge\delta I\) and consequently \(\delta M\le\Theta\le M\). Multiplying the right side of \eqref{mb:eq-zero-current-woodbury} by \(K_F-DFM^{-1}(DF)^T\) cancels the two terms containing \(\Theta^{-1}\), by the definition of \(\Theta\). This proves the formula.

For the tube, the local differences and finite row moments give relative form perturbations of size \(C\varepsilon\) for \(A_s\) and \(K_F\) and a perturbation of the same size in the balanced site bound. In each covariant first difference the change is multiplication by \(O(\varepsilon/\ell)\); for a row the weighted sums of the changed coefficients are \(O(\varepsilon)\) in normalized row units. The Wilson form perturbation is controlled by its principal difference and curvature multiplication terms. These bounds and the real gaps give the site and field inverses by the Neumann formula. They give \(F\) in the balanced norm and \(\Theta\) in the normalized row norm, since \(DF\), from the site-top space with the physical \(L^2\) pairing in \eqref{mb:physical-pairings} to \(X^*\), is bounded by \(\|F\|_2\). Thus \(M^{-1}\) and \(\Theta^{-1}\) continue by the same argument. The displayed identities are algebraic with bilinear transpose.
\end{proof}

We record the weighted restoration estimate used with the current. In the same distance weights as above,
\begin{equation}\label{mb:eq-weighted-restoration}
 \|\ell^{-2}w\Lambda w^{-1}v\|_2+
 \|wD\Lambda w^{-1}v\|_X+\|wSw^{-1}v\|_X\le C\|v\|_X.
\end{equation}
For this purpose, the balanced estimate shows that \(w Fw_{\rm top}^{-1}\) and its conjugate with \(-s\) differ from \(F\) by \(O(|s|)\) in the top-to-fine \(L^2\) norm. Expand \(wM w^{-1}\) using these two factors, in the top weight. Its inverse is then a uniform small perturbation of \(M^{-1}\). Thus \(wRw^{-1}\) is bounded on fine \(L^2\). Set \(\lambda=\Lambda w^{-1}v\). It has \(q\lambda=0\) and \(\|w\Delta\lambda\|_2\le C\|v\|_X\), by the first-difference product rule for \(D^T\). Apply the zero-average pillar estimate with \(h=\ell^{-1}w\). The row error of \(q(h\lambda)\), since \(q\lambda=0\), is bounded by \(C(\operatorname{Lip}\ell+|s|)\|\ell^{-2}w\lambda\|_2\) in the grade-two norm. Testing \(\Delta\lambda\) against \(h^2\lambda\) and absorbing the corresponding first-difference error bounds \(\|\ell^{-2}w\lambda\|_2+\|\ell^{-1}wD\lambda\|_2\). The covariant curl of \(D\lambda\) is a curvature times \(\lambda\), while its divergence is \(\Delta\lambda\). The component-gradient estimate applied to \(wD\lambda\) therefore bounds the remaining part of \(X\): both weight commutators in curl and divergence are \(O(|s|\ell^{-1}wD\lambda)\). The zero-average bare collar keeps the nonzero stencils inside the action region. This proves \eqref{mb:eq-weighted-restoration} with only first differences of the weight. The same estimate holds on the smaller complex tube by perturbing the weighted operator and using the same zero-row equation and fixed real norms.

\begin{proposition}[Current completion]
\label{mb:prop-current-completion}
For a fine current \(J\), write
\[
 (W[J]\lambda)(b)=\tfrac i2[\lambda(b_-)+
                   \operatorname{Ad}(U_b)\lambda(b_+),\,J(b)] .
\]
The contraction of the second jet of the finite relative gauge by \(J\) is
\begin{equation}\label{mb:eq-current-completion}
 B[J]=W[J]\Lambda+\Lambda^T W[J]^T-
          \tfrac12\Lambda^T(D^T W[J]+W[J]^TD)\Lambda .
\end{equation}
For the actual Wilson current \(j_U\),
\[
                 H_W-B[j_U]=S^TH_WS .
\]
The bound \(\|wB[J]w^{-1}\|_{X\to X^*}\le C\gamma_{\rm t}\) holds with \(\gamma_{\rm t}=\sup_{b\text{ touching }\Omega_1}\ell(b)^3|J(b)|\). The empty touching stock gives \(\gamma_{\rm t}=0\).
\end{proposition}

\begin{proof}
For \(u\) set \(\ell_u=(\Lambda u)(b_-)\) and \(r_u=\operatorname{Ad}(U_b)(\Lambda u)(b_+)\). Apply the finite gauge \(e^{-i\Lambda u}\) to \(e^{i\eta u}U\). Its Hermitian bond logarithm relative to \(U\), divided by \(\eta\), is
\((i\eta)^{-1}\log(e^{-i\ell_u}e^{i\eta u}e^{i r_u})\).
Its first derivative is \(Su\), and its second differential on \(u,v\) is
\[
 \tfrac i2([u,\ell_v+r_v]+[v,\ell_u+r_u])
     -\tfrac i{2\eta}([\ell_u,r_v]+[\ell_v,r_u]).
\]
Indeed the degree-two part of the logarithm of the product is half the sum of commutators in the written order. Trace invariance identifies the first term with the first two terms of \eqref{mb:eq-current-completion}. For the second, use \(D\Lambda u=(\ell_u-r_u)/\eta\) and
\[
 [\ell_u-r_u,\ell_v+r_v]+[\ell_v-r_v,\ell_u+r_u]
                  =2([\ell_u,r_v]+[\ell_v,r_u]).
\]
This gives the last term with its factor \(1/2\). The chain rule for the gauge-invariant inner Wilson action gives the identity at \(j_U\).

In a bilinear pairing move each displayed \(D^T\) and \(\Lambda^T\) onto its test vector. Each of the resulting terms pairs \(W[J]\lambda\) against a field with \(X\)-norm controlled by \eqref{mb:eq-weighted-restoration}. Pointwise endpoint comparison and the current grade give
\(\|\ell W[J]\lambda\|_2\le C\gamma_{\rm t}\|\ell^{-2}\lambda\|_2\)
on the touching stock. Indeed both site factors \(\Lambda u,\Lambda v\) vanish at bare sites by the bare singleton rows of \(q\). A bond with neither endpoint in \(\Omega_1\) therefore contributes zero. In the weighted pairing the endpoint weight ratio is bounded by a fixed constant; apply \eqref{mb:eq-weighted-restoration} with both signs. No difference is applied to \(J\).
\end{proof}

Let \(C_\omega\) denote the Hessian of the actual per-length straightened constraint, contracted with the physical row multiplier \(\omega\). If its quadratic Taylor term is \(C^{(2)}\), the convention is
\(\langle u,C_\omega u\rangle=2\langle\omega,C^{(2)}(u)\rangle_{\rm rows}\).
In the raw coordinates the same term is \(\Xi_f C^{(2)}\), and the transpose of a raw lift in counting row coordinates is \(\Xi_f^3 H^T\). Their counting pairing therefore equals the physical pairing just written. The two-leg row bounds give
\begin{equation}\label{mb:eq-two-leg-energy}
 \|\ell C_\omega u\|_2
       \le C\Big(\sup_{c\ \mathrm{nonbare}}\xi_c^3|\omega(c)|\Big)
                         \|\ell^{-1}u\|_2 .
\end{equation}
In fact the second derivative in one leg has a normalized fixed-leg sum on the full carrier. In the transpose contraction the fine factor \(\eta^{-4}\) cancels that leg's density \((\eta/\xi_c)^4\) against \(\xi_c^4\). The remaining per-length factor has size \(\xi_c\) when the two arguments are unscaled field values, so a multiplier of grade three leaves a field mass of grade two. Bounded row incidence and comparable assigned lengths give both absolute sums of the resulting \(\ell C_\omega\ell\) matrix. This proves \eqref{mb:eq-two-leg-energy} by the Schur inequality, and also proves its weighted version because the row diameter in the layer metric is bounded. The actual bare mean is linear and has zero second jet.

\begin{theorem}[Normalized Schur matrices and coupled response]
\label{mb:thm-coarse-schur}
On the admitted real and complex coefficient domains, take the touching current sufficiently small and define, in this order,
\begin{align*}
 K_J&=K_0-B[J],&
 G_J&=K_J^{-1},&
 T_J&=a_f^{1/2}QG_JQ^Ta_f^{1/2},\\
 H_J&=G_JQ^Ta_f^{1/2}T_J^{-1}a_f^{1/2},&
 \omega&=H_J^TJ,&
 K_{1,J}&=K_J-S^T C_\omega S .
\end{align*}
Write $L^2_r$ for the physical row Hilbert space of \eqref{mb:physical-pairings}. The second variant has \(G_{1,J},T_{1,J},H_{1,J}\) defined from \(K_{1,J}\) by the same inverse, row matrix and lift formulas. For either variant write \(K,G,T_f,H\) for these four objects and put \(V_f=QGQ^T\). All weighted norms
\[
 G:X^*\to X,\quad T_f^{\pm1}:L^2_r\to L^2_r,
 \quad H:Y\to X
\]
are uniform in mesh, depth and volume. \(QH=I\), and \(H\) is the \(K\)-energy-minimizing lift at real data. In counting orthonormal row coordinates \(\Xi_f^2\), the entries of \(T_f^{-1}\) are \(O(e^{-m d(c,d)})\). The entries of \(V_f^{-1}\) as a counting array are \(O(\xi_c^{-3}\xi_d e^{-m d(c,d)})\), and as a physical row kernel are \(O(\xi_c^{-3}\xi_d^{-3}e^{-m d(c,d)})\).

For sources \(f\in X^*\), \(b\in Y\), the solution of \(Ku-Q^T\mu=f,\ Qu=b\) is
\begin{equation}\label{mb:eq-field-saddle}
 u=(I-HQ)Gf+Hb,\qquad \mu=V_f^{-1}(b-QGf).
\end{equation}
The \(X\) and \(Y^*\) output norms are bounded by \(C(\|f\|_{X^*}+\|b\|_Y)\). Removing \(Q^Ta_fQ\) changes the multiplier to \(\mu-a_fb\). The corresponding datum-to-multiplier map is \(V_f^{-1}-a_f\). These formulas are holomorphic and have exponential cell-to-cell bounds in the stated norms.
\end{theorem}

\begin{proof}
First consider \(K_0\). Conjugating its local parts by \(w=e^{s\psi}\) changes the weak form by \(O(|s|)\) in \(X\to X^*\), by the difference product rule and complete-row oscillation bound. For \(DPD^T\) use the weighted \(L^2\) estimate of \(P\) in the proof of \eqref{mb:eq-weighted-restoration}. It is \(O(|s|)\) close to \(P\), and the two adjacent first differences are bounded \(L^2\to X^*\) and \(X\to L^2\), with commutators of size \(O(|s|)\). Inverting the conjugated \(K_0\) is thus a uniform perturbation of its real gap. The first current subtraction is small by \cref{mb:prop-current-completion}.

For the normalized row matrix at zero current the upper bound follows from the energy trace bound of \(a_f^{1/2}Q\). For a physical row vector \(\beta\) use the exact field lift \(E_Q a_f^{-1/2}\beta\) of \cref{mb:lem-field-lift}, whose \(X\)-norm is at most \(C\|\beta\|_{\rm rows}\). Cauchy--Schwarz in the real \(K_0\) energy bounds
\[
 \|\beta\|_{\rm rows}^2
  =\langle Q^Ta_f^{1/2}\beta,E_Qa_f^{-1/2}\beta\rangle
  \le C\langle\beta,T_0\beta\rangle^{1/2}\|\beta\|_{\rm rows}.
\]
This proves the positive gap. Conjugating \(a_f^{1/2}Q\) with fine and row weights introduces an \(O(|s|)\) normalized difference. The normalized \(T_J\), conjugated by the row weight, is \(O(|s|+\gamma_{\rm t}+\varepsilon)\) from the real reference \(T_0\), where \(\varepsilon\) is the tube radius. Thus its inverse is uniform. The lift formula and \(QH_J=I\) now give a uniformly weighted \(H_J\). The first lift readout estimate proved immediately below gives a nonbare \(\omega\) of size \(C\gamma_{\rm t}\). Use \eqref{mb:eq-two-leg-energy} and \eqref{mb:eq-weighted-restoration} for the second subtraction, then apply the identical trace and perturbation argument to obtain \(T_{1,J}^{-1}\) and \(H_{1,J}\).

Choose \(\psi\) to be the distance to the input cell, truncated above at the output distance. The weighted bounds give the exponential cell estimates. Conversion from the normalized counting array to \(V_f^{-1}\) multiplies the \(c,d\) entry by \(\xi_c^{-2}a_f(c)^{1/2}a_f(d)^{1/2}\xi_d^2\), and conversion to a physical kernel divides by \(\xi_d^4\). This proves both exponents of the physical lengths. Substitution proves \eqref{mb:eq-field-saddle}. At real data \(Ku\) for \(u=Hb\) is in \(\operatorname{ran}Q^T\) and hence \(K\)-orthogonal to \(\ker Q\); strict positivity proves the minimizing assertion. The multiplier shift follows by substituting \(Qu=b\) into the removed term. Every analytic inverse above is a Neumann perturbation in fixed reference norms.
\end{proof}

\begin{corollary}[The two physical current readouts]
\label{mb:cor-current-readouts}
For \(i=0,1\), put \(H_{0,J}=H_J\) and \(\omega_i=H_{i,J}^T J\). For either variant
\[
 \sup_c\xi_c^3|\omega_i(c)|\le C\gamma,\qquad
 \sup_{c\ \mathrm{nonbare}}\xi_c^3|\omega_i(c)|\le C\gamma_{\rm t},
 \qquad \gamma=\sup_b\ell(b)^3|J(b)|.
\]
The first readout alone enters \(K_{1,J}\). Both variants depend on \(J\) through its restriction to the touching stock. The unrestricted readout has the exact additional bare component obtained by identifying each nontouching current bond with its bare row.
\end{corollary}

\begin{proof}
The proof for the first variant uses only its weighted lift already obtained before constructing \(K_{1,J}\). For a unit colour \(e\), let \(b^{(c)}(c)=\xi_c^{-1}e\), zero elsewhere. Its \(Y\)-norm is bounded above and below by fixed constants. The identity
\[
 \xi_c^3\langle e,\omega_i(c)\rangle
      =\langle H_{i,J} b^{(c)},J\rangle
\]
uses the exact physical row volume. On \(B_y\), \(\|\ell J\|_2\le C\gamma\) by the cell volume bound. The weighted lift estimate bounds \(\|\ell^{-1} H_{i,J}b^{(c)}\|_{2,B_y}\) by \(Ce^{-m d(y,c)}\). Summing the pairings uses the metric growth with rate strictly smaller than \(m\). Hence the constant is independent of the number of cells. A datum supported on nonbare rows lifts to a field vanishing on the bare rows since \(QH_{i,J}=I\) and the bare row is the fine coordinate. A current outside the touching stock therefore contributes only to its identical bare row, with both measures \(\eta^4\). It drops out of \(C_\omega\). Apply the estimate to the restricted current for the nonbare bound, first for \(H_J\) and then for the second variant.
\end{proof}

For clarity about the independent-current extensions, put \(\delta J=J-j_U\). The preceding identities give, for either variant,
\begin{equation}\label{mb:eq-off-current-defect}
 K DG_sR=DR-B[\delta J]DG_sR,\qquad
 GDR=DG_sR+GB[\delta J]DG_sR .
\end{equation}
Indeed \(S DG_sR=0\), \(QDG_sR=0\), and \(RD^TDG_sR=R\). Transposing and composing with the row factors of the lift gives its slice defect. In particular it vanishes when \(J=j_U\) on the touching stock. If a physical pair additionally satisfies \(j_U=Q^T\omega_{\rm phys}\), the first lift gives \(H_J^Tj_U=\omega_{\rm phys}\). We next verify the Lagrangian Ward identity for this pair.

Write \(\mathfrak q(a)\) for the actual per-length nonlinear constraint in the reference chart and \(\Gamma_a\lambda\) for the gauge tangent in that chart. Differentiating the exact average and node-factor covariance gives
\[
 D\mathfrak q(a)\Gamma_a\lambda
     =\overline D(a)\tau_a\lambda,\qquad
 D\mathfrak q(0)=Q,\quad \Gamma_0=D,\quad
 \tau_0=q,\quad \overline D(0)=\overline D_{\rm ext}.
\]
Here the two moving coarse endpoint factors are included in \(\overline D(a)\). At an outward row the further one-step mean of the finer endpoint factors enters that operator. Thus its value at zero is the extended row difference on own-top values. The Wilson Ward identity \(\langle j_U,D\lambda\rangle=0\), the equality \(j_U=Q^T\omega_{\rm phys}\), and onto \(q\) give \(\overline D_{\rm ext}^T\omega_{\rm phys}=0\).

For \(q\lambda=0\), differentiate the displayed covariance at zero in direction \(v\) and pair with \(\omega_{\rm phys}\). The term differentiating \(\overline D(a)\) vanishes because \(\tau_0\lambda=0\); the term differentiating \(\tau_a\) vanishes by the coarse transpose identity. What remains is
\[
 \langle v,C_{\omega_{\rm phys}}D\lambda\rangle
   +\langle j_U,\Gamma'_0[v]\lambda\rangle=0.
\]
The differentiated action Ward identity gives the same equation with \(H_W\) in place of \(C_{\omega_{\rm phys}}\). Subtract to obtain the full identity \((H_W-C_{\omega_{\rm phys}})D\lambda=0\) on \(\ker q\), with arbitrary test \(v\). Since the range of \(I-S\) lies in \(D\ker q\) and the form is symmetric, \(S^T(H_W-C_{\omega_{\rm phys}})S=H_W-C_{\omega_{\rm phys}}\). On that pair this gives
\[
 K_{1,J}=H_W-C_{\omega_{\rm phys}}+DRD^T+Q^Ta_fQ .
\]
This identification uses the stated critical-pair identity as input. The construction of the two Green operators and their independent-current bounds needs only the admitted coefficient data.

The constrained Green operator is
\[
 \mathfrak G=G_{1,J}-G_{1,J}DRD^TG_{1,J}-H_{1,J}QG_{1,J}.
\]
On shell \eqref{mb:eq-off-current-defect} gives \(QG_{1,J}DR=0\) and \(RD^TG_{1,J}DR=R\). Thus \(\mathfrak G f\) belongs to \(\ker Q\cap\ker RD^T\). On this common slice, pairing \(K_{1,J}\mathfrak G f\) with a test field gives exactly the pairing with \(f\), since both subtractions are in the respective transpose ranges. This is the inverse on the stated slice. The same formula defines the independent-current extension. The first lift on shell is the stationary field for the Wilson quadratic form on \(Qu=b,\ RD^Tu=0\): its penalized equation has multiplier \(V_f^{-1}b\), the gauge term vanishes, and \(S\) acts as the identity on both tests and solutions. The stationary value is \(\tfrac12\langle b,(V_f^{-1}-a_f)b\rangle_{\rm rows}\).

\subsection{A local augmented system with independent sources}

The site spaces \(\mathcal Y_s,\mathcal Z_s,\mathcal Y_s^\flat\) are those in the proof of \cref{mb:lem-site-woodbury}. A site divergence \(h_0+D^Th_1\) has \(\mathcal Y_s^\flat\)-norm at most \(C(\|\ell^2h_0\|_2+\|\ell h_1\|_2)\). This follows by applying \(D\) to \(\ell^2v\) in the test pairing. In particular the norms below admit divergence data with no pointwise differences of their coefficients.

\begin{proposition}[Unfolded response]
\label{mb:prop-unfolded-response}
Use the same real or complex admission as \cref{mb:thm-coarse-schur}, and set \(W=W[J^{\rm t}]\), where \(J^{\rm t}\) is the restriction to the touching stock. In the second variant use the first lift's nonbare \(\omega\) from that theorem. For \(\epsilon=0,1\) put
\[
 e_\omega=\epsilon C_\omega(A-D\lambda),\qquad
 h_J=W^TA-\tfrac12(D^TW\lambda+W^TD\lambda)-D^T e_\omega.
\]
The finite system with arbitrary sources
\begin{equation}\label{mb:eq-unfolded-response}
 \begin{aligned}
 K_FA-Dv-W\lambda-e_\omega-Dt&=f_0,\\
 A_s\lambda+v-D^TA&=g,\qquad q\lambda=c,\\
 A_sv-q^T\Xi_s^{-2}\beta&=h,\\
 A_st+q^T\Xi_s^{-2}\beta'-h_J&=h',\\
 A_s\zeta-t&=r,\qquad q\zeta=d_0
 \end{aligned}
\end{equation}
has a unique solution, with bound
\[
 \begin{aligned}
 &\|A\|_X+\|\lambda\|_{\mathcal Z_s}+\|\zeta\|_{\mathcal Z_s}+\|v\|_{\mathcal Y_s}+\|t\|_{\mathcal Y_s}+
        \|\beta\|_{\rm top}+\|\beta'\|_{\rm top}\\
 &\qquad\le C\bigl(\|f_0\|_{X^*}+\|g\|_2+\|r\|_2+
      \|h\|_{\mathcal Y_s^\flat}+\|h'\|_{\mathcal Y_s^\flat}
      +\|\Xi_s^{-2}c\|_{\rm top}+\|\Xi_s^{-2}d_0\|_{\rm top}\bigr).
 \end{aligned}
\]
Adding \(-Q^T\mu\) in the first equation and \(QA=b\) yields the analogous bound with \(\|\mu\|_{Y^*}\) and \(\|b\|_Y\) added to the two sides. At zero auxiliary sources the field inverse is \(G_J\) or \(G_{1,J}\), respectively. These statements concern the actual system on its prescribed full finite domain.
\end{proposition}

\begin{proof}
Fix \(A\). Write \(f=D^TA+g\) and \(v_0=G_sh\). The second and third equations have the unique solution
\[
 \beta=M^{-1}\{F^T(f-v_0)-\Xi_s^{-2}c\},
 \qquad v=v_0+F\beta,\qquad \lambda=G_s(f-v).
\]
The constraint follows from \(qG_sF=\Xi_s^2 M\). The balanced site inverse and \(M\) gap bound \(v,\beta\) in their stated norms. On a complex tube use the right lift
\[
 E_s(U)=E_s(U_0)\{q_U E_s(U_0)\}^{-1}.
\]
The matrix in braces is $I+O(\varepsilon)$ in the top norm
$\|\Xi_s^{-2}\,\cdot\,\|_{2,s}$, by the actual row
perturbation and the compact real lift estimate. Its Neumann
inverse gives $q_UE_s(U)=I$ and the same $\mathcal Z_s$
bound. Subtract this site lift \(E_sc\) from \(\lambda\) and use the zero-average second-difference estimate with source
\[
       \Delta(\lambda-E_sc)=f-v-q^Ta_sc-\Delta E_sc.
\]
The actual row norm and the bound on \(E_s\) show this source is \(L^2\)-bounded by the declared norms. This gives \(\lambda\) in \(\mathcal Z_s\). For a given \(h_J\) the last two equations have the unique solution
\[
 t_0=G_s(h'+h_J),\quad
 \beta'=M^{-1}\{F^T(t_0+r)-\Xi_s^{-2}d_0\},\quad
 t=t_0-F\beta',\quad \zeta=G_s(t+r).
\]
Subtracting \(E_sd_0\) now gives the same \(\mathcal Z_s\) estimate for \(\zeta\).

At zero auxiliary sources \(v=PD^TA\), \(\lambda=\Lambda A\) and \(t=RG_s h_J\). Since \(DRG_s=\Lambda^T\), the terms \(-W\lambda-Dt-e_\omega\) in the field equation equal \(-B[J]A-S^T e_\omega\), with \(e_\omega=\epsilon C_\omega SA\). This proves exact elimination to \(K_J\) or \(K_{1,J}\). For nonzero auxiliary sources subtract these zero-auxiliary solutions at the same \(A\). The \(v,\lambda\) differences are bounded by \(\|g\|_2+\|h\|_{\mathcal Y_s^\flat}+\|\Xi_s^{-2}c\|_{\rm top}\) from the preceding formulas. The resulting change in \(h_J\) is bounded in \(\mathcal Y_s^\flat\) by those same norms. Indeed \(\|\ell W\lambda\|_2\le C\gamma_{\rm t}\|\ell^{-2}\lambda\|_2\), and testing \(W^TA\) against \(\ell^2v\) gives \(C\gamma_{\rm t}\|A\|_X\|v\|_{\mathcal Y_s}\). Replace \(A\) by \(D\lambda\) for the third term in \(h_J\), and use the divergence estimate and \eqref{mb:eq-two-leg-energy} for the other two terms. Thus the errors \(Dv,Dt,W\lambda,e_\omega\) have \(X^*\)-norm at most the auxiliary source bound. The actual field inverse, or \eqref{mb:eq-field-saddle} with the added constraints, bounds \(A\); the site equations then give all other estimates. This proof uses the same site algebra and relative perturbation estimates on the smaller complex tube. In particular the finite-carrier equations do not require invertibility of their compressions with new artificial boundary conditions.
\end{proof}


% Construction source: draft/proper-site.tex
\section{The proper site inverse in a strong norm}\label{mb:proper-site}

The site gauge parameter requires a second derivative in a
H\"older norm. This conclusion uses a regular frame on each
complete interior box. An actual-current bound in an admitted
Coulomb frame is one sufficient way to verify this input.

Let $U_0$ be a real proper reference on the original determining
graph. On the nested complete interior boxes it obeys the small
curvature bound and admits real frames satisfying
\eqref{mb:eq-local-covariant-strong-coefficients}, uniformly after
rescaling by the assigned length. Alternatively, it admits the
Coulomb frames and actual-current bounds of
Corollary~\ref{mb:cor-current-holder}, with enough additional
outer boxes for that corollary. Proposition~\ref{mb:prop-coulomb-frame}
constructs these real Coulomb frames when its enlarged curvature
stocks are available. Let $q_0$ be its own-pillar site average,
and evaluate the operators $G_s,F,M,R$ of
Proposition~\ref{mb:prop-site-gram} at this reference.
Write $R_0$ for its projector $R$.
The site-gap and normalized-Gram estimates supply these
inverses, as in Proposition~\ref{mb:prop-site-gram} and
Lemma~\ref{mb:lem-site-woodbury}.
Their exponential column bounds use the physical
orthonormal top basis. The two pairings in $F^{\mathsf T}$ must
therefore be retained.

Fix $0<\alpha<1$. On each nonbare pillar, take the nested complete
boxes specified for the strong interior estimates. H\"older
increments compare values after the fixed short coordinate
transport inside the same box. They include increments across
adjacent pillar faces. Define
\begin{align}
 \|\rho\|_{\alpha,2}
   &=\sup_y\{\xi_y^2\|\rho\|_{\infty,B_y}
                    +\xi_y^{2+\alpha}[\rho]_{\alpha,y}^{U_0}\},
                \label{mb:site-source-norm}\\
 \|\lambda\|_{Y^{2,\alpha}}
   &=\sup_y\{\|\lambda\|_{\infty,B_y}
        +\xi_y\|D_0\lambda\|_{\infty,B_y}
        +\xi_y^2\|D_0^2\lambda\|_{\infty,B_y}
        +\xi_y^{2+\alpha}[D_0^2\lambda]_{\alpha,y}^{U_0}\}.
                \label{mb:site-solution-norm}
\end{align}
At a bare pillar $\xi_y=\eta$, use the fixed finite collars of
present stencils. The condition $q_0\lambda=0$ then fixes the
actual bare-site value to zero. No smoothness through a rough
exterior boundary is required.

\begin{proposition}[Strong proper site inverse]\label{mb:proper-site-inverse}
Under the preceding real curvature and admitted-frame hypotheses,
with the fixed separation large enough for the column summations,
\begin{equation}\label{mb:proper-site-estimates}
 \|R_0\rho\|_{\alpha,2}\le C_\alpha\|\rho\|_{\alpha,2},\qquad
 \|G_sR_0\rho\|_{Y^{2,\alpha}}
                         \le C_\alpha\|\rho\|_{\alpha,2}.
\end{equation}
For $\rho\in\operatorname{ran}R_0$, the field $\lambda=G_s\rho$
is the unique solution on the actual graph of
\begin{equation}\label{mb:proper-site-equations}
          q_0\lambda=0,\qquad R_0D_0^{\mathsf T}D_0\lambda=\rho.
\end{equation}
\end{proposition}
\begin{proof}
Proposition~\ref{mb:prop-site-pointwise} gives the normalized
site-column estimates
\[
 |F_z(x)|\le C\xi_y^{-2}e^{-m d(y,z)},\qquad
 [F_z]_{\alpha,y}^{U_0}
             \le C_\alpha\xi_y^{-2-\alpha}e^{-m d(y,z)}.
\]
These are estimates for the physical orthonormal column indexed
by $z$. A source of unit norm in \eqref{mb:site-source-norm}
has physical $L^2$ norm at most a constant on each pillar.
Pairing it with $F_z$ and summing the exponential cell bounds
therefore gives uniformly bounded top coefficients.
The normalized Gram inverse preserves this bound by its
two-sided exponential sums. Summing the output columns proves
the first estimate in \eqref{mb:proper-site-estimates} for
$R_0=I-FM^{-1}F^{\mathsf T}$. For the H\"older seminorm,
take the increment of each output column before summing.
The source is not differentiated.

Set $u=G_sR_0\rho$. Since
$F^{\mathsf T}=\Xi_s^{-2}q_0G_s$ and $F^{\mathsf T}R_0=0$,
we have $q_0u=0$. The defining resolvent equation then gives
\begin{equation}\label{mb:proper-site-unprojected}
                         D_0^{\mathsf T}D_0u=R_0\rho.
\end{equation}
The exponential energy blocks of $G_s$ imply
\[
 \xi_y^{-2}\|u\|_{2,B_y}
       +\xi_y^{-1}\|D_0u\|_{2,B_y^+}
                         \le C\|\rho\|_{\alpha,2}.
\]
Indeed a resolvent block has factor
$\xi_y\xi_z e^{-m d(y,z)}$, and the source has bounded
cell $L^2$ norm. The remaining scale ratio is summable by
\eqref{mb:rate-reserve}. The derivative block is treated in
the same way.

Rescale a nonbare interior box by $\xi_y$. The bounded-force
bootstrap in Proposition~\ref{mb:prop-local-scalar} gives a $C^{1,\alpha}$ bound
on a smaller box: choose a finite $p>4/(1-\alpha)$ and use
\eqref{mb:proper-site-unprojected}. Its right-hand side has
the required $C^\alpha$ bound by the first estimate already
proved. The admitted-frame hypothesis supplies the uniformly
bounded $C^{1,\alpha}$ reference potential on the larger box;
in the Coulomb-frame alternative this is
Corollary~\ref{mb:cor-current-holder}.
Corollary~\ref{mb:cor-local-covariant-schauder}
on one further nested box gives the second derivatives and
their H\"older seminorm. Rescaling gives exactly the powers
of $\xi_y$ in \eqref{mb:site-solution-norm}.

At a bare cell, the physical atomic $L^2$ estimate bounds the
finitely many values in its fixed collar. A difference of order
$r$ costs at most $C\eta^{-r}$, and a H\"older quotient costs
at most $C\eta^{-\alpha}$. The weights in
\eqref{mb:site-solution-norm} cancel those factors. This proves
the same estimate on present bare stencils, with $q_0u=0$
retaining the actual zero values.

For uniqueness, suppose $q_0\lambda=0$ and
$R_0D_0^{\mathsf T}D_0\lambda=0$. Since $R_0$ is the
orthogonal projection onto the complement of $\operatorname{ran}F$,
write $D_0^{\mathsf T}D_0\lambda=Fa$. The vanishing site
average gives $\lambda=G_sFa$. Applying $q_0$ yields
\[
                0=q_0G_sFa=\Xi_s^2Ma.
\]
The normalized Gram is invertible, so $a=0$ and $\lambda=0$.
Existence follows from \eqref{mb:proper-site-unprojected}.
\end{proof}


% Construction source: audit/local-families.tex
\section{Connected families for the finite inverses}
\label{mb:sec-local-families}

A family coefficient has a finite connected set of tagged cubes, together with the endpoint, mesh, row and domain data read on that set. A tag includes its level. A cube on the exterior is allowed as a geometry tag, with the bare scale. Besides the fine connection, a coefficient reads the complete plaquette and row stocks, the row weights and the zero-extension and free-coordinate masks on its carrier. In particular a row with some endpoints outside a patch retains its actual coefficients.

For a cell \(B_y\), use the force norm \(\mathcal F_{2,y}(f)=\xi_y\|f\|_{2,B_y}\) and the energy norm \(X_y\) of the preceding section, with the same conventions for scalar energy. For spaces involving derivatives choose a smooth square partition subordinate to fixed enlargements of the assigned cells. One such partition is obtained by taking sampled smooth bumps equal to one on each cell, supported in fixed enlargements, and dividing each by the square root of the sum of their squares. Separation bounds the number of overlaps and the ratios of scales at every overlap. The difference of order \(a\) of the bump is \(O(\xi_y^{-a})\) on its support. Thus the product rule gives bounded multiplication on the local energy spaces, their duals and the second-difference spaces. On \(L^p\) and row spaces one can use their actual cell partition. Write \(\phi_y\) for these localization multipliers. The identity used at an intermediate slot is \(\sum_y\phi_y^2=I\).

For a coefficient from \(z\) to \(y\), let \(\sigma(C)\) be a fixed additive constant plus the connected length of its tagged carrier. Include both canonical endpoint tags. The additive constant pays for the fixed short join between carriers at a common localization slot. In particular \(\sigma(C\vee C')\le \sigma(C)+\sigma(C')\). Use the norm
\begin{equation}\label{mb:eq-family-norm}
 \|\mathbf T\|_{s,\vartheta}
 =\max\left\{
 \sup_y\sum_{z,C}e^{s d(y,z)+\vartheta\sigma(C)}
       \|\phi_y T_C\phi_z\|,
 \ \sup_z\sum_{y,C}e^{s d(y,z)+\vartheta\sigma(C)}
       \|\phi_y T_C\phi_z\|
 \right\}.
\end{equation}
Here each coefficient norm includes the supremum over that coefficient's admitted local field and current data, taken before either endpoint sum. The geometry and declared read set are fixed during that supremum. The operator norms use the two declared normalized local spaces, including their physical measures. Supports and any derivative collars of \(\phi_y,\phi_z\) are included in the filing. For each product insert the square partition and retain the joined carrier for every intermediate slot. The two triangle inequalities and the separate absolute sums prove the product bound. One may keep the sum of the intermediate spatial distances in the weight of each word instead of its endpoint distance. This gives the same product estimate and bounds decay in the length of a connecting path through the successive slots. The norm of a summed matrix alone would give no such bound.

\begin{lemma}[Local positive Green families]
\label{mb:lem-positive-carrier-walks}
Fix a finite carrier rate \(\vartheta\). On the separated hierarchy, with a sufficiently large fixed separation, \(G_s\) and \(G_F\) have exact normally convergent connected families from local force to energy and from local energy dual to energy. The bounds in \eqref{mb:eq-family-norm} hold for a fixed \(s>0\), uniformly in mesh, depth and compatible volume. For \(G_s\) they also hold from unweighted \(L^2\) to \(\mathcal Z_s\), and from \(\mathcal Y_s^\flat\) to \(\mathcal Y_s\), in the physical local norms. Every term uses the prescribed positive operator compressed by zero extension. The same families are holomorphic on a common smaller weighted \(C^1\) tube around the admitted real reference.
\end{lemma}

\begin{proof}
We give the construction for \(\mathcal A=A_s\) or \(K_F\). On a set \(D_0\) of actual variables let \(E_{D_0}\) be zero extension. The positive form \(E_{D_0}^T\mathcal A E_{D_0}\) inherits the global lower bound in the energy of the zero extension. For a crossing row, the row value is evaluated on that extension. There is no inf-sup claim for a newly renormalized truncated row. Multiplication by \(e^{s\psi}\), with \(\psi\) a bounded one-Lipschitz function of the layer metric, preserves this support. Its conjugation of the restricted form changes the energy-dual bound by \(C|s|\). For the difference part this follows from
\[
 D(wu)(b)-w(b_-)Du(b)
       =\eta^{-1}\bigl(w(b_-)-w(b_+)\bigr)\operatorname{Ad}(U_b)u(b_+),
 \qquad w=e^{s\psi},
\]
with the same estimate for every translated component difference. For a row subtract its scalar weight at the top. Its oscillation in relative units is \(O(|s|)\), and the normalized row and column bounds control the two row factors. In the Wilson form the covariant curl has the same product rule, and the curvature coefficients have a finite shift and size \(O(\xi_y^{-2})\). These observations prove the conjugation estimate also on the domain of zero extensions. The same estimates for the changed coefficients, with size \(O(\varepsilon)\) in the tube, give the complex inverse by a Neumann argument. Distance weights now give columns bounded by \(Ce^{-s_0d(y,z)}\) from \(\mathcal F_{2,z}\), or from the local energy dual, to \(X_y\). For the latter source one tests the compressed variational equation against the weighted solution. Cutoff multiplication of a functional is defined by multiplying its test vector.

We next choose cutoffs using only the local geometry. Take a large \(m\) and use the comparison length \(\ell\) of \eqref{mb:eq-slow-length}, putting \(r_x=m\ell(x)\). Take \(m\delta\) small and \(\rho\delta\) sufficiently large depending on \(L\), with the ambient bare scale outside \(\Omega_1\). One can choose the common slope after the largest cutoff radius parameter used in the finitely many constructions. The estimates preceding the cutoff choice are uniform under such a decrease of the slope and increase of the separation. To verify the read collar, a point at level at most two below the level of \(x\) is separated from \(x\) by a complete intervening layer. Its distance term alone exceeds \(\xi(x)\) if \(\rho\delta\) is large. A point of the same or higher level cannot improve the zero-distance candidate. Hence all decisions for \(\ell(x)\) are made inside a ball of radius \(C_L\ell(x)/\delta\), including decisions about exterior labels. This collar is included even if it exceeds the support of the cutoff. Smooth difference estimates for \(\ell\) itself are unnecessary in this construction.

Perform the center counting and all local patch constructions on one compatible cover on which the enlarged balls and their decision collars embed. In particular a short fundamental period does not truncate the ambient fine centers in a cutoff ball. The periodic-sector calculation below returns the resulting operators and their family estimates to the original period.

Choose a fixed smooth nonnegative bump \(\varphi\), supported at radius \(1/2\) and equal to one at radius \(1/4\), invariant under the hypercubic group. At every ambient fine center \(x\) put
\[
 p_x=\eta^4r_x^{-4},\qquad
 N(v)=\sum_x p_x\varphi((v-x)/r_x)^2,\qquad
 \chi_x(v)=N(v)^{-1/2}\varphi((v-x)/r_x).
\]
For bond slots sample this formula at their geometric midpoints. If \(|x-v|\le a r_x\) for a fixed \(a\), then \(|r_x-r_v|\le am\delta r_x\). Thus the radii are comparable there. Counting fine centers in the smaller and larger radius balls gives \(c\le N(v)\le C\) and \(\sum_{x:v\in B(x,a r_x)}p_x\le C\). These constants are uniform because \(r_x\ge m\eta\). Differentiating the sampled formula takes differences only in \(v\), and each derivative of the bump costs \(r_x^{-1}\). The lower bound for \(N\) and the product rule give \(|D^a\chi_x|\le C_a r_x^{-a}\) for \(a\le2\), and in fact for any fixed \(a\) used below. We have the exact identity \(\sum_xp_x\chi_x^2=I\).

Let \(D_x\) be the actual variables in \(B(x,2r_x)\), and \(I_x=E_{D_x}(E_{D_x}^T\mathcal A E_{D_x})^{-1}E_{D_x}^T\). Include its full finite interaction collar. The support of \(\chi_x\) and each stencil reading it lie strictly inside the artificial boundary. The separation and local scale comparison ensure that the whole enlarged patch has at most two adjacent scales and that the only rough physical boundary is in the bare layer. The exact commutator satisfies
\begin{equation}\label{mb:eq-slow-family-commutator}
 \|[\mathcal A,\chi_x]u\|_{\mathcal F_{2,y}}
          \le \frac{C}{m}\bigl(\|\nabla_Uu\|_{2,B_y^+}
                              +\xi_y^{-1}\|u\|_{2,B_y^+}\bigr).
\end{equation}
For the interior covariant Laplacian, expand two differences by the product rule. The terms have either a first cutoff difference and a first difference of \(u\), or two cutoff differences and \(u\). The finite curvature terms give a single cutoff difference of relative size \(O(1/m)\). In the row mass, insert the cutoff value at each row location between the two row factors. On the complete row support the difference from this value is \(O(1/m)\); the two normalized absolute sums control the resulting \(L^2\) operator. At a physical bare boundary, the present matrix entries are \(O(\eta^{-2})\) and connect only finitely many neighboring bare-scale coordinates. Their cutoff differences are \(O(1/m)\), so the weighted force bound follows from the bare mass norm directly. No missing plaquette is filled in. The commutator reads \(u\) only on the cutoff stock and its stencil collar, away from the artificial boundary.

Set
\[
 \mathcal M=\sum_xp_x\chi_x I_x\chi_x,\qquad
 \mathcal T=\sum_xp_x[\mathcal A,\chi_x] I_x\chi_x .
\]
Zero extension and the interior support of \(\chi_x\) give
\(\mathcal A\mathcal M=I+\mathcal T\) exactly. The local inverse column and \eqref{mb:eq-slow-family-commutator} give exponential cell bounds \(Ce^{-s_1d(y,z)}\) for the first summand and \(Cm^{-1}e^{-s_1d(y,z)}\) for the second, each multiplied by \(p_x\). For a fixed \(y\) the sum of contributing \(p_x\) is bounded, by the overlap estimate on the enlarged patch. The same holds with \(z\) fixed. Sum the remaining exponential with the strict growth reserve of \cref{mb:metric-growth,mb:weighted-composition}. In particular the number of fine centers or the volume of the patch does not enter the bound.

File each summand on the patch, its complete row and stencil collars, the normalization and length decision collars, and its endpoint tags. Cubes of side \(M_0\xi_j\) at the adjacent scales, with \(M_0\ge C_L/\delta\), cover this set with at most \(B_L\) tags after bounded joins. Increase separation relative to \(LM_0\). The constant \(B_L\) is independent of the radius parameter \(m\). It follows that \(\|\mathbf T\|_{s,\vartheta}\le(C/m)e^{\vartheta B_L}\), and \(\|\mathbf M\|_{s,\vartheta}\le Ce^{\vartheta B_L}\). Choose \(m\) so the former is less than \(1/2\). This is consistent with first taking \(\delta\) small relative to \(m^{-1}\) and then increasing \(M_0,\rho\): the local inverse columns and overlap constants are independent of \(m,\delta,M_0,\rho\) once the fixed slope and rate conditions hold. The exact family is the expansion
\begin{equation}\label{mb:eq-positive-walk}
                    \mathcal A^{-1}=\mathcal M\sum_{n\ge0}(-\mathcal T)^n .
\end{equation}
The rightmost error may act on a local energy-dual source by the variational column estimate, while every subsequent error acts on the force space. The zero-error term uses the variational column directly. This gives the energy-dual assertion without changing the contraction.

For a site \(L^2\) source, the interior equation on the output cutoff is \(\Delta u=f-q^Ta_squ\), plus first and second cutoff terms. The normalized row sums bound the last term in local \(L^2\) by \(C\xi_y^{-1}\) times local scalar energy. Apply the local component estimate
\[
 \|\nabla_U Du\|_{2,B_y}
   \le C\{\|\Delta u\|_{2,B_y^+}
       +\xi_y^{-1}\|Du\|_{2,B_y^+}
       +\xi_y^{-2}\|u\|_{2,B_y^+}\}.
\]
This follows by applying the covariant curl-divergence identity to a smooth cutoff of \(Du\): its curl is curvature times \(u\), and the cutoff derivatives have size \(C/\xi_y\). Use the atomic estimate at a true bare stencil. Consequently the leftmost inverse with cutoff in \eqref{mb:eq-positive-walk} maps \(\mathcal F_{2,z}\) to \(\mathcal Z_{s,y}\) with an additional \(\xi_y^{-1}\). The unweighted input \(L^2\) contributes \(\xi_z\), and the intervening normalized factors keep only the ratio \(\xi_z/\xi_y\). Pay it once by the strict scale reserve. For \(\mathcal Y_s^\flat\to \mathcal Y_s\), scalar energy on \(B_y\) is comparable to \(\xi_y^{-1}\mathcal Y_{s,y}\) and its dual to \(\xi_y^{-1}\mathcal Y_{s,y}^\flat\). This gives the ratio \(\xi_y/\xi_z\), paid in the same way. These proofs use only the displayed second-difference estimates and the fixed real coercive references. They extend to the common tube using the same form and local coefficient bounds.
\end{proof}

\subsection{Coarse inversion by a grouped short power}

The matrices to be inverted next act on top or row slots with a bounded number of counting orthonormal coordinates per assigned cell. The number of fine variables in the reads can be arbitrarily large. Normalize the top coordinate by \(J_s=\Xi_s^2\) and the field row coordinate by \(J_f=\Xi_f^2\). In particular
\[
 \widehat F=FJ_s^{-1}=G_s q^T\Xi_s^{-4},\qquad
 \widehat M=\widehat F^T\widehat F,\qquad
 \widehat\Theta=\widehat M-(D\widehat F)^TG_FD\widehat F .
\]
The transpose here uses counting coordinates at the indicated normalized slot and the physical coordinates at every fine slot. The density of \(q^T\Xi_s^{-4}\) on the complete own pillar of a unit counting top is \(O(\xi_y^{-4})\). Its \(\mathcal Y_{s,y}^\flat\)-norm is bounded, since the \(\ell^2\) in the test and the physical volume factor in \(L^2\) each contribute \(\xi_y^2\). The positive site family therefore gives \(\widehat F:\ell^2_{\rm top}\to \mathcal Y_s\) and \(D\widehat F:\ell^2_{\rm top}\to\mathcal F_2\), with the same bounds for their physical transposes into the corresponding dual slots. Taking their products and using \(G_F:\mathcal F_2\to X\) gives the exact input families for \(\widehat M,\widehat\Theta\). Both have positive uniform real gaps by \cref{mb:lem-site-woodbury}.

\begin{lemma}[Inversion of a positive coarse family]
\label{mb:lem-coarse-carrier-inverse}
Let \(A=A^T\) have an exact connected family on such a coarse coordinate graph, with a real gap \(aI\le A\le b_0 I\) and a common complex tube. Suppose its coefficients admit the norm \eqref{mb:eq-family-norm} at some positive spatial rate and at a carrier rate with a fixed reserve over the desired output rate. The bound uses local coefficient suprema inside the sums. For a short
power combined entry, assume that every admitted real datum on its
declared read has a real extension with the same uniform gap and input
family bound. The extension may use another compatible volume or
hierarchy if it preserves all the geometry, phases, masks, weights,
and averaging reads of the entry. It need not preserve a periodic
identification of the read on the lifted lattice. A complex datum is
admitted on a smaller tube about a real datum with such an extension,
using the cutoff continuation in the proof. The matrices used below
have the real extensions described after this lemma. For sufficiently separated geometry, with fixed aligned reblocking, \(A^{-1}\) has an exact family at the output rate and a positive spatial rate. The construction also yields grouped powers for local inverse square roots and logarithms. The bounds are uniform in depth and compatible volume.
\end{lemma}

\begin{proof}
Write \(H=I-A/b\) with \(b>2b_0\), and symmetrize each coefficient by \(T\). On a periodic quotient first average the unmarked input family over its full deck orbit, carrying the tile phases with the tags. Carrier length is preserved by that translation. The short-carrier cutoff below therefore commutes with deck averaging, and its sum \(S\), as well as \(H\), preserves the periodic sector before powers are formed. This preserves the family bound because both row and column sums are used. In a provisional tiling of side \(M_0\xi_j\), split the family by its full length, writing
\[
 S=\sum_{\sigma_0(C)\le R_0} H_C,\qquad E=H-S.
\]
Use \(R_0\) large. The tail norm at an intermediate carrier rate is at most \(K e^{-cR_0}\) by the reserve at the input rate. At the real reference \(H\) lies between \(I/2\) and \((1-a/b)I\); for a sufficiently small tail \(0\le S\le q_* I\) for a fixed \(q_*<1\). The tail and these estimates hold for the full graph, including compressions already specified as the original input domain of \(A\).

An old tag of level \(j\) has metric diameter at most \(C_L M_0\) on the indicated adjacent layers. Its bounded joins have that same size bound. Thus every coefficient of \(S\) has its whole read set within \(C_L M_0(R_0+1)\) metric distance of either endpoint. For \(1\le j\le n\), combine the entries of \(S^j\) as matrices before taking their absolute values. Their range is at most \(B_j=C_LM_0j(R_0+1)\), and a combined entry reads only the two radius-\(B_j\) endpoint balls. This follows by following any short word through its successive intermediate cells and including the full carrier at each step. The balls include the geometric decisions and all boundary masks needed to know which short coefficients enter the sum. Truncation by endpoint separation would not give this read set.

Choose separation so \(2B_n\) is shorter than a full intervening layer crossing. Within these balls the metric and one of two adjacent lattice units are comparable, so each ball has at most \(C_L(1+B_j)^4\) coarse coordinates, with a fixed factor for the cell dimension. At a common real pair, the sum of squares of the block entries in one row or column of \(S^j\) is bounded by \(Cq_*^{2j}\), by the Hilbert norm of that matrix. Cauchy--Schwarz bounds its absolute row and column sums by \(C_L(1+B_j)^2q_*^j\). If the supremum over local data is taken separately for each entry, use the individual bound \(q_*^j\) and count all possible entries. This gives the bound \(C_L(1+B_j)^4q_*^j\), with the same fixed real spectral reserve at any admitted real extension. The distinction matters because the extremizing data of different entries can differ.

Take the individual estimate on the lifted entry before summing its
periodic images. An averaged phase is formed by translating the filing
rule and its field arguments together. Each phase is an exact family
for the same operator on the extended data, with the same bounds.
The completed field therefore need not be deck invariant. At periodic
data the averaged short matrix is deck equivariant, as required for
the subsequent projection. Summing over input lifts at one fixed
output lift then gives the same bound for the projected row, and the
transpose argument gives the column bound.

For use on local complex domains, take a collar at the scale of the assigned cells around the full read set. Sum the squares of the smooth localization profiles whose supports meet the read set to obtain a cutoff equal to one there and vanishing before the outer edge of the collar. Its first differences are \(O(1/\xi)\) by bounded overlap and comparable scales. Multiply the local relative \(C^1\) field by this cutoff and extend by zero, relative to the fixed real reference. The full weighted \(C^1\) norm costs a fixed factor by the product rule. Cutoff extension of a current uses only its amplitude bound. The collar and the local tube divided by this fixed factor ensure that the data extend to the admitted global tube. This construction uses the fixed underlying hierarchy and does not change geometry outside the read set. On a weighted \(C^{1,\alpha}\) domain, use the same smooth cutoff with \(|D^2\chi|\le C\xi^{-2}\). Its scaled first-difference H\"older norm is bounded by interpolation between \(|D\chi|\le C\xi^{-1}\) and this second-difference bound. The discrete product rule then bounds the scaled \(C^{1,\alpha}\) norm of \(\chi A\) by a fixed multiple of that of \(A\), including pairs crossing the collar where the cutoff vanishes. Thus the extension also preserves the stronger domain used for the fourth-order inverse. On it the complex short matrix obeys \(\|S\|_{2\to2}\le q'<1\) by reducing the original tube. Its individual complex entry bounds therefore use \(q'\) in place of \(q_*\), even when the local data are varied separately. More precisely, after fixing \(n,R_0\) the difference of each short power from its value at the real reference is holomorphic in its local field, and its entry variation is at most \(C_j\varepsilon\) in a smaller radius \(\varepsilon\). This follows by telescoping \(S^j-S_*^j\), or by the one-complex-variable Cauchy estimate on each line in the local tube. The polynomial entry count above gives a uniform row and column bound \(\Delta_j^+\) for this supremum. The tube can be chosen so \(\Delta_n^+\le(q')^n\) and \(\sum_{j<n}\Delta_j^+\le(q')^n\). There is no factor from the number of fine coefficient variables.

The long part is the exact family
\[
 H^j-S^j=\sum_{i=0}^{j-1} S^i E H^{j-1-i}.
\]
Each term keeps its joined carrier and has total family norm at most \(CjK^j e^{-cR_0}\) by the product estimate. Choose \(R_0=D n\) for fixed \(D\) with \(cD>2\log(1+K)+2\). Next choose aligned parent tiles of side \(M_c\xi_i\), \(M_c\) a multiple of \(M_0\), with \(M_c\ge C_L(B_n+1)\). File each combined short entry on the full two balls. It then has at most a fixed number of new tags. For a long carrier traverse its old connected tree in at most twice its length. Divide the traversal into pieces of at most \(M_c/M_0\) edges. Each piece, with its old collars, meets a bounded number of parent cubes on at most two adjacent levels if \(\rho\gg LM_c\). Hence the new tag count is at most
\begin{equation}\label{mb:eq-reblock-count}
                    C_L(1+(M_0/M_c)\sigma_0(C)).
\end{equation}
Endpoint tags, level tags and short joins are included in this estimate. The parent mapping does not enlarge the family norm beyond a fixed factor once its carrier weight is paid from the reserved old rate. Likewise, for any fixed parameter radius \(\mathcal S>2\), choose \(M_c\) so the factor \(\mathcal S^{\# C}\) of a carrier is paid by the reserve in \eqref{mb:eq-reblock-count}. The factor for a short combined entry is then bounded independently of \(n\).

In the new filing of \(H^n=S^n+(H^n-S^n)\), the norm of the short part is at most \(C_L(1+B_n)^4(q')^n\), up to the fixed tag and parameter factors. The long part is bounded by \(CnK^n e^{-cDn}\). For the short spatial weight take \(s\le \log2/(B_n+1)\), still smaller than the input spatial rate. Choose \(n\) large enough that the combined family has norm less than \(1/2\); then choose \(M_c,\rho\) as above and impose the strict spatial reserve for all scale ratios used in the construction. The bound remains valid under that increase of separation. The constants used to choose \(n\) are from the initial positive construction and the real gaps, and do not depend on the final parent tile. Finally choose the smaller common complex tube for the finitely many short powers. If no local coefficient suprema are needed, the quadratic power of \(1+B_n\) from the common-pair row estimate suffices, but the preceding stronger construction will be used below.

Use the same filing for \(H^j\), \(j<n\), and then
\begin{equation}\label{mb:eq-coarse-carrier-inverse}
 A^{-1}=b^{-1}\left(\sum_{j<n}H^j\right)\sum_{v\ge0}(H^n)^v .
\end{equation}
At complete data this is an ordinary matrix identity. With the joined carriers the geometric series converges absolutely in the family norm. All local coefficients are holomorphic, since the positive primitive inverses and the finite short products are holomorphic on their declared read sets, and the series is normally convergent. The same grouping can be used for any scalar power series in \(H\) with bounded coefficients, in particular the square root and logarithm constructions below.
\end{proof}

For the actual coarse matrices the short-read extension requires only
a real energy gap. Enlarge the two endpoint balls and all declared
decision and stencil collars to an axis-parallel cube of radius
\(R\xi_y\) and, where needed, one further fixed enlargement. Here
\(R\) is fixed after the short word length and the old carrier cutoff.
Increase the separation so that the enlarged cube contains at most
two adjacent assigned scales. Include every complete plaquette and
every projected branch stock in that collar. Its radius increases the
short read by a fixed factor and costs a bounded number of parent tags.

If the enlarged cube is interior, keep the hierarchy, all row weights,
and all selection masks on a compatible lift where the cube embeds.
Use the fine coordinate comb to choose a unitary frame on the cube.
Each transformed link is the product of unitarily transported
plaquette holonomies in a strip between two comb paths,
with at most \(C R\xi_y/\eta\) factors.
Thus its principal potential \(a\), defined by \(U=e^{i\eta a}\),
satisfies \(\|a\|_\infty\le C_L R\kappa/\xi_y\) on the cube.
Reserve an inner collar containing the later relative complex cutoff
and all its derivative stencils. Take a scalar cutoff \(\chi\)
equal to one on this entire collar, including every read link,
and vanishing in the outer collar before the cube boundary,
with first differences bounded by
\(C/(R\xi_y)\). Extend by \(e^{i\eta\chi a}\) in this frame and by
the identity beyond its support. For a complete plaquette \(p\) in
the cutoff region, the product expansion gives
\[
 \eta^{-2}|\log U_{\rm ext}(\partial p)|
 \le C\big(\eta^{-2}|\log U(\partial p)|
       +\|D\chi\|_\infty\|a\|_\infty+\|a\|_\infty^2\big)
 \le C_L(\kappa+R^2\kappa^2)\xi_y^{-2}.
\]
Indeed in the linear sum of the four oriented link potentials one may
subtract a common plaquette value of \(\chi\); the remainders are
bounded by one link length times the cutoff derivative. Both
plaquette product remainders are quadratic in the link angles.
All stencils meeting the cutoff are complete within the enlarged
cube. Undo the comb frame throughout the inner collar, extending that unitary site
gauge arbitrarily on the lift. Curvature is unchanged by this last
operation. The small-curvature recurrence on complete backward
averaging stocks admits the real branches, including the projected
branch tests. Copy the independent current on the read and set it
to zero elsewhere, which preserves its weighted amplitude bound.
Reducing the local curvature ceiling after fixing \(R\) puts this
extension in the already proved real energy-gap class.

If the enlarged cube meets the physical exterior, separation makes
its entire active part bare, including the decision collars. Take a
larger level-zero aligned box within the supplied collar and restrict
the original active set to it, with no higher regions. Retain the
original action-boundary convention and copy the data on all the
resulting stencils. For an ambient compression keep the complete
ambient crossing stencils in the enlarged collar. Every retained
curvature test was a test in that collar. The new box boundary is
beyond the entire read and its relative complex cutoff and derivative
collar. Inside the read the layer length stays
\(\eta\), the site and bond rows are the identity, and the physical
boundary and every cutoff and mask decision are unchanged. Hence the
short coefficients are those of this admitted real bare problem.
Proper outer coefficients have their enlarged short reads in the
interior top stock by the assumed outer separation.

For successive coarse inversions make these choices in their
dependency order: first the site matrices, then the row matrix,
and then the outer precision. At each stage choose its finite
read radius and reduce the local real curvature ceiling enough
that its completion lies in every earlier input family class.
Choose nested complex tube radii in the same order. Earlier
gap and family bounds remain valid on the smaller classes.
The maximum of this finite list of read radii fixes the required
separation. None of these choices depends on the mesh or the
number of averaging levels.

\mbneedspace{10\baselineskip}
\subsection{Actual current families and the reconstruction slots}

\begin{theorem}[Connected coefficients for the named propagators]
\label{mb:thm-current-families}
Fix a finite carrier rate and choose the aligned tiles, separation and smaller common complex coefficient tube of the preceding lemmas. For a sufficiently small touching current the maps \(P,R,\Lambda,S,G_0,G_J,G_{1,J}\), both normalized Schur inverses and lifts, the two multiplier readouts and the coupled response have exact connected families. In the norms of \eqref{mb:eq-family-norm} they are bounded in their stated energy, balanced site, force and row slots. The same holds for \(\mathfrak G\), for the zero-current constrained response \((I-H_{0}Q)G_0\) with \(H_0=G_0Q^T(QG_0Q^T)^{-1}\), and for the mass-removed row response. These assertions include local \(X^*\)-sources and local \(L^2\)-force sources.
\end{theorem}

\begin{proof}
Obtain \(\widehat M^{-1}\) and \(\widehat\Theta^{-1}\) at a stronger intermediate carrier rate by \cref{mb:lem-coarse-carrier-inverse}. The site and field positive families used in their input are first obtained at a still stronger rate. Use aligned parent tiles throughout. Compose to get \(P=\widehat F\widehat M^{-1}\widehat F^T:L^2\to \mathcal Y_s\) and \(R=I-P\) on \(\mathcal Y_s,L^2\). The map \(\Lambda=G_s RD^T\) takes \(X\) to \(\mathcal Z_s\); then \(D\Lambda:X\to X\). The first part of this composition uses \(D^T:X\to L^2\). At a force input \(f\), use instead \(\|D^T f\|_{\mathcal Y_s^\flat}\le C\|\ell f\|_2\), already proved by the test product rule. Thus \(DRG_s=\Lambda^T:\mathcal Y_s^\flat\to\mathcal F_2\) and \(S^T=I-DRG_sD^T\) is bounded on \(\mathcal F_2\) in the family norm. These bounds arise from the actual site family, including its balanced \(\mathcal Y_s^\flat\to \mathcal Y_s\) bound. The other orientation, \(S:X\to X\), follows from the \(\mathcal Z_s\) output. Formula \eqref{mb:eq-zero-current-woodbury} gives \(G_0\) from \(\mathcal F_2\) or \(X^*\) to \(X\), with \(G_F:X^*\to X\) also used in its rightmost factor.

Before taking local coefficients, write the exact current completion using \(W=W[J^{\rm t}]\). Indeed every occurrence of \(W\) in \eqref{mb:eq-current-completion} has a site factor \(\Lambda\), or its transpose, and \(q\Lambda=0\) makes that site field zero on bare singletons. Hence replacing \(J\) by \(J^{\rm t}\) leaves the complete \(B[J]\) unchanged. The multiplication \(W[J^{\rm t}]\) takes arbitrary \(\mathcal Z_s\) inputs to \(\mathcal F_2\) with norm \(C\gamma_{\rm t}\). Its transpose takes \(X\) to \(\mathcal Y_s^\flat\) by the weighted test pairing. Equation~\ref{mb:eq-current-completion}, with every \(W\) replaced by \(W[J^{\rm t}]\) and its \(D^TW\) term interpreted as a divergence source, therefore gives a family for \(B[J]:X\to\mathcal F_2\) with norm \(C_\vartheta\gamma_{\rm t}\). Every such coefficient reads the local current bond. Hence \(G_J=G_0\sum_{n\ge0}(B[J]G_0)^n\) converges in the family norm after decreasing the fixed touching-current ceiling. At an \(X^*\) input the rightmost \(G_0\) acts on that slot and every subsequent loop acts on the force slot. The first difference \(G_J-G_0\) has the additional factor \(\gamma_{\rm t}\).

Normalize \(\mathcal Q_n=J_f a_f^{1/2}Q\). The actual injection \(\mathcal Q_n^T e_c\) from a unit counting datum has density \(O(\xi_c^{-3})\) on its full row carrier, so has bounded \(\mathcal F_2\) norm and a bounded local family. The transpose takes \(X\) to counting rows by the row bounds. Thus \(\mathcal T_0=\mathcal Q_n G_0\mathcal Q_n^T\) has the exact family and positive real gap proved in \cref{mb:thm-coarse-schur}. Apply \cref{mb:lem-coarse-carrier-inverse} once more at the desired rate. The current difference \(\mathcal T_J-\mathcal T_0\) has norm \(C_\vartheta\gamma_{\rm t}\), so an ordinary family Neumann series gives \(\mathcal T_J^{-1}\), and
\[
  H_J=G_J\mathcal Q_n^T\mathcal T_J^{-1}J_fa_f^{1/2}.
\]
This formula keeps the row weight on the input.

For the multiplier first use the exact split of \cref{mb:cor-current-readouts}: pair the lift with \(J^{\rm t}\), and add the explicit bare-row readout of the nontouching current. Test each term of the lift family on the datum \(b^{(c)}\) and split its pairing with \(J^{\rm t}\) by cells. A coefficient for \(\xi_c^3(H_J^T J)(c)\) has the carrier of the lift, joined to the actual current cell \(B_y\), with root \(c\). Its absolute size is bounded by the normalized column norm times \(C\gamma_{\rm t}\), since \(\|\ell J^{\rm t}\|_{2,B_y}\le C\gamma_{\rm t}\). Sum first over all such current cells and carriers at the fixed root. The weighted column bound gives a uniformly summable rooted family. The added bare-row readout is local and vanishes at nonbare roots. Thus the touching-current estimate holds coefficient by coefficient before the root sum. The two-leg estimate \eqref{mb:eq-two-leg-energy} now gives a family for \(C_\omega\) with norm \(C_\vartheta\gamma_{\rm t}\) from \(X\) to \(\mathcal F_2\), since its actual two legs and their row root read only a bounded star. Consequently \(S^T C_\omega S\) has that same bound. Inverting the second variant by the corresponding force-space Neumann series gives \(G_{1,J}\). For its row inverse use the difference from \(\mathcal T_0\) and the same normalized row formula. The same rooted pairing gives the second readout. The first readout was used in constructing this second operator.

All remaining claimed families are finite products of these families. For \(\mathfrak G\) use \(D^T:X\to L^2\), \(R\) on \(L^2\) and \(D:L^2\to X^*\), as well as the lift. For the datum-to-multiplier response the physical angle is \(\Xi_f b\), and the physical current amplitude is \(\Xi_f^3\mu\). Their family identity is
\begin{equation}\label{mb:eq-mass-removed-family}
 \Xi_f^3\{(QG_iQ^T)^{-1}-a_f\}\Xi_f^{-1}
      =(\Xi_f a_f^{1/2})(\mathcal T_i^{-1}-I)(\Xi_f a_f^{1/2}),
 \qquad i=J,\ (1,J).
\end{equation}
This follows by writing \(T_f=J_f^{-1}\mathcal T_iJ_f\). It retains the row mass subtraction. The constrained response and the remaining blocks of \eqref{mb:eq-field-saddle} follow by inserting their actual local row maps. This completes the claimed constructions.
\end{proof}

\begin{corollary}[Comparison of determining data]
\label{mb:cor-determining-comparison}
Use one filing rule, including its aligned tile phase, on a common lift of two admitted geometries. Suppose the two systems agree on a given open read region, including all fine links, independent currents, original inner plaquettes, determining rows and their weights, domain masks, and geometry labels there. The coefficients with full carriers in that common region agree. The difference of the localized complete operators on matching endpoint slots is bounded by the absolute sums of the families whose carriers leave it. Write $d_i$ for the layer metric of system $i$ and $r_i$ for the $d_i$-distance from the output slot to the complement of the common read region, after subtracting the fixed collar allowances. The difference is bounded by
\[
 C\sum_{i=1}^2 e^{-s d_i(y,z)-c(r_i)_+/M_c}.
\]
Here $s,c>0$ are reserved rates. If the common region contains a radius-$r$ ball measured before first exit, then $r_i\ge r$ in both systems, since their geometry agrees along every path before that exit. One may then use $2C e^{-s\min(d_1,d_2)(y,z)-cr/M_c}$. A single metric suffices when the two hierarchies agree globally.
\end{corollary}

\begin{proof}
The read rule for a primitive coefficient first decides its comparison lengths and normalization on its full geometry collar. It then uses only the actual compressed finite matrix, its cutoff and endpoint data. The row trace, difference, local current and constraint coefficients have the same literal property. In a short combined entry use the two endpoint balls enlarged by the full data collar. They can be fixed in Euclidean coordinates using radius \(C_L B_j\xi_y\) at the \(y\) endpoint and \(C_L B_j\xi_z\) at the other. Every metric path of length at most \(B_j\) stays at the two indicated adjacent scales by separation, so these Euclidean balls contain all short reads. This fixes their construction without needing data at a remote alternative path. Joined products and the rooted current pairing preserve equality when all those reads agree. Thus the terms match individually, including whether they are present. All families converge absolutely, so their common terms cancel in the difference. A connected carrier leaving the common radius has at least \(c r/M_c\) tags, up to the fixed endpoint and collar allowances, because each tag and short join has metric diameter at most \(C_L M_c\). Reserve part of the carrier exponent for this bound separately in each hierarchy, using its own metric and first-exit distance. Add the two remaining absolute sums. The same argument with the input slot gives the corresponding column comparison.
\end{proof}

The families can be used on periodic quotients without cutting a literal path at a periodic face. First construct all primitive and coarse stages on one common compatible lift large enough for the finite patches and aligned tile choices. The real positive compressions and all local coefficients are defined there. The complete operators on that lift commute with deck shifts. For example, choose isometric pullback of base coordinates into the deck-invariant subspace, with the normalization by the square root of the number of sheets. Each complete deck-equivariant matrix intertwines this injection, as does its bilinear transpose. An individual local coefficient need not preserve the periodic sector; its full deck orbit does. Inverses and products of the complete equivariant matrices preserve that sector by uniqueness. Average the families over simultaneous deck shifts, translating the filed tags and phases as well, before projection. Convexity preserves both absolute sums. A projected row sum is bounded by the lifted row sum over all input lifts at one fixed output lift. Deck covariance gives the corresponding bound for columns. The projected tag length and distances can only decrease, while the full lifted reads remain declared. There is no factor from the number of sheets.

One may use a common filing on all compatible periods at fixed levels. Fix the aligned tag phases on the universal lattice and use their finite uniform average under translations by the top mesh, modulo the final tile width. Include the phase as part of each coefficient's label. For any finite collection of words the lifted coefficients are independent of period once all their reads embed. The family bounds, including the old carrier tail, dominate the projection sums over longer words. Thus on any period the same rule is obtained by projection of the lifted coefficients, with winding reads retained. The finite-cover sums agree with the complete operator by the preceding periodic-sector identity; absolute convergence gives the equality as the lift increases. This argument also gives literal term agreement on common lifts. A period so small that a chosen support wraps is treated by projection, with its winding and periodic identifications in the coefficient. It is not treated as an embedded copy of that support.

In real gauge changes the coefficient on each output and input slot is conjugated by its induced color map. For example a reversal of a fine bond sends its left logarithmic increment to \(-\operatorname{Ad}(U_b^{-1})A(b)\); on a determining row the same expression uses the actual averaged row link. For a site the color map has no reversal sign. These maps preserve the declared bilinear pairings because \(\operatorname{Ad}(g)^T=\operatorname{Ad}(g^{-1})\) and each isometry preserves the slot measure. All primitive cutoffs are scalar and each term reads the actual complete rows needed for these endpoint factors. Covariance therefore holds at admitted images, with the tags and their phases carried along. A single fixed phase gives termwise isometry covariance for isometries preserving that phase. The complete sums transform under every admitted isometry of the graph and domains. Complex covariance is this bilinear identity in the appropriate charts.

\begin{corollary}[Independent carrier parameters]
\label{mb:cor-carrier-parameters}
After a sufficiently large fixed aligned reblocking, assign a complex parameter \(s_\Delta\) to each parameter cube and multiply a joined carrier coefficient by the product of the parameters of its distinct cubes. On every prescribed fixed polydisc \(|s_\Delta|<\mathcal S\), the resulting family converges normally, with the same physical local spaces and a reserved rate. At all ones its sum is the original operator. Setting parameters outside a wall to zero retains precisely the terms filed in the wall, in addition to any explicitly held guard. Separate wall components have no terms connecting them after their guards and complete read sets are included and intersecting components are merged.
\end{corollary}

\begin{proof}
Apply \eqref{mb:eq-reblock-count} to each old connected carrier.
Increase the parameter tile width until the cost per old tag,
\[
                    (\log\mathcal S)C_L M_0/M_c,
\]
fits in the old carrier reserve. The constant term costs only \(\mathcal S^{C_L}\). This yields the absolute row and column bounds with the coefficientwise suprema on every smaller field tube. Normal convergence follows by the Weierstrass criterion there. For a product of operators the parameter monomial of the joined coefficient uses its union of cubes. Repeated visits to a cube do not create new factors in this convention. The product of the absolute majorants for the constituents bounds this marking. Evaluation at all ones preserves operator products. A surviving coefficient at a zero wall parameter has its entire filed support in the allowed wall and held guard. Its connected carrier proves the splitting. Corollary~\ref{mb:cor-determining-comparison} proves equality on matched wall reads.
\end{proof}

On a short period, descend the unmarked families first, including their full lifted read collars. Choose the parameter partition on the base torus afterward. At a fixed level measure lengths in that level's grid units and write \(W=M_c\). For an axis with \(N\ge3W\) sites, use \(\lfloor N/W\rfloor\) cyclic intervals and distribute the remainder so their lengths lie between \(W\) and \(2W\). If \(N<3W\), use \(\min(2,N)\) balanced cyclic intervals. Products of these intervals are the parameter boxes; on an exactly divisible long period they may be cubes. A ball of radius \(CW\) meets at most \(\lceil2C\rceil+3\) intervals in a long axis and at most two in a short one. Apply this count to the traversal pieces in the proof of \eqref{mb:eq-reblock-count}. Separation leaves at most two adjacent levels in each piece, so the same bound holds with a changed \(C_L\), uniformly in rectangular periods. In particular, on one level with all four axes short there are at most sixteen parameter boxes. Mark a distinct base box once, regardless of the number of lifted visits that project to it. The parameter partition and phase are included among the matched data for wall comparison. The matrix-function extension below uses its separately stated ordinary-product convention.

The ordinary uniform estimates, such as \cref{mb:thm-coarse-schur}, use a coefficient and current ceiling determined by the real form gaps and the geometric growth reserve, before introducing any walk tile. The stronger families at a prescribed carrier rate use a finite list of tile choices and may shrink the complex tube and current ceiling after those choices. The real curvature ceiling for independently varied local reference data may also be reduced after the short-read radius is fixed, as in the real extension proof above. The ordinary summed estimates keep their previously chosen ceilings. All localized bounds and sums then remain independent of mesh, depth and compatible volume. A pointwise entry bound in an unnormalized coordinate can still contain its physical mesh powers. Such an entrywise bound, and a whole-volume bound used only to verify convergence, do not enter any of the weighted local bounds or rates.


% Construction source: audit/strong-covariance.tex
\section{Strong local Green and restoration estimates}
\label{mb:sec-strong-green}

Use the fixed real comparison frames on complete nonbare boxes, and the bare atomic convention at true exterior stencils. Take finitely many nested boxes at each nonbare cell, with enough room for the row carriers and the successive interior cutoffs. Their radii are fixed multiples of \(\xi_y\) and lie in at most two adjacent layers by separation. The coefficient premises for the finite-\(p\) estimates are bounded scaled potentials and scaled first differences in these frames. The estimates with two H\"older derivatives below require a scaled \(C^{1,\alpha}\) potential bound on the real reference and the relative perturbation. One can use an admitted strong frame or an admitted Coulomb frame with the actual current hypotheses of \cref{mb:cor-current-holder}. The independent current in the propagator is a separate variable.

For a field force and field output use
\[
 \|f\|_{\mathcal F_{p,y}}=\xi_y^{3-4/p}\|f\|_{p,B_y^+},\qquad
 \|u\|_{\mathcal A_{\alpha,y}}=
     \xi_y\|u\|_\infty+\xi_y^2\|\nabla_Uu\|_\infty+
                   \xi_y^{2+\alpha}[\nabla_Uu]_\alpha^U .
\]
The H\"older comparisons use short coordinate transports on complete local boxes. On the complex tube they are equivalent to the same comparisons in the fixed reference frames. A site function with amplitude of grade two has scaled first-order norm
\[
 \|v\|_{\mathcal V_{p,y}}=
      \xi_y^{2-4/p}\|v\|_p+\xi_y^{3-4/p}\|Dv\|_p .
\]
A site source \(D^Tf+h\) is measured by the infimum over representations of \(\|f\|_{\mathcal F_{p,y}}+\xi_y^{4-4/p}\|h\|_p\), with the complete source collar. Its \(D^T\) is variational. A bounded site source has grade-four amplitude \(\sup_y\xi_y^4\|h\|_{\infty,B_y^+}\).

\begin{lemma}[The interior Wilson stencil]
\label{mb:lem-wilson-component-stencil}
On a complete nonbare stencil, with the Wilson quadratic normalization for which its flat Hessian is the curl energy, the actual operator \(H_W+DD^T\) equals the component covariant Laplacian plus a finite-shift kernel \(Z_U\) with absolute row and column bounds \(C\kappa\xi_y^{-2}\) on a scale-\(\xi_y\) box. The identity and bounds hold on the admitted complex tube with the corresponding curvature ceiling.
\end{lemma}

\begin{proof}
For a plaquette, transport the four tangent insertions to its initial vertex and divide their oriented sum by \(\eta\). Squaring this linear curl gives the Hessian at unit plaquette holonomy. The exact Wilson insertion formula gives the difference as a sum of products of two transported insertions paired with \(U(\partial p)-I\), with coefficients \(O(\eta^{-2})\) after division by the fine volume. It is therefore a finite shift of size \(C\kappa\xi_y^{-2}\). Replacing the four path transports in this curl by the standard covariant forward curl changes a coefficient by \(O(\eta^{-1}|U(\partial p)-I|)\): the only alternative routes around the square differ by the plaquette holonomy. Multiplying the two curl matrices gives mixed coefficients \(O(\eta^{-2}|U(\partial p)-I|)\), since each curl entry has size \(C/\eta\) and each stencil has a bounded number of entries. This step takes no difference of the curvature.

For the standard curl write \(T_\mu=\operatorname{Ad}(U_\mu)\tau_\mu\), where \(\tau_\mu\) shifts by one forward step. Its curl square plus the gradient-divergence square has diagonal \(\eta^{-2}\sum_\nu(2-T_\nu-T_\nu^{-1})\). An off-diagonal entry is a signed difference of the two ordered shifts \(T_\mu T_\nu^{-1}\) and \(T_\nu^{-1}T_\mu\), divided by \(\eta^2\). Those two routes differ by one elementary plaquette and hence give a kernel of the same curvature size. Bounded incidence gives the column bounds as well. All transposes here use transport inverses, an algebraic identity also for complex links.
\end{proof}

\mbneedspace{16\baselineskip}
\begin{lemma}[Strong columns of the positive primitives]
\label{mb:lem-strong-positive-columns}
For \(0<\alpha<1\), finite \(p>4/(1-\alpha)\), and a fixed carrier rate, \(G_F:\mathcal F_p\to\mathcal A_\alpha\) has a connected family in the norm \eqref{mb:eq-family-norm}. For every fixed \(2<p<\infty\), \(G_s\) has a family from the variational site source just defined to \(\mathcal V_p\). From bounded grade-four site sources it has a family to
\[
       \xi_y^2\|u\|_{\infty,B_y^+}
       +\xi_y^3\|Du\|_{\infty,B_y^+}
       +\xi_y^{3+\alpha}[Du]_{\alpha,y}^U .
\]
For site columns \(\widehat F_z=G_sq^T\Xi_s^{-4}e_z\) this gives the point and H\"older bounds with factor \(e^{-s d(y,z)}\) and the physical scales in the displayed norm. Constants are uniform in mesh, depth and compatible volume and hold on a common smaller complex coefficient tube.
\end{lemma}

\begin{proof}
Start with a positive local compression from \cref{mb:lem-positive-carrier-walks}. The normalized \(\mathcal F_p\) source embeds into its \(L^2\) force slot by the cell volume bound. Its local energy column has the exponential rate \(s_0\). On an interior box the equation is the component covariant Laplacian by \cref{mb:lem-wilson-component-stencil}, with lower-order curvature and row terms. In a fixed comparison frame subtract the ordinary lattice Laplacian. Pairing the forward and inverse shifts gives first-order coefficients \(O(\|a\|_\infty)\) and zeroth-order coefficients \(O(\|\nabla a\|_\infty+\|a\|_\infty^2)\) in rescaled units. For example the zeroth term from one direction is the sum of the forward \((\operatorname{Ad}e^{i\eta a}-I)/\eta\) and its backward inverse counterpart, divided by \(\eta\); their linear terms leave one first difference of \(a\). The curvature term of \cref{mb:lem-wilson-component-stencil} is a bounded finite-shift multiplication.

The row mass has bounded point forcing on each box. Indeed the complete normalized row has coefficient density \(O((\eta/\xi_c)^4)\), and hence \(Qu(c)\) is bounded by \(C\xi_c^{-2}\|u\|_{2,\operatorname{supp} Q_c}\). The physical transpose of \(Q\) on one row has bounded point density and \(a_f(c)\asymp\xi_c^{-2}\). The result is a point force of grade three controlled by the cell energy. The same argument for \(q^Ta_sq\) gives grade four from the balanced site energy.

In the unit-scale rescaled box the field \(H^1\) estimate gives the lower-order first differences in \(L^2\). The local finite-\(p\) estimates give \(W^{2,2}\), and strict lattice Sobolev embedding gives gradient \(L^3\). A second application at exponent \(3\) gives \(W^{2,3}\), gradient \(L^8\) and bounded field. For \(p>8\) use one further application at \(8\) before the final one at \(p\). For \(p\le8\) the final \(p\) application is already valid. Every lower-order term is estimated on the next larger complete box. The strict Morrey estimate gives \(\alpha<1-4/p\), with the stated gradient and potential grades after rescaling. These are the finite-\(p\) and covariant interior estimates of \cref{mb:prop-flat-lp,mb:prop-local-scalar}. The original energy and source terms keep their exponential column bounds since each box enlargement costs a bounded metric distance. The bounded-site-source case uses the same bootstrap componentwise. A true bare stencil uses its finite number of coordinates and \(\xi_y=\eta\), so the scaled value, first difference and H\"older quotient follow from its \(L^2\) value bound directly.

For a site divergence source use the balanced energy column. In a comparison frame the connection lower-order terms may be written as divergence of a bounded coefficient times \(u\), plus bounded multiplication of \(u\). For a forward first difference this is the identity
\[
 a(x)\frac{u(x+\eta)-u(x)}{\eta}
 =\frac{a(x)u(x+\eta)-a(x-\eta)u(x)}{\eta}
     -\frac{a(x)-a(x-\eta)}{\eta}u(x)
\]
in one coordinate direction. The covariant divergence differs from the flat one by another bounded multiplication of the specified source. Smooth cutoff multiplication preserves the divergence representation by the same product rule. Starting with \(H^1\) gives \(u\in L^3\); the local divergence estimate at \(3\) gives \(W^{1,3}\) and hence \(u\in L^8\). If the desired \(p\) is smaller, stop at \(p\). Otherwise apply the estimate at \(8\) and use the resulting bounded \(u\) before the final finite exponent. This gives \(\mathcal V_p\) with the same exponential rate after the fixed scale-ratio reserve. It differentiates no input \(f\).

Use these stronger local columns in the slow parametrix. Each interior commutator has either \(D\chi_x\,\nabla_Uu\), \(D^2\chi_x\,u\), a bounded curvature coefficient times a cutoff difference, or a complete-row cutoff difference as in \eqref{mb:eq-slow-family-commutator}. The same bounds hold in \(\mathcal F_p\), in the site source class, and in the bounded site force class respectively. In the site divergence case the commutator is even a grade-four \(L^p\) source from \(\mathcal V_p\). All these local error norms are at most \(C_{p,\alpha}/m\) after rescaling. The square partition and weighted center sum from the positive walk proof yield a contraction on each input force space once \(m\) is fixed sufficiently large. Thus \eqref{mb:eq-positive-walk} gives exact strong families. For \(\widehat F_z\) the input density is grade four, supported on the complete own pillar, so the bounded-site-source conclusion gives its columns. The stated coefficient perturbation preserves every estimate on the smaller tube. The cutoff and collar choices are independent of the fine number of variables.
\end{proof}

Two useful physical conversions follow. In the unnormalized site array the inverse of \(qG_s^2q^T\) has size \(C\xi_y^{-4}e^{-s d(y,z)}\): it equals \(\Xi_s^{-2}M^{-1}\Xi_s^{-2}\) and the change from normalized counting coordinates contributes \(\xi_y^{-2}\xi_z^2\). Also the counting fine array of \(D P D^T\), at fine bonds based in \(B_y,B_z\), has size
\begin{equation}\label{mb:eq-projection-point-kernel}
 C\,(\eta/\xi_z)^4\xi_y^{-2}e^{-s d(y,z)} .
\end{equation}
In fact its physical kernel is the composition of the point columns \(D\widehat F\), the normalized \(\widehat M^{-1}\) and their transpose columns. Each end contributes its grade-three amplitude and the intermediate counting sums are exponential, giving \(C\xi_y^{-3}\xi_z^{-3}\) times the exponential. Multiply by \(\eta^4\) for the counting array and use the reserved scale ratio \(\xi_z/\xi_y\). With \(P,DP,PD^T\) use the same argument with grades two or three according to the differentiated ends. In particular these are smoothed projection kernels, with no diagonal Coulomb singularity.

\subsection{A fourth-order construction for strong restoration}

The formula \(G_sR\) gains two H\"older derivatives on a H\"older source, but its second-order positive pieces can have an averaging-mass jump at a pillar face. To obtain the same gain for each local coefficient, use the actual fourth-order positive matrix
\[
  q^*=q^T\Xi_s^{-4},\qquad
  A_4=\Delta^2+q^*q,\qquad G_4=A_4^{-1}.
\]
Here \(q^*\) pairs the raw top values with counting measure. The site form norm is \(V_2=\|\Delta u\|_2+\|\ell^{-1}Du\|_2+\|\ell^{-2}u\|_2\). It controls \(\nabla_U Du\) on complete stencils by the curl-divergence estimate, with the bare atomic convention elsewhere. The real form \(\|\Delta u\|_2^2+\|qu\|_{\rm count}^2\) is equivalent to \(V_2^2\). For the lower bound subtract \(E_s qu\). The remainder has zero averages and its zero-average estimate is bounded by its \(L^2\) Laplacian. The lift has \(V_2\)-norm at most \(C\|qu\|_{\rm count}\). The upper bound follows from the physical row bound at grade zero and the weight \(\ell^{-2}\) in the fine \(L^2\) norm. The form equivalence carries over by a small relative perturbation of \(\Delta:V_2\to L^2\) and \(q:V_2\to\ell^2_{\rm count}\), including their zero-extension restrictions.

For this paragraph use the stronger coefficient hypothesis stated at the start of the section. A fourth-order source has a representation
\begin{equation}\label{mb:eq-fourth-source}
                 \rho=\Delta f+D^Tg+h
\end{equation}
with \(f\) in grade-two \(C^\alpha\), \(g\) of bounded grade three and \(h\) of bounded grade four. The norm is the infimum over representations of the sum of these normalized norms on the appropriate complete boxes. Its local energy-dual bound follows by testing \(\Delta f,D^Tg,h\) against \(\Delta u,Du,u\) in the physical pairing. The normalized cell \(L^2\) sizes of \(f,g,h\) are \(O(1),O(\xi_y^{-1}),O(\xi_y^{-2})\). This source class is preserved by smooth localization. For instance
\[
 [\Delta,\chi]v=D^T\big((D_{\rm flat}\chi)(v_-+\operatorname{Ad}(U_b)v_+)\big)
                               -(\Delta_{\rm flat}\chi)v
\]
on the actual graph, with the same present scalar stencils in \(\Delta_{\rm flat}\). At the tail the divergence term supplies \((\chi_--\chi_+)(v_-+\operatorname{Ad}(U_b)v_+)/\eta^2\); its \(v_-\) term is canceled by \((\Delta_{\rm flat}\chi)v\). The head gives the transported opposite contribution. This proves the identity, including for crossing edges on the zero extension. Thus commuting \(\chi\) across \(\Delta f\) differentiates only \(\chi\) in the new divergence and bounded point sources. If the reference connection is changed within the tube, expand the difference of its two Laplacians as first-order shifted terms and a zeroth term. The first-order terms have the divergence representation from the preceding proof. Their coefficients and first differences are \(O(\varepsilon/\ell)\) and \(O(\varepsilon/\ell^2)\), respectively. Consequently the norms for \eqref{mb:eq-fourth-source} at the varying and fixed bases are equivalent without differentiating \(f\) or \(g\).

\begin{lemma}[Strong restoration family]
\label{mb:lem-strong-restoration-family}
On a common weighted \(C^{1,\alpha}\) tube, \(G_sR\) has an exact connected family from grade-two \(C^\alpha\) to the \(Y^{2,\alpha}\) space of \eqref{mb:site-solution-norm}. Accordingly \(\Lambda\) has such a family from \(\mathcal A_\alpha\) to \(Y^{2,\alpha}\) and \(S\) from \(\mathcal A_\alpha\) to \(\mathcal A_\alpha\). All norms are in \eqref{mb:eq-family-norm} with fixed \(s,\vartheta\) after the separation and tube choices.
\end{lemma}

\begin{proof}
For a positive compression of \(A_4\), begin with its exponential \(V_2\) energy column. Here the weight needs two differences. Smooth the truncated distance to the source cell by the same squared local partition as in the positive walk, with a fixed radius multiple of \(\ell\). Over the profiles meeting \(y\), the center distances differ by a bounded amount. Since their weights sum to one, subtract one center value before differentiating the weighted sum. The resulting \(\psi\) differs from the truncated distance by \(O(1)\) and obeys \(|D\psi|\le C\ell^{-1}\) and \(|D^2\psi|\le C\ell^{-2}\). The second-difference product rule gives
\[
 \|e^{-s\psi}\Delta(e^{s\psi}u)-\Delta u\|_2
     \le C|s|(\|\ell^{-1}Du\|_2+\|\ell^{-2}u\|_2).
\]
The full-row relative weight oscillation is \(O(|s|)\), so the conjugated \(A_4\) form differs by \(O(|s|)\) in \(V_2\to V_2^*\). Its positivity or small complex perturbation therefore gives the exponential energy column from the source in \eqref{mb:eq-fourth-source}. Both norms are dimensionless in four dimensions.

If \(u\) is that local solution, put \(v=\Delta u-f\). On the interior cutoff stock the two equations are exactly
\[
           \Delta v=D^Tg+h-q^*qu,\qquad \Delta u=f+v .
\]
Here \(q^*qu\) is bounded of grade four by the actual complete row and the local energy bound. The first equation and a squared cutoff test bound \(v\) in local \(H^1\) by its local \(L^2\) norm, \(g\) and the bounded source, using the real reference in a small complex tube. The divergence bootstrap of the preceding proof then gives \(v\) in \(W^{1,p}\) for any fixed finite \(p\). Take \(p>4/(1-\alpha)\) to obtain a grade-two \(C^\alpha\) bound. Apply the scalar point-force bootstrap and the covariant Schauder estimate to the second equation, in the scaled \(C^{1,\alpha}\) reference class. This gives \(u\) in \(Y^{2,\alpha}\). The discontinuous averaging density appeared only in the first equation. Nested interior cutoffs retain the exponential cell estimate, and at the true bare scale the energy column alone gives the rescaled finite differences. No output norm is used across the artificial compression boundary.

The exact fourth-order cutoff identity is
\[
 [\Delta^2,\chi]u=\Delta[\Delta,\chi]u+[\Delta,\chi]\Delta u .
\]
For strong \(u\), the first term is \(\Delta\) of a grade-two \(C^\alpha\) function bounded by \(C/m\) times the strong norm. Use the ordinary two-difference product rule and the first three cutoff differences, so that both its point and H\"older bounds have this factor. In the second term apply the displayed divergence form for \([\Delta,\chi]\) to the bounded \(\Delta u\), using only its values. This gives a grade-three divergence coefficient and a grade-four bounded term, again with factor \(C/m\). The complete-row commutator of \(q^*q\) gives a bounded grade-four term with the same factor. At a true bare boundary the atomic formulas use the original finite matrix and each coefficient has its appropriate bare power. Thus the slow positive parametrix is a contraction on the source class \eqref{mb:eq-fourth-source}. Its exact inverse has the strong local family \(G_4:\rho\mapsto Y^{2,\alpha}\).

For the short-power real extensions of \(M_4\) only its energy
family is needed. The same fourth-order compression and cutoff
expansion just used yields that family on the real curvature class.
In fact the two-difference product rule bounds
\(\|[\Delta,\chi]u\|_2\le (C/m)\|u\|_{V_2}\).
In the first term of \([\Delta^2,\chi]\) test the outer
\(\Delta\) against the test function. In the second move the
skew-adjoint commutator to the test function. This gives
\(\|[\Delta^2,\chi]\|_{V_2\to V_2^*}\le C/m\).
The complete-row cutoff difference gives the same bound for
\([q^*q,\chi]\). Use the preceding exponentially weighted
energy columns and the same center sums to obtain the convergent
walk on \(V_2^*\to V_2\), increasing \(m\) if needed for both
norms. These are the same filed matrix coefficients. The stronger
regularity hypothesis is needed on the local collars where their
strong output norms are estimated. The rest of a real completion
for the short-power spectral bound needs only the energy estimates.

The counting coarse matrix \(M_4=qG_4q^*\) is positive with a uniform gap. For a counting datum \(\beta\), dualize against \(E_s\beta\), use \(qE_s\beta=\beta\) and \(\|E_s\beta\|_{V_2}\le C\|\beta\|_{\rm count}\). This gives the lower bound by Cauchy--Schwarz in the real \(A_4\) form. The upper bound follows from the row trace. Its actual family comes from \(G_4\) and the local row maps, with \(q^*\) giving a bounded source of grade four and \(q\) tracing a grade-zero site output. Apply \cref{mb:lem-coarse-carrier-inverse}. The exact identity is
\begin{equation}\label{mb:eq-strong-restoration-identity}
        G_sR=(G_4-G_4q^*M_4^{-1}qG_4)\Delta .
\end{equation}
At real coefficients the right side applied to \(f\) is the unique zero-average minimizer of \(\|\Delta u-f\|_2^2\). Its normal equation gives \(\Delta u=Rf\) by \(\Delta(\ker q)=\operatorname{ran}R\). This is exactly \(G_sRf\). The identity continues bilinearly. Composition of the local families in \eqref{mb:eq-strong-restoration-identity} uses \(f\mapsto\Delta f\) as a source in \eqref{mb:eq-fourth-source}, without differentiating the input norm. Finally \(D^T\) takes \(\mathcal A_\alpha\) to grade-two \(C^\alpha\) on complete boxes, and \(D\) takes \(Y^{2,\alpha}\) to \(\mathcal A_\alpha\). They are finite stencils with the stated grade bounds. This proves the two restoration families and their complete identities.
\end{proof}

\subsection{The two strong named propagators}

\begin{theorem}[Strong current and row response families]
\label{mb:thm-strong-families}
Fix \(0<\alpha<1\), finite \(p>4/(1-\alpha)\) and a finite carrier rate. Work on the weighted \(C^1\) coefficient class of \cref{mb:lem-strong-positive-columns}. After the fixed geometry and tube choices, and a sufficiently small touching current for these choices, \(G_0,G_J,G_{1,J}\) have exact connected families \(\mathcal F_p\to\mathcal A_\alpha\). They also have families from bounded field divergence sources \(\nabla_U^T F\), with \(F\) of grade two, to the point field norm of grade one. If \(F\) has the scaled grade-two \(C^\alpha\) norm, the output is \(\mathcal A_\alpha\). In particular one can use \(p=8,\ \alpha=1/4\) for the first-order regular-pair families.

For either current variant, the coupled response and lift have the same strong families with normalized datum \(\Xi_f b\), normalized multiplier \(\Xi_f^3\mu\), and field force norm \(\mathcal F_p\). Thus for sources and data in a single cell or row, each of these normalized output norms has the exponential factor of its distance to the source. The two current readouts and the mass-removed coarse operator retain the families of \cref{mb:thm-current-families}. The constrained Green operators are obtained by their stated algebraic subtractions, including the \(\mathfrak G\) gauge term.
\end{theorem}

\begin{proof}
Use the exact touching-current split before expanding coefficients, as in \cref{mb:thm-current-families}. First use the two scalar families of \cref{mb:lem-strong-positive-columns} and the normalized site Gram inverse. The point column \(\widehat F\) has site amplitude of grade two and \(D\widehat F\) has amplitude of grade three. When \(\widehat F^T\) pairs a grade-two field with one such column, the physical volume contributes \(\xi_y^4\), canceling the two grade-two factors. Hölder's inequality gives the same cancellation for the normalized \(\mathcal V_p\) input. Consequently \(R=I-\widehat F\widehat M^{-1}\widehat F^T\) is bounded as a family on \(\mathcal V_p\) and on the grade-two site space with bounded first difference of grade three. The exact adjoints then satisfy
\begin{equation}\label{mb:eq-strong-adjoint-slots}
 \begin{aligned}
 \Lambda^T&=DRG_s:
   \{\text{bounded site grade four}\}\to
   \{\text{bounded field grade three}\},\\
 S^T&=I-DRG_sD^T:\mathcal F_p\to\mathcal F_p .
 \end{aligned}
\end{equation}
The second map uses the finite-\(p\) divergence estimate. In particular \(\Lambda^TD^T=I-S^T\) applies to a bounded field force through its finite-\(p\) norm.

The positive vector family and \eqref{mb:eq-zero-current-woodbury} give \(G_0:\mathcal F_p\to\mathcal A_\alpha\). Indeed \(D\widehat F\) gives bounded field grade-three forcing, and its transpose pairs that grade with grade-one fields on physical cells, so the volume factors again cancel. The intermediate \(\widehat\Theta^{-1}\) acts in normalized counting coordinates. For \(u\in\mathcal A_\alpha\) put \(\lambda=\Lambda u\). Its input \(D^T u\) is bounded of grade two, and the smoothed columns of \(P\) show that \(R\) preserves this slot. Apply the bounded site point-force family to \(R D^T u\), changing the normalized input from grade four to grade two. The extra factor \((\xi_z/\xi_y)^2\) for the local column is absorbed by the scale reserve. This gives \(\lambda\) of bounded grade zero and \(D\lambda\) of bounded grade one. These are the site outputs used in the following current form:
\[
 B[J]u=W[J^{\rm t}]\lambda+
 \Lambda^T\bigl(W[J^{\rm t}]^Tu-\tfrac12W[J^{\rm t}]^TD\lambda\bigr)
              -\tfrac12\Lambda^TD^TW[J^{\rm t}]\lambda.
\]
The first term is bounded field grade three. In the last term \(W[J^{\rm t}]\lambda\) has that same grade and enters the finite-\(p\) divergence slot of \eqref{mb:eq-strong-adjoint-slots}. Each coefficient reads its current bond and the complete background at all ends. The result is a family from \(\mathcal A_\alpha\) to \(\mathcal F_p\) of norm at most \(C_{p,\alpha,\vartheta}\gamma_{\rm t}\). Thus the first Neumann series of \cref{mb:thm-current-families} converges in these strong slots.

For the second variant, the first lift gives the same rooted multiplier family as before, since its strong estimate includes the cell pairing required in \cref{mb:cor-current-readouts}. Both absolute normalized sums of the two-leg Hessian are bounded by \eqref{mb:eq-two-leg-energy} and its pointwise row proof. In particular \(C_\omega\) takes a bounded grade-one field to a bounded grade-three force. Its two legs and the multiplier retain their joined data carrier. The right factor \(Su=u-D\lambda\) is bounded of grade one by the preceding site estimate. On the left use \(S^T\) in the finite-\(p\) slot \eqref{mb:eq-strong-adjoint-slots}. Thus \(S^TC_\omega S\) has norm \(C_{p,\alpha,\vartheta}\gamma_{\rm t}\) from \(\mathcal A_\alpha\) to \(\mathcal F_p\). The second Neumann series converges after the corresponding fixed decrease of the ceiling. The stronger \(C^{1,\alpha}\) domain in \cref{mb:lem-strong-restoration-family} is used for the separate second-derivative site conclusion.

Here are the divergence input slots. Use the normalized field space \(X_{1,p,y}=\xi_y^{1-4/p}\|u\|_p+\xi_y^{2-4/p}\|\nabla_Uu\|_p\). For \(G_F\nabla_U^TF\), the same variational energy start and local divergence bootstrap as for \(G_s\) give \(X_{1,p}\) with \(F\) of grade-two \(L^p\). The equation is componentwise by \cref{mb:lem-wilson-component-stencil}, and the additional row force is pointwise bounded by the energy. The slow commutator of \(G_F\) now maps \(X_{1,p}\) to \(\mathcal F_p\); only the rightmost factor of the positive walk needs the divergence source estimate. A grade-two \(C^\alpha\) source \(F\) gives an \(\mathcal A_\alpha\) output at that rightmost factor and in the no-error term, by the interior divergence Schauder estimate. In that estimate put the connection's lower terms and the finite curvature kernel in finite \(L^p\) forcing, using the \(W^{1,p}\) bound just obtained. The rest is ordinary divergence of a \(C^\alpha\) field. The corresponding statement at the bare boundary follows by finite summation on the atomic stencil.

For the current corrections on \(X_{1,p}\), the actual \(D^T u\) is grade-two \(L^p\), and \(R\) is bounded on that site slot by the same column pairing. The positive site point-forcing bootstrap gives two \(L^p\) differences of \(\Lambda u=G_sRD^T u\). This is the site estimate proved above with a grade-two \(L^p\) point source rather than grade four, and with output grade zero. The extra scale ratio between the two normalizations is paid once in the spatial reserve. Morrey gives bounded \(\lambda\) and \(D\lambda\). Therefore the displayed \(B[J]\) formula and the \(C_\omega\) bound hold also from \(X_{1,p}\) to \(\mathcal F_p\). For \(S^T C_\omega S\) use the pointwise grade-one output of \(Su\) and the same finite-\(p\) adjoint. This gives the divergence-input statements by using the divergence source at the rightmost inverse factor. All subsequent loops use \(\mathcal F_p\to\mathcal A_\alpha\). The row trace and transpose have bounded strong cell norms by their complete row moments. The normalized row inverse and \eqref{mb:eq-field-saddle} therefore give the strong lift and saddle.

For \(G_{1,J}DRD^TG_{1,J}\), the first \(G_{1,J}\) takes a point force to \(\mathcal A_\alpha\). Its \(D^T\) output is grade-two \(C^\alpha\) because the comparison connection is Lipschitz. The smoothed columns of \(P\) have bounded first differences, so \(R\) preserves this H\"older slot. The remaining \(D\) is a field divergence of this site datum. Explicitly, for \(T_\mu\) as in \cref{mb:lem-wilson-component-stencil}, \((I-T_\mu)z=(T_\mu^{-1}-I)(T_\mu z)\). Thus the tensor with component \(\delta_{\mu\nu}T_\mu z\) gives \(Dz\) under the declared forward-difference transpose on full stencils. At a true bare boundary retain the finite remainder \(Dz-\nabla_U^TF\) on its actual stencil. It is a bounded grade-three point force, so the ordinary bare force estimate controls it. The \(G_{1,J}\) divergence estimate gives the desired output. With a bounded divergence source at the rightmost factor, use \(X_{1,p}\) there and the grade-two \(L^p\) slot in the middle. This proves the two corresponding estimates for \(\mathfrak G\). The same column pairings give the pointwise force family \(DPD^T:\mathcal A_\alpha\to\mathcal F_\infty\) and the raw lift transpose bound.
\end{proof}

\mbneedspace{12\baselineskip}
\begin{corollary}[Output bounds for the compressed inverses]
\label{mb:cor-compressed-regularity}
Use the actual or the full single-comb rows and the ambient
coefficient stock of \cref{mb:prop-compressed-gap}. Assume
the regular local frame bounds of \cref{mb:prop-local-scalar}.
Let \(G_{\rm c}\) be the inverse of the compressed \(A_s\), or the
inverse of the compressed zero-current \(K_0\). In the former case
\(\nabla\) denotes the site difference and \(L_{\rm c}\) the
scalar site Laplacian on the prescribed graph or Dirichlet
compression. In the latter case \(\nabla\) denotes the
component difference and \(L_{\rm c}\) the compressed
component covariant Laplacian.

For bounded \(f\) supported in the cell \(B_z\), \(u=G_{\rm c}f\)
satisfies
\begin{align}
 &\xi_y\|u\|_{\infty,B_y}+\xi_y^2\|\nabla u\|_{\infty,B_y}
   +\xi_y^{2+\alpha}[\nabla u]_{\alpha,y}
   +\xi_y^3\|L_{\rm c}u\|_{\infty,B_y}\notag\\
 &\qquad\le C_\alpha e^{-s d(y,z)}\xi_z^3\|f\|_\infty
 \label{mb:eq-compressed-point}
\end{align}
for \(0<\alpha<1\). For \(u=G_{\rm c}\nabla^T F\) with \(F\)
supported in the complete source collar of \(B_z\), one has
\(\xi_y\|u\|_{\infty,B_y}\le C e^{-s d(y,z)}\xi_z^2\|F\|_\infty\).
If \(F\) is \(C^\alpha\) on that full stock, including its
support boundary, the first three terms of
\eqref{mb:eq-compressed-point} are bounded by
\[
 C_\alpha e^{-s d(y,z)}
       \bigl(\xi_z^2\|F\|_\infty+\xi_z^{2+\alpha}[F]_{\alpha,z}\bigr).
\]
Use complete interior stencils and the prescribed bare
boundary convention. These bounds continue to a fixed smaller
weighted \(C^1\) tube on the same ambient stock, including
the crossing plaquettes, around the real reference, with norms
measured in its comparison frames.

The normalized compressed row matrix
$a_f^{1/2}QG_{0,\rm c}Q^Ta_f^{1/2}$ has a uniform positive
real gap. Its inverse and the exact lift
\[
 H_{\rm c}=G_{0,\rm c}Q^T(QG_{0,\rm c}Q^T)^{-1}
\]
have exponential local bounds, with holomorphic continuation
on the same smaller tube. For a datum supported on a row
$c$, the value, first difference and gradient H\"older terms
of $H_{\rm c}b$ in the normalized output above are bounded
by $C_\alpha e^{-s d(y,c)}|\xi_c b(c)|$.
\end{corollary}

\begin{proof}
The weighted energy form estimate for either compressed
positive inverse gives a cell energy bound of order
\(e^{-s_0d(y,z)}\xi_z^3\|f\|_\infty\). In fact the source has
energy-dual norm at most \(C\xi_z\|f\|_{2,B_z}\). For a
divergence source the pairing moves its difference onto the
test field, so its energy-dual norm is bounded instead by
\(C\|F\|_{2,B_z^+}\le C\xi_z^2\|F\|_\infty\).
This uses the full prescribed stencil in the source collar.
The crossing plaquette forms are local and occur only in
the bare stock. The extra scalar Dirichlet diagonal commutes
with the spatial weights. For the projected field form use
the same site-column bound for \(R\) as in the real gap proof.
The conjugated perturbation estimates are uniform for the
real and complex coefficient classes in the statement.

On complete interior boxes the site equation is
\[
 L_{\rm c}u=f-q^Ta_squ .
\]
For the field equation the component stencil gives
\[
 L_{\rm c}u=f+DPD^Tu-Q^Ta_fQu-Z_Uu .
\]
Both row masses have the stated bounded point forcing from
the energy and the full normalized row densities.
For \(DPD^Tu\) use \(D\widehat F\widehat M^{-1}
(D\widehat F)^Tu\). Its point columns \(D\widehat F\) follow
from the site bootstrap with the bounded injection
\(q^T\Xi_s^{-4}e_z\), with the same scaling and complete
rows as in \cref{mb:lem-strong-positive-columns}. Their
transpose pairs with the local field energy by the physical
cell volume, giving a point force of grade three. The finite
curvature shifts in \(Z_U\) are bounded on the local \(L^p\)
spaces. In unit cell coordinates the site and vector local
bootstraps therefore give the first three point-output bounds.
Inserting these bounds in the two equations gives the
Laplacian bound. At a true bare boundary \(\xi_y=\eta\) and
the same conclusions follow from the local energy and the
finitely many stencil entries in those units.

For divergence input, retain the divergence in the local
equations and apply the finite-\(p\) first-order bootstrap.
This gives the point field for \(p>4\), including with the
bounded projection and row terms just estimated. If \(F\)
is H\"older, put the other lower-order terms in finite \(L^p\)
on a smaller cutoff stock. The flat resolvent then bounds
their gradient H\"older norm by \eqref{mb:eq-gradient-morrey},
and the divergence term uses the two-difference multiplier
of \cref{mb:prop-flat-schauder}. The difference between the
flat and covariant divergences is bounded multiplication by
the connection potential. This proves the claimed derivative
range and uses only the fixed scaled first-difference bounds
on the complex coefficients. The bare divergence estimate
again uses finite stencils. Applying these local arguments
to the weighted columns costs only fixed collar widths in
the metric. If the site column is read in its grade-two
output and grade-four source conventions, conversion to the
scales in \eqref{mb:eq-compressed-point} costs
\(\xi_z/\xi_y\); reserve this single ratio using
\eqref{mb:rate-reserve}. The normalized single-comb rows
have the same finite moment and perturbation bounds by
\cref{mb:prop-compressed-gap}, so all these computations
apply to that full differential as well.

For the row matrix, dualize against the uniformly bounded exact row
lift in the compressed energy form, as in
Theorem~\ref{mb:thm-coarse-schur}. This proves the normalized
positive gap. The just-proved weighted $G_{0,\rm c}$ bound
and the complete row bounds give an $O(s)$ change of the
normalized row matrix under exponential conjugation.
Its Neumann inverse therefore has exponential decay.
The physical row-transpose injection of one normalized
datum has bounded grade-three point density. Compose its
compressed point-source estimate with that row inverse
to obtain the stated lift bounds. The same relative
perturbation proves their complex continuation.
\end{proof}

\paragraph{Bounds independent of localization choices.}
The ordinary point and gradient estimates for the two inverse operators can use the current ceiling of the weighted energy theorem before any walk tile is chosen. For a cell source, the weighted energy inverse bounds the field in \(X_y\), and \eqref{mb:eq-weighted-restoration} bounds its restored site parameter in \(\mathcal Z_{s,y}\), with the same exponential localization after a scale reserve. On unit-scale boxes, strict Sobolev embedding gives the field and \(D\lambda\) in \(L^3\) and \(\lambda\) in \(L^{12}\). The first current correction is consequently in field \(L^8\): \(W[J^{\rm t}]\lambda\) is in \(L^{12}\), the two site point sources \(W[J^{\rm t}]^Tu\) and \(W[J^{\rm t}]^TD\lambda\) are in \(L^3\), and the last source is a divergence with coefficient in \(L^{12}\). A site \(L^3\) point source gives a derivative in \(L^8\) by two finite-\(p\) differences and strict embedding. The site divergence source gives the same derivative bound by the variational estimate. Projection by \(R\) preserves these bounds by its smoothed columns. The second correction is initially in field \(L^3\) by both absolute sums of its row Hessian and the finite-\(p\) bound of \(S^T\).

In the actual field equation, \(DPD^T u\) is a bounded point force by the point columns of \(D\widehat F\) paired with the cell energy. The other terms are the two current corrections just bounded and the given force. Thus the local \(W^{2,3}\) vector estimate gives bounded field and gradient in \(L^8\). The site equation for \(\Lambda u\) now gives bounded \(\lambda,D\lambda\), since its right side \(R D^T u\) is in \(L^8\) in scale units. Both current corrections consequently belong to every required finite \(L^p\), by \eqref{mb:eq-strong-adjoint-slots} and the two-leg point bound. The final \(W^{2,p}\) estimate gives the asserted point and gradient H\"older bounds for \(G_J,G_{1,J}\) on \(\mathcal F_p\). Each estimate is local on fixed boxes and retains the exponential cell weight. For a divergence input, keep it in variational form throughout. The first vector bootstrap gives \(W^{1,3}\), then \(W^{1,8}\) and bounded field, from the same initial current bounds. This yields the same final finite-\(p\) derivative bound for bounded divergence data, and the divergence Schauder step for H\"older data. The ordinary point estimate for the component Laplacian at zero current follows directly from its equation with the bounded row, curvature and \(DPD^T\) sources. At a nonzero independent current these arguments give the finite-\(p\) forcing bound instead. They require no derivative of that current.


% Construction source: assembly/physical-gap.tex
\section{Removing the row mass in the outer precision}
\label{mb:sec-physical-gap}

Use the separated inner hierarchy in top units. Retain a
union of complete \(L\)-blocks of the next lattice, separated
from \(\Omega_k^c\) by the fixed required collar. Let \(X_{\rm out}\)
be the level-\(k\) bonds touching those blocks, and let
\(\iota\) extend fields on \(X_{\rm out}\) by zero to the
independent determining row space. All incident double
blocks, outer pivots, contours and curl plaquettes used
below lie in the interior top stock, with their complete
fine backward carriers in \(\Omega_k\). These conditions
hold for the proper fluctuation geometry with sufficiently
large fixed separation. The outer pairing is counting
measure, since the supporting top rows have length one.
Inner row pairings remain physical.

Form the ordered current inverses and lifts of
Theorem~\ref{mb:thm-coarse-schur}. Write
\[
 \Sigma_1=(QG_{1,J}Q^{\mathsf T})^{-1},\qquad
 \omega_0=H_J^{\mathsf T}J,\qquad
 \omega_1=H_{1,J}^{\mathsf T}J,\qquad
 a_{\rm top}=\iota^{\mathsf T}a_f\iota .
\]
The first multiplier enters \(K_{1,J}\). The second
will enter the outer curvature term in the proper
extension. On physical stationary pairs they agree,
as proved below.

\begin{lemma}[The top curl after row-mass subtraction]
\label{mb:lem-physical-schur}
At a real reference in the curvature-small class with
a sufficiently small admitted independent current, put
\(V=\mathcal M^kU\) on the complete outer stock. For
\(a\) supported on \(X_{\rm out}\), extend it by zero
to the curl plaquettes \(\mathcal P\) of
Lemma~\ref{mb:lem-outer-geometry}. Then
\begin{equation}\label{mb:eq-physical-curl}
 \begin{split}
 \langle a,(\iota^{\mathsf T}\Sigma_1\iota-a_{\rm top})a\rangle
    &\ge c\|\mathrm d_V a\|_{2,\mathcal P}^2
                          -C(\kappa+\gamma)\|a\|_2^2,\\
 \|\iota^{\mathsf T}\Sigma_1\iota-a_{\rm top}\|_{2\to2}&\le C .
 \end{split}
\end{equation}
The constants are independent of the depth, the proper
inner geometry and the admitted volume.
\end{lemma}

\begin{proof}
The required flux estimate is
\begin{equation}\label{mb:eq-covariant-flux}
 \|\mathrm d_V Qv\|_{2,\mathcal P}
                       \le C\|\mathrm d_Uv\|_2+C\kappa\|v\|_X ,
\end{equation}
where the left side uses the four actual top rows of each
retained plaquette. To prove it, first use the flat
coefficient formula on the top grid of spacing one and
put \(N=\eta^{-1}\). Write \(A_\eta\) for the average over
the centred \(N^4\) sites and
\(S_\mu=\eta\sum_{j=0}^{N-1}T_\mu^j\) for the segment mean
with ordinary fine shift \(T_\mu\). Then \(Q_\mu=A_\eta S_\mu\)
sampled at the unit-grid origins. The exact identity
\(T_\mu^N-I=S_\mu(T_\mu-I)/\eta\) gives
\[
  (\mathrm d_c Qv)_{\mu\nu}
       = A_\eta S_\mu S_\nu(\mathrm d v)_{\mu\nu}.
\]
The nonnegative flux row has row sum one and counting
column sum \(\eta^4\), also as an upper bound when rows
are restricted. In fact the translates of the blocks give
that column sum for \(A_\eta\), and the two segment means
preserve it. Thus the flux map is a contraction from the
physical fine \(L^2\) norm to the unit-grid \(L^2\) norm.

For a real field choose a common comb frame on the complete
fixed-size physical cube about each top plaquette, containing
all four fine carriers. All assigned scales in this stock
are one. The fine potential is \(O_L(\kappa)\), and the
coarse bonds of \(V\) have angles \(O_L(\kappa)\) in the
same frame, by the summable raw transport comparison.
The actual field row differs there from the flat row by
\(C_L\kappa\eta^4\) per coefficient. Its physical \(L^2\)
functional norm difference is \(O_L(\kappa)\). The fine
curl differs from the flat curl by \(O_L(\kappa)\) as a
map on local fine \(L^2\), because its transport difference
is \(O_L(\eta\kappa)\) before division by \(\eta\).
The outer curl is a unit-spacing difference and costs a
fixed constant on row errors. Insert these estimates in
the flat identity. Summing their squares uses the bounded
overlap of the complete cubes, and yields
\eqref{mb:eq-covariant-flux}. Unitary changes of frame
preserve these local estimates.

Take \(v=H_{1,J}\iota a\). Then \(Qv=\iota a\),
\(\|v\|_X\le C\|a\|_2\), and
\(K_{1,J}v=Q^{\mathsf T}\Sigma_1\iota a\).
Subtracting the row-mass form on this field gives exactly
\[
 \langle a,(\iota^{\mathsf T}\Sigma_1\iota-a_{\rm top})a\rangle
 =\langle v,(H_W+D_URD_U^{\mathsf T}-B[J]
                       -S^{\mathsf T}C_{\omega_0}S)v\rangle_f .
\]
The gauge term is nonnegative. The current form, the
second-jet contraction and the weighted restoration are
bounded by \(C\gamma\|v\|_X^2\), using
Proposition~\ref{mb:prop-current-completion},
Corollary~\ref{mb:cor-current-readouts} and
\eqref{mb:eq-constraint-small}. The Wilson bound
\eqref{mb:eq-wilson-form} gives a positive curl square
up to \(C(\kappa+\gamma)\|v\|_X^2\).
By \eqref{mb:eq-covariant-flux},
\(\|\mathrm d_Uv\|_2^2\ge
 c\|\mathrm d_Va\|_{2,\mathcal P}^2-C\kappa^2\|v\|_X^2\).
This proves the lower estimate. The normalized row inverse
and unit top lengths give the operator upper bound.
\end{proof}

Let \(q_V\) be the actual raw one-step plain constraint
logarithm on the complete outer stock, and let \(\Phi_V\)
be its pivot graph, with derivative \(C_V\) and free
restriction the identity. Its second derivative is
\[
 D^2\Phi_V(0)[u,w]
       =-h_VD^2q_V(0)[C_Vu,C_Vw],
\]
with \(h_V\) the pivot-supported right inverse. For a subset
\(Y_{\rm out}\) of the retained free coordinates, \(e_Y\)
extends by zero on the other free coordinates. We retain
all incident pivots when applying \(C_Ve_Y\), even if
they lie outside the selected free subset. A row that
reads an internal link of a retained block is incident
to that block. A crossing link is read only by its
longitudinal face row, since internal contours use no
crossing links. Thus zero extension creates no pivot
outside the declared incident stock.

\begin{theorem}[The proper outer fluctuation gap]
\label{mb:thm-physical-outer-gap}
With the preceding proper stock and the hypotheses of
Lemma~\ref{mb:lem-physical-schur}, define the symmetric
outer matrix by
\begin{equation}\label{mb:eq-physical-outer-form}
 \begin{split}
 \langle u,A_Y(\tau)w\rangle
 ={}&\tau\langle C_Ve_Yu,
       (\iota^{\mathsf T}\Sigma_1\iota-a_{\rm top})C_Ve_Yw\rangle\\
 &+\tau\langle\iota^{\mathsf T}\omega_1,
                              D^2\Phi_V(0)[e_Yu,e_Yw]\rangle\\
 &+\langle b_VC_Ve_Yu,b_VC_Ve_Yw\rangle .
 \end{split}
\end{equation}
For sufficiently small fixed \(\kappa,\gamma\) and real
\(\tau\) in a fixed neighbourhood of one, this matrix
has bounds \(c_LI\le A_Y(\tau)\le C_LI\).
On the subspace with zero reconstructed field on each
retained internal tree, the same estimates hold without
the last term. Its inverse is holomorphic and uniformly
bounded on a fixed smaller weighted \(C^1\) complex
background and independent-current tube, and for complex
\(\tau\) near one. The bounds are independent of depth,
the proper inner region, the free subset, and compatible
volume.
\end{theorem}

\begin{proof}
The pivot graph has uniformly bounded first derivative,
with \(\|C_Ve_Yu\|_2\ge\|u\|_2\) by its free restriction.
Its fixed one-step carriers and second-derivative bounds
give
\[
 |\langle\iota^{\mathsf T}\omega_1,
          D^2\Phi_V(0)[e_Yu,e_Yw]\rangle|
                  \le C_L\gamma\|u\|_2\|w\|_2
\]
by the second readout bound and bounded overlap.
Combine \eqref{mb:eq-physical-curl} with the geometric
inequality \eqref{mb:eq-outer-geometry} on \(C_Ve_Yu\).
Its row term vanishes. The two \(O_L(\kappa+\gamma)\)
errors are absorbed in the curl-contour gap.
For tree-zero fields use the version of that inequality
without contours. Tree links are internal and are never
pivots, so the tree-zero subspace is fixed in the free
coordinates. The upper bound follows from the operator
bounds of all factors. An increment of \(\tau\) costs
\(C_L|\tau-1|\) in norm, so the bounds persist.

For the complex domain, the ordered field inverses,
normalized coarse Schur matrices, and their current
readouts are holomorphic on the common admitted tube.
On a fixed smaller tube their changes are \(O(\varepsilon)\)
in their energy, normalized row and weighted current
norms when the joint coefficient increment is
\(\varepsilon\). This follows by the inverse identity
and Cauchy's formula applied to the bounded normalized
operators and readouts, respectively. Restricting to
unit top rows turns these into the stated outer norms.
The actual hierarchical coarse field varies by
\(O_L(\varepsilon)\) in raw unit-spacing coefficients,
by the summable averaging and node-frame bounds.
The outer graph and contour derivative have fixed
one-step analytic carriers, so the matrix difference
is at most \(C_L(\varepsilon+|\tau-1|)\), uniformly in
\(Y_{\rm out}\). A smaller tube gives a norm difference
less than half the real gap, proving the claim by
a Neumann series.
\end{proof}

\subsection{Identification on the physical chart}

At a given physical critical pair, assume \(J=j_U\) is the
actual inner Wilson current and \(j_U=Q^{\mathsf T}\omega\).
If the field is critical on the fixed Landau constraint
slice, the full stationarity identity follows because
\(QS=Q\) and the Wilson Ward identity kills the gauge
part \(I-S\). The row lift proves uniqueness of \(\omega\).
Both \(QH_J=I\) and \(QH_{1,J}=I\) then give
\(\omega_0=\omega_1=\omega\). The physical Hessian
\(H_L=H_W-C_\omega\) uses the same constraint jet.
The physical Ward identity for \(H_L\) and the second
lift's horizontality give
\begin{equation}\label{mb:eq-physical-lift-readout}
 H_{1,J}^{\mathsf T}H_LH_{1,J}=\Sigma_1-a_f .
\end{equation}
Indeed \(H_{1,J}^{\mathsf T}K_{1,J}H_{1,J}=\Sigma_1\)
by its defining equation. The gauge term vanishes,
the row penalty contributes \(a_f\), and the remaining
form on that slice is \(H_L\).

Let a twice differentiable critical branch for this fixed
fine action have coarse parameter $B$ and reduced datum
$b(B)$. Use the relative exponential coordinates at the
given critical pair $U$, and freeze its reference slice as
$\ker(R_UD_U^{\mathsf T})$. The restriction of $Q$ to this
slice is onto by the restored row lift. Near zero, its
fixed right inverse continues by a finite-dimensional
Neumann inverse for $D\mathfrak q(a)$ on the same slice.
The restricted stationarity multiplier is therefore unique
and continuously differentiable along the given branch.
Differentiating this stationarity equation yields the
constrained Lagrangian saddle with row value $Db(0)$.
By uniqueness of that saddle its tangent is
\[
                         H_{1,J}Db(0).
\]
The second chain rule then gives the effective action
Hessian
\begin{equation}\label{mb:eq-coarse-hessian-identification}
 D^2\mathcal A_{\rm eff}(0)
  =Db(0)^{\mathsf T}(\Sigma_1-a_f)Db(0)
                           +\langle\omega,D^2b(0)\rangle_r .
\end{equation}
In fact differentiating the datum constraint twice
replaces the fine second-variation term paired with
\(j_U\) by \(\langle\omega,D^2 b\rangle_r\) minus the
fine constraint Hessian paired with \(\omega\).
That subtraction is precisely \(C_\omega\).

On a complete interior top stock with all lower-layer
data fixed, the specified relative left coarse coordinates
give \(Db(0)=\iota\) and \(D^2b(0)=0\) on the variations
under consideration. The interior equation uses the same
left increment \(e^{iB(c)}\) relative to \(V(c)\).
An outward equation can have an extra endpoint gauge
factor from the lower-layer data. It is constant when
those data are fixed, whereas an absolute logarithm
would still have a second jet at nontrivial \(V(c)\).
Equation~\ref{mb:eq-coarse-hessian-identification}
retains the extra term in such coordinates.

The outer chain rule on the actual pivot graph adds
\(\langle\iota^{\mathsf T}\omega,D^2\Phi_V(0)\rangle\).
Together with the relative contour Hessian this identifies
$A_Y(1)$ in \eqref{mb:eq-physical-outer-form}
with the physical action-plus-contour Hessian.
For $A_Y(\tau)$ the factor $\tau$ multiplies the
action part, including its multiplier graph term; the
contour gauge weight stays fixed. When the reference contours are
one, the relative and absolute contour Hessians agree.
In the homogeneous geometry this gives the effective curl
bound and the next-step Gaussian gap with the same uniform
constants. There is also a homogeneous independent-current
extension obtained by using \(\omega_0\) in the second line
of \eqref{mb:eq-physical-outer-form}. Its gap and continuation
have the same proof using the first readout estimate.
The two definitions agree on physical pairs by the
exact stationarity identity. The linear estimates require
only the specified Hessian and chart data; the existence
of the critical branch is a separate nonlinear input.


% Construction source: audit/conditional-matrices.tex
\section{Conditional matrix functions and complex Gaussians}
\label{mb:sec-conditional-matrix}

Use normalized counting coordinates on the outer free bonds. Let \(A_Y\) be the precision on the indicated selected free domain, with zero extension of those coordinates when it is formed from the full outer pivot graph. Thus compression precedes inversion. A subsequent coordinate restriction \(\chi_W\), \(W\subseteq Y\), takes the covariance marginal after conditioning on \(Y\), at complete parameters. The parameter extensions below have guards that depend on \(W\); away from complete parameters no marginal consistency between different choices of \(W\) is asserted. On a tree-zero domain the missing tree coordinates are included in the mask throughout.

For the proper outer form \eqref{mb:eq-physical-outer-form}, the two premises for the following construction are the real gap of \cref{mb:thm-physical-outer-gap} and a connected family for that same precision. To obtain its family use \eqref{mb:eq-mass-removed-family} for the second row Schur, and the second-lift rooted readout of \cref{mb:thm-current-families} for its pivot curvature. The actual pivot derivative, pivot second jet and relative contour coefficient have fixed one-step carriers. The current root and its fine input cell stay in each coefficient. The compression mask is read on that carrier before the two outer coordinate restrictions are applied. The resulting family has the required row and column sums with local coefficient suprema. The unit outer scale and the bounded number of outer variables per cell allow the coarse short-power construction. In a proper hierarchy the short endpoint balls can be chosen inside the last inner region by increasing the fixed separation after their radii have been chosen. Their background cutoff extension uses the outer unit scale there, and their long words retain all mixed-scale reads. These verifications also apply to a homogeneous inner hierarchy with its specified precision and its own stipulated current lift in the pivot contraction. The coefficient and real-gap premises concern each particular matrix on its own domain.

\begin{theorem}[Filed covariance and logarithm]
\label{mb:thm-filed-covariance}
Let the real symmetric precision \(A_Y\) on normalized outer free coordinates have \(aI\le A_Y\le b_0I\), and let its exact connected family, local coefficient tube and selection mask satisfy \cref{mb:lem-coarse-carrier-inverse}. This includes \eqref{mb:eq-physical-outer-form} on its admitted real and complex domains. For each fixed cube radius \(\mathcal S>2\) there are fixed aligned parameter cubes and complete root guards such that the following extensions exist.

Use one common admitted real comparison on $Y$ for the covariance sector. The separate real completions of individual entries establish local supremum bounds; they do not replace this common comparison. The filed inverse and logarithm are normally convergent on \(|s_\Delta|<\mathcal S\). At all ones they give \(A_Y^{-1}\) and the logarithm continued from its positive real value. The logarithm has a uniform diagonal trace per outer cell. For every finite selected coordinate set \(W\subseteq Y\) there is a holomorphic symmetric covariance \(D_W(s)\) on the same polydisc, with
\begin{equation}\label{mb:eq-selected-covariance}
             D_W(1)=\chi_W A_Y^{-1}\chi_W^T .
\end{equation}
For a real positive comparison matrix \(D_{0,W}\) with \(b^{-1}I\le D_{0,W}\le CI\),
\begin{equation}\label{mb:eq-gaussian-relative-sector}
 \big\|D_{0,W}^{-1/2}(D_W(s)-D_{0,W})D_{0,W}^{-1/2}\big\|<\tfrac14
\end{equation}
holds uniformly in \(W\), depth and compatible volume on the common smaller coefficient tube. The difference can be made smaller than any prescribed inverse power of the parameter cube width. The logarithm and covariance coefficients retain their declared reads. Zeroing parameters outside a wall gives the corresponding guarded wall dependence and factorization across separated guarded components.
\end{theorem}

\begin{proof}
Use the construction of \cref{mb:lem-coarse-carrier-inverse} with \(H=I-A_Y/b\), \(b>2b_0\), short matrix \(S\) and long differences \(E_j=H^j-S^j\). In particular \(S_*\) at the real comparison data is symmetric and \(0\le S_*\le q_*I\) with \(q_*<1\). A prefix \(S^j\) is filed on the full endpoint balls before absolute summation. At a common real datum its row and column sums are at most \(C(1+B_j)^2 q_*^j\). With independent local data before the sums use \(C(1+B_j)^4(q')^j\), \(q_*<q'<1\), with the real-gap collars and the smaller complex tube there. For fixed \(R_0\) these are respectively quadratic and fourth-power polynomial factors in \(j\), with constants depending on \(R_0\). Both forms are uniform in the inner mesh. The normalized coordinate count here is the count of outer free slots.

Take \(R_0=D n\) and file the combined \(n\)th power \(H^n=S^n+E_n\). Mark each of its coefficients by the product of its declared cube parameters, with parameters of the held guard set to one, and denote the resulting matrix by \(T_s\). The new cube width \(M_c\) can be taken comparable to a fixed multiple of \(M_0n(R_0+1)\). The short combined coefficients touch a bounded number of these cubes. The old long carrier reserve pays the other marks by \eqref{mb:eq-reblock-count}. Thus on the whole parameter polydisc
\[
 \|T_s\|_{\rm Sch}\le
   C_{\mathcal S,L}(1+B_n)^4(q')^n+CnK^n e^{-cDn}=:r_n .
\]
Use \(D\) fixed sufficiently large and then \(n\) large. The bound tends to zero exponentially up to a polynomial in \(n\). The same estimates apply to the parameterized \(E_j\), \(j<n\), with bounds \(CjK^j e^{-cDn}\). The geometry can be increased after these finite choices, and all the normalized local estimates retain their fixed margins.

For the logarithm use the joined words in
\begin{equation}\label{mb:eq-filed-log}
 \log A_Y=(\log b)I-\sum_{j=1}^{n}H^j
                         \sum_{v\ge0}\frac{(H^n)^v}{nv+j}.
\end{equation}
The identity at complete parameters follows by grouping the absolutely convergent scalar series of \(I-H\) in the matrix. Away from complete parameters the right side means the chosen filed extension, with the parameters of each joined carrier. Its absolute convergence follows from the grouped contraction and the finite prefix bounds. Each diagonal trace is a sum of closed carriers through that outer coordinate, and the family bound and bounded number of coordinates per outer cell prove the trace assertion. The inverse uses the analogous grouped expansion \eqref{mb:eq-coarse-carrier-inverse}. The two algebraic identities at all ones use the full matrices.

For the covariance construction take the complete radius-\(B_n\) guards around the selected roots \(W\), including every short-prefix read and its domain and geometry collars. Hold the short prefixes at their unmarked values, and for \(j<n\) put \(A_{j,s}=S^j+(E_j)_s\), with \(E_0=0\). The coefficients of \((E_j)_s\) retain their own joined parameter carriers. Use \(c_\ell=4^{-\ell}\binom{2\ell}{\ell}\) and put
\begin{equation}\label{mb:eq-covariance-root-series}
 \begin{aligned}
 R_W(s)&=b^{-1/2}\chi_W
             \sum_{j<n} A_{j,s}\sum_{v\ge0}c_{nv+j}T_s^v,\\
 D_W(s)&=R_W(s)R_W(s)^T .
 \end{aligned}
\end{equation}
These are ordinary matrix products of the indicated marked blocks, and \(T\) is bilinear. The powers in this definition can repeat a cube parameter. This is a matrix-function extension with held prefixes. The distinct-cube marking of an already filed operator family in \cref{mb:cor-carrier-parameters} is the other convention. Both retain the same support union for every product. By \(0<c_\ell\le1\) and \(r_n<1/2\) the series in \eqref{mb:eq-covariance-root-series} converges normally with the weighted local sums. At all ones \(A_{j,1}=H^j\), \(T_1=H^n\). Grouping the binomial series of \((I-H)^{-1/2}\) gives \(R_W(1)=\chi_W A_Y^{-1/2}\). The square with \(T\) proves \eqref{mb:eq-selected-covariance}.

For the comparison sector put \(R_*=b^{-1/2}\sum_{j<n}c_jS_*^j\). Since \(S_*\ge0\), \(R_*\ge b^{-1/2}I\). Its norm is bounded by \(b^{-1/2}\sum_{j\ge0}c_jq_*^j\). Hence \(D_{0,W}=\chi_W R_*^2\chi_W^T\) has both claimed real bounds, by zero extension of a vector in \(W\). The long prefixes and the \(v\ge1\) terms in \eqref{mb:eq-covariance-root-series} have norm at most a fixed polynomial in \(n,R_0\) times \(r_n+nK^n e^{-cDn}\), since the short-prefix sums have polynomial bounds. The finite unmarked prefix differs from the real comparison prefix by at most \(C\sum_{j<n}\Delta_j^+\). Decrease the common coefficient tube after fixing \(n\), so this term is bounded by the same tail budget. It follows that \(R_W(s)-\chi_W R_*\) is small in operator norm, uniformly in \(W\). Expand \(R_W R_W^T-\chi_W R_*^2\chi_W^T\), use the uniform bound for \(R_*\) and the lower bound for \(D_{0,W}\), and obtain \eqref{mb:eq-gaussian-relative-sector}. The polynomial growth of \(M_c\) in \(n\) gives the claimed stronger budget by taking \(n\) larger before decreasing the tube.

On a wall, a marked factor survives only if every one of its parameter cubes is allowed, with the held guards included. Its full filed carrier remains available there. The unmarked prefix starting at \(W\) stays in its root guard. Connectedness then forbids a surviving path between distinct components once intersecting guarded read sets have been joined. The covariance on their selected roots is block diagonal and its normalized integrals factor. Agreement of complete wall and guard reads proves coefficient equality by \cref{mb:cor-determining-comparison}. For a small period use the common cover and the deck average from the family proof, retaining all lifted root balls and parameters on their periodic images. Descend the square-root matrix before selecting \(W\). Its periodic sector is reducing for it and its bilinear transpose. With isometric pullback \(I_{\rm per}\) this gives
\[
 I_{\rm per}^T R_{\rm cov} R_{\rm cov}^T I_{\rm per}
                       =R_{\rm base}R_{\rm base}^T .
\]
Thus the subsequent selection gives the covariance of the actual base precision at complete parameters. The absolute row and column sums descend with no sheet-count factor.
\end{proof}

\begin{corollary}[The finite-range first-step precision]
\label{mb:cor-first-step-local-gaussian}
In the actual one-step chart of \cref{mb:cor-contour-gap}, let \(K_A\) be the Wilson and constraint-curvature part of \eqref{mb:eq-one-step-contour} and let \(K_G=(b_VC_V)^T(b_VC_V)\). For \(\tau\) near one the precision \(K_c=\tau K_A+K_G\), compressed to any admitted selection of its free coordinates, has the local inverse root, covariance and logarithm of \cref{mb:thm-filed-covariance}. Here one can use localization cubes of side \(O_L(n)\), and the local background and independent current variables are admitted on fixed absolute small tubes in the one-step units.
\end{corollary}

\begin{proof}
The centered independent-current action in the exact pivot chart \(\Phi_V\) is
\[
 A_V(\Phi_V(u))-A_V(0)
       +\langle J-j_V,\Phi_V(u)\rangle-\langle J,C_Vu\rangle,
 \qquad j_V=DA_V(0).
\]
Its first derivative is zero, and its second derivative is \(C_V^T H_W(V) C_V+\langle J,\Phi_V''(0)\rangle\). Since \(\Phi_V''(0)=-h_VD^2q_V(0)[C_V\,\cdot,C_V\,\cdot]\), this is \(K_A\) with its actual constraint curvature. In the contour term use the product of the contour mean at \(\Phi_V(u)\) and the inverse of the reference contour mean. Its value is the identity at zero, so the Hessian of the trace gauge term is the Gram form \(K_G\). At the real reference with \(\tau=1\), the gap and upper bound are those of \cref{mb:cor-contour-gap}.

The real gap for a field supported on a free coordinate subset uses only the fixed complete block neighborhoods of its support. Its pivot reconstruction uses one-step carriers, and the curl-contour inequality is the sum of estimates on these blocks and their neighbors. Its real-frame comparison and small curvature errors therefore require a fixed additional collar \(R_{\rm gap}=C_L\). Compression is evaluation on the zero extension within the original pivot chart, including the full crossing rows. This gives the same real bounds on every such local coefficient matrix. Real contour gauges act unitarily on the physical free coordinates, so the same real bounds hold in those coordinates.

The pivot derivatives, Wilson stencil and relative contour map in \eqref{mb:eq-one-step-contour} have bounded finite range \(R_K=C_L\) in one-step units, including all data reads. Each entry is holomorphic in a fixed finite stock of the actual local connection and independent current. If \(\mathbf V=e^{iE}V\) and \(\mathbf J=J+Z\), bounded derivatives of these chart coefficients on a smaller local tube give the positive entry envelope
\[
 \|K_c-K_r\|_{\mathrm{Sch},\sup}
      \le C_L(\|E\|_\infty+\|Z\|_\infty+|\tau-1|)
      =:\delta_K^+ .
\]
Here the supremum is taken before summing the endpoints; bounded incidence gives the same \(C_L\) in both sums. This bound uses the relative contour completion and the constraint-curvature term in \(K_A\). The increment of \(\tau\) multiplies just \(K_A\).

Choose \(b>2\|K_r\|\) uniformly, set \(H_r=I-K_r/b\), \(H_c=I-K_c/b\), and \(r=1-\gamma/b<1\) from the real lower gap \(\gamma\). Take the full-carrier cutoff above the range of the primitive entries. Then \(E_j=0\), \(S=H_c\), and the endpoint balls for \(H_c^j\) have radius at most \((j+1)R_K+R_{\rm gap}\). Combining the powers before taking absolute values gives \(C_L(1+j)^2r^j\) at one common real pair, and \(p_j^+=C_L(1+j)^4r^j\) with separate local suprema. The latter bound follows entry by entry by taking the indicated local real compression with its complete gap collar. If \(C_H\) bounds the positive envelopes of \(H_r,H_c\), telescoping their powers gives the coefficientwise difference bound
\[
 \Delta_j^+\le b^{-1}\delta_K^+\,jC_H^{j-1}.
\]
After fixing the finite \(n\), choose the absolute field, current and \(\tau\) radii so \(\Delta_n^+\le r^n\) and \(\sum_{j<n}\Delta_j^+\le r^n\).

Take the cube side to be a power of \(L\) above a fixed multiple of \((n+1)R_K+R_{\rm gap}\), and keep the union of the corresponding padded cubes at each root as its held guard. The side is \(O_L(n)\); the number of cubes met by one combined \(n\)th-power entry is bounded independently of \(n\). Its parameterized norm is at most \(2\mathcal S^{C_L}p_n^+\). The sum of the prefix norms is bounded by \(\sum_{j\ge0}p_j^++1\). These are precisely the grouped-step and held-prefix bounds used in \eqref{mb:eq-covariance-root-series} and \eqref{mb:eq-filed-log}. Their exponentially small error beats every fixed inverse power of the new cube width, on the stated common smaller absolute tube.

These bounds also hold when the real and complex pair is given only on an output domain and its declared reads. For a combined entry rooted at \(y\), the compression to the radius-\(nR_K\) endpoint ball contains every length-\(n\) coefficient path from \(y\). Hence its \(n\)th power has the same \(y\)-row as the original polynomial on that row's support. The ball's real-gap collar is already in the entry carrier. For the held reference prefix on \(W\), take one common compression to the radius-\(nR_K\) neighborhood of \(W\) inside its guard. The same path observation identifies every prefix row and gives the real comparison bounds there. Thus deleting entries whose complete carriers lie outside the wall reduces the positive majorants, and the fixed-root comparison uses only the supplied real data in the guard.
\end{proof}

\begin{corollary}[Normalized local Gaussian]
\label{mb:cor-complex-gaussian}
In the sector \eqref{mb:eq-gaussian-relative-sector}, on \(d=\dim W\) real coordinates the holomorphic Gaussian density is
\[
 d\mu_{D_W}(u)
    =(2\pi)^{-d/2}\det(D_W)^{-1/2}
                       \exp(-u^TD_W^{-1}u/2)\,du .
\]
Its determinant branch is the one continued from \(D_{0,W}\). Its integral is one and its total variation is at most \(C^d\), uniformly on the stated complex tube and parameter polydisc. Integration against any fixed bounded Borel test on these real coordinates gives a holomorphic function. For \(W=\varnothing\) the integral is one.
\end{corollary}

\begin{proof}
Write \(D_W=D_{0,W}^{1/2}(I+E)D_{0,W}^{1/2}\) with \(\|E\|<1/4\). Use the convergent matrix logarithm of \(I+E\) to define the determinant factor. The singular value lower bound gives \(|\det(I+E)|^{-1/2}\le(4/3)^{d/2}\), and \(\|(I+E)^{-1}-I\|\le1/3\). After the real change \(u=D_{0,W}^{1/2}v\), the absolute value of the exponential is bounded by \(\exp(-|v|^2/3)\). Thus its integral is absolutely bounded by \(C^d\) and is holomorphic in the complex symmetric entries on every smaller sector. It equals one on the open real positive slice, so it equals one throughout the sector by holomorphic uniqueness in the complexification of that real space. The same dominating Gaussian gives holomorphic polynomial moments and parameter derivatives on smaller tubes. The trace bound \(|\operatorname{Tr}\log(I+E)|\le 2d\|E\|\) follows from the logarithm series. Since a wall cube contains at most \(C_L M_c^4\) outer coordinates, a sector budget \(o(M_c^{-4})\) yields the corresponding small log-determinant error per wall cube. The count is on selected outer free coordinates, independent of the inner fine mesh.
\end{proof}


% Construction source: draft/retained-seam.tex
\section{Additional rows at a retained inward seam}\label{mb:retained-seam}

The proper determining set has no inward rows. An optional
convention retains some inward bonds. At each additional seam
node choose one as a tying bond, oriented toward that node.
Its equation determines the node factor. The independent remaining
equations are the quotients of the other inward bonds by the
tying bond. We call them edge rows. A node with only one retained
inward bond contributes no edge row.

In this section all constraint coordinates are raw logarithms,
and both row spaces use counting measure. Let $K_p$ and $K_e$
be the differentials of the proper and edge rows. Their
transposes map counting row covectors to the physical fine-field
dual. If an edge $e$ is the quotient of $b_i$ by its tying bond
$b_1$, the quotient rule gives
\begin{equation}\label{mb:seam-differential}
 K_eA(e)=\widehat{\dot a_q[A]}(b_i)
       -\operatorname{Ad}(V_0(e))\widehat{\dot a_q[A]}(b_1).
\end{equation}
Here $a_q$ is the actual reduced logarithm at level $q$,
$V_0(e)=W_q(b_i)W_q(b_1)^{-1}$, and $W_q=\mathcal M^qU_0$.
A reversed stored bond uses
$\widehat a_q(\bar b)=-\operatorname{Ad}(W_q(b)^{-1})a_q(b)$.

Fix an odd $L\ge13$. Leave $\Omega'_0=\Omega_0$.
For $q\ge1$, retain precisely the native aligned level-$q$
blocks $B\subset\Omega_q$ with
$\operatorname{dist}_\infty(B,\Omega_q^c)\ge3\xi_q$.
Their union is $\Omega'_q$. Distances are periodic; a full
torus has distance infinity from its empty complement and is
unchanged. No additional localization grid is needed.
Choose the original separation floor larger than five, and larger still as required by the fixed proper-inverse collars.
Then
\begin{equation}\label{mb:peeled-scales}
 \Omega_{q+1}\subset\Omega'_q\subset\Omega_q,
              \qquad \xi'(x)\le\xi(x)\le L\xi'(x).
\end{equation}
Indeed an original level-$q$ block meeting $\Omega_{q+1}$
loses at most $\xi_q$ from that region's separation from
$\Omega_q^c$, so it is retained. Every discarded point is
within $4\xi_q$ of $\Omega_q^c$. For a point in
$\Omega'_{q+1}$, a shortest path to $\Omega_q^c$ first
travels at least $3\xi_{q+1}$ before exiting
$\Omega_{q+1}$ and then crosses the original separating
layer. Subtracting the discarded strip gives a remaining
separation of at least $(\rho+3L-4)\xi_q$.
Thus the peeled sequence retains the required separation.
A site loses at most one level, and complete peeled pillars
subdivide the old ones.

\begin{proposition}[Uniform edge Schur correction]\label{mb:edge-schur}
Suppose $U_0$ is real and the original pair $(U_0,j_{U_0})$
and peeled pair $(U_0,0)$ obey the respective regular-pair
coefficient hypotheses of \cref{mb:thm-layered-gap,mb:thm-strong-families}, with a common small
radius $\varepsilon_s$. The first current is the actual inner Wilson
current; the second is an independent zero current. Require
these admissions on every complete local box used by either
proper inverse. Set
\begin{align*}
 H_{\rm seam}&=H_W+DRD^{\mathsf T},&
 T&=H_{\rm seam}+K_p^{\mathsf T}K_p,& G_p&=T^{-1},\\
 H_p&=G_pK_p^{\mathsf T}(K_pG_pK_p^{\mathsf T})^{-1},&
 C_p&=G_p-H_pK_pG_p,& M_e&=K_eC_pK_e^{\mathsf T}.
\end{align*}
Theorem~\ref{mb:thm-layered-gap} gives $T>0$ with
$a_f=\Xi_f^{-2}$: because $K_p=\Xi_fQ$ uses counting row
coordinates, $K_p^{\mathsf T}K_p=Q^{\mathsf T}\Xi_f^{-2}Q$.
The proper field and lift estimates used below are the
zero-current case of \cref{mb:thm-strong-families}. There are $c,C,s>0$, independent of mesh, depth,
volume and the retained subset, such that
\begin{equation}\label{mb:edge-gram-gap}
 cI\le M_e\le CI,\qquad
             |M_e^{-1}(e,e')|\le Ce^{-s d(y_e,y_{e'})}.
\end{equation}
Define
\begin{align}
 H_a(b_p,b_e)&=H_pb_p+C_pK_e^{\mathsf T}M_e^{-1}
                         (b_e-K_eH_pb_p),\notag\\
 C_a&=C_p-C_pK_e^{\mathsf T}M_e^{-1}K_eC_p.
                                           \label{mb:edge-correction}
\end{align}
For $K=(K_p,K_e)$ these maps satisfy
\begin{equation}\label{mb:edge-identities}
 KH_a=I,\quad KC_a=0,\quad C_aH_{\rm seam}=I-H_aK,
                     \quad\ker C_a=\operatorname{ran}K^{\mathsf T}.
\end{equation}
They have the proper raw-datum-to-field and force-to-field
pointwise, first-difference and local gradient H\"older bounds,
with a smaller positive metric rate. If the edge space is zero,
the correction is zero. On a common smaller complex tube the
matrices $G_p$ and $M_e^{-1}$ remain bounded and holomorphic.
The maps $H_a,C_a$ continue holomorphically with the stated
norm bounds and algebraic identities.
\end{proposition}
\begin{proof}
Let $H'_p$ be the proper raw lift on the peeled sequence.
For a field set
\[
 \|A\|_{\mathcal E}^2
   =\sum_y\max\{\xi_y\|A\|_{\infty,B_y},
                    \xi_y^2\|\nabla_{U_0}A\|_{\infty,B_y^+}\}^2.
\]
Use the fixed complete collars from the proper estimates.
Their overlap is bounded. The exponential lift bounds and
\eqref{mb:rate-reserve} give
$\|H_pb\|_{\mathcal E}+\|H'_pb'\|_{\mathcal E'}
 \le C(\|b\|_2+\|b'\|_2)$.
By \eqref{mb:peeled-scales}, every peeled cell meeting an
old enlarged cell has a comparable scale, with constants
depending only on $L$ and the fixed enlargement. Bound the
old maximum by the sum of squares over all those peeled
cells. Their bounded overlap gives
$\|A\|_{\mathcal E}\le C_L\|A\|_{\mathcal E'}$.

We construct trial data for an edge row. Write its nontying
inward bond as $(y_i,x)$, where
$x=y_i+\sigma\xi_qe_\mu$. In the removed strip the three
consecutive bonds
\[
 (y_i,x),\quad(x,x+\sigma\xi_qe_\mu),\quad
 (x+\sigma\xi_qe_\mu,x+2\sigma\xi_qe_\mu)
\]
are proper interior rows of the peeled geometry. Their three
successive inner endpoints lie in the first native level-$(q+1)$
block entered by the inward bond, since $L\ge13$.
That boundary block is removed from $\Omega'_{q+1}$.
The original separation places this block and $y_i$ in
$\Omega'_q$. Give these three rows the data
$(1,-2,1)v$ in the frames transported along that line.
Use the covariant minus sign if a stored orientation is reversed.
Call this assignment $J_ev$. Distinct directions at a seam node
use distinct bonds. Parallel inward bonds on the same block
line are at least $L-1$ bonds apart. Thus, for $L\ge13$, the
three-bond supports are disjoint and avoid tying bonds. Hence
$\|Jb\|_2^2=6\|b\|_2^2$ for $Jb=\sum_eJ_eb(e)$.

At a flat reference the straight-path average of this pattern
vanishes at the next level. The paths through the first
endpoint block meet its initial segments of lengths one, two
or three. Their sums are $1,-1,0$; the two nonzero sums occur
once each. For the next coarse bond the two nonzero sums are
$-1,1$. The averaging coefficients are equal, so each pair
cancels. The transverse averaging factor multiplies zero.
This cancellation uses the boundary placement of the pattern,
not an arbitrary three-bond translate.

At a curved reference, compare the paired transports in their
complete double block. Their difference is a loop with scaled
curvature $O(\varepsilon_s)$, so the transport comparison costs
$C_L\varepsilon_s$. The derivative of the actual straightened
one-step mean differs from that straight-path map by the same
amount in its finite row and column sums. To see this, work in
the admitted local comb frame and differentiate the mean
equation and both endpoint factors. At zero curvature the
endpoint derivatives cancel the internal contour derivatives.
The remaining derivatives are analytic on a fixed smaller
finite chart, so the curvature error is $O_L(\varepsilon_s)$.
The derivative of the implicit mean has a uniformly bounded
inverse on that chart. Real gauge covariance returns the
estimate to the original frame.

Put $A_b=H'_pJb$. Complete reads shared by the two geometries
give the same reduced logarithm, by the exact averaging
recursion. Thus \eqref{mb:seam-differential} reads $b$ on each
nontying bond and zero on its tying bond. An old proper row
shared with the peeled set reads zero. If one of an old
row's level-$j$ endpoint blocks survives in $\Omega'_j$,
the row remains an interior or outward row. An omitted old
row therefore has both endpoint blocks outside $\Omega'_j$.
Every original bare row is shared, so an omitted row has
$j\ge1$. Its full predecessor double block lies in
$\Omega'_{j-1}$ by the
original separation, so its predecessor rows are proper
interior rows there. It reads the preceding pattern through
that complete double block. The cancellation just proved and the uniform
actual row-composition bounds therefore give
\begin{equation}\label{mb:edge-trial}
 K_eA_b=b,\qquad
 \|K_pA_b\|_2\le C_L\varepsilon_s\|b\|_2,\qquad
 \|A_b\|_{\mathcal E}\le C\|b\|_2.
\end{equation}
Bounded double-block overlap supplies both Schur sums. No
factor depending on the number of levels is introduced.

Correct into the proper kernel by
$A=A_b-H_pK_pA_b$. The local edge-row bound and
\eqref{mb:edge-trial} give
$K_pA=0$, $\|A\|_{\mathcal E}\le C_A\|b\|_2$ and
$\|K_eA-b\|_2\le C_e\varepsilon_s\|b\|_2$.
Choose $C_e\varepsilon_s\le1/2$. On this kernel the row mass vanishes,
and the Wilson curl comparison and orthogonality of $R$ give
$\langle A,TA\rangle_f\le C_T\|A\|_{\mathcal E}^2$.
The physical cell volume contributes $\xi_y^4$, cancelling
the four inverse length powers from the squared first
differences. Adjacent scales have bounded ratio.

The constrained inverse of $T$ has the variational formula
\[
 \langle f,C_pf\rangle_f
   =\sup_{K_pu=0}\{2\operatorname{Re}\langle f,u\rangle_f
                                      -\langle u,Tu\rangle_f\}.
\]
Take $f=K_e^{\mathsf T}b$ and $u=tA$. The real pairing
$\langle b,K_eA\rangle$ is at least $\|b\|_2^2/2$.
Choosing $t=(2C_TC_A^2)^{-1}$ proves
$M_e\ge(4C_TC_A^2)^{-1}I$.
Its upper bound and exponential kernel follow by composing
the proper constrained inverse with the local edge maps.
The counting-to-physical transposes include the fine density.

For the inverse decay, conjugate $M_e$ by
$e^{s\min\{d(y_e,y_{e_0}),N\}}$.
The exponential kernel and its first distance moment bound
the conjugation error by $Cs$, uniformly in $N$.
Choose $s>0$ so that $Cs\le c/2$ and so that the retained
kernel rate has the reserve in \eqref{mb:rate-reserve}.
The conjugated inverse has norm at most $2/c$.
Apply it to an edge coordinate at $e_0$ and let $N$ increase.
This proves \eqref{mb:edge-gram-gap}. If necessary increase
the fixed separation so that the final smaller rate is also
summable. There are only finitely many edge species at a top.

Direct multiplication in \eqref{mb:edge-correction} proves
the first three identities in \eqref{mb:edge-identities}.
The form $T$ is positive on $\ker K$, and $C_a$ is its
constrained inverse there. Its kernel is therefore the
physical annihilator of $\ker K$, namely
$\operatorname{ran}K^{\mathsf T}$. Composing the proper
exponential kernels proves the asserted field bounds. For a
local gradient H\"older estimate, regard each correction as
$C_pf$ and apply the corresponding proper force estimate.
The row datum is not differentiated.

Finally the proper inverse and local rows have small uniform
perturbations on a common smaller complex tube. The same is
true of $M_e$ in the normalized weighted row norm. A Neumann
series preserves invertibility and the displayed algebraic
formulas. Their transposes are bilinear, so the continuation
is holomorphic.
\end{proof}

\mbneedspace{20\baselineskip}
\begin{corollary}[The corrected seam saddle]\label{mb:seam-saddle}
Suppose the real reference is constrained critical for the
combined raw constraint $\mathfrak q$, with
$\mathfrak q'(0)=K$ and actual multiplier
$\omega_0=H_a^{\mathsf T}j_{U_0}$. Let $B$ be defined by
\[
 \langle v,Bu\rangle_f
       =\langle\omega_0,\mathfrak q''(0)[u,v]\rangle.
\]
Fix $0<\alpha<1$ and use the strong field norm
$\|u\|_{\mathcal A_\alpha}=\sup_y\|u\|_{\mathcal A_{\alpha,y}}$
on the complete boxes above. Suppose the combined raw constraint
uses the actual admitted mean and quotient charts with their
two-leg moment bounds. If the scaled actual current on the
touching stock is sufficiently small, the saddle
$(H_{\rm seam}-B)u-K^{\mathsf T}\nu=f$, $Ku=b$ has the unique solution
\begin{equation}\label{mb:seam-saddle-inverse}
 u=(I-C_aB)^{-1}(C_af+H_ab),\qquad
 \nu=H_a^{\mathsf T}((H_{\rm seam}-B)u-f).
\end{equation}
It obeys the proper force, raw-datum and multiplier estimates,
also in a fixed smaller exponential weight.
\end{corollary}
\begin{proof}
Here is the bound for the added quotient rows. At a row of
length $\xi_q$, write the two oriented reduced raw logarithms
as $a_i,a_1$. Its relative edge logarithm is
\[
 i^{-1}\log\bigl(e^{ia_i}e^{-i\operatorname{Ad}(V_e)a_1}\bigr).
\]
The second differential consists of the difference of the two
second jets and the symmetric cross commutator of their first
jets, with coefficient $-i/2$. Each raw first jet has absolute
sum at most $C\xi_q$ and each one of its coefficients has norm
at most $C\xi_q(\eta/\xi_q)^4$. Each raw second jet has
fixed-leg sum at most $C\xi_q^2(\eta/\xi_q)^4$.
The cross commutator has the same fixed-leg bound by multiplying
the first two estimates. These are the actual row moments used
in \eqref{mb:eq-two-leg-energy}, converted to raw rows.
The physical fine transpose divides by $\eta^4$.
Inserting $|u|\le C\xi_q^{-1}\|u\|_{\mathcal A_\alpha}$
therefore gives a force bounded by
$C\xi_q^{-3}|\omega_0(e)|\|u\|_{\mathcal A_\alpha}$.
At one level only a bounded number of full seam carriers read
a given fine bond. Different levels have separated seam
boundaries, so this incidence bound is uniform in depth.
Thus the combined Hessian has the same force bound as the
proper constraint Hessian.

For a nonbare or edge unit raw datum, $KH_a=I$ makes its lift
vanish on every nontouching bare bond. Its current readout
therefore uses only the touching current. The column bound
for $H_a$ and the physical scale cancellation
$\xi^{-1}\xi^{-3}\xi^4=1$ give
$|\omega_0(e)|\le C\varepsilon_s$ on every row whose second jet
is nonzero. Bare proper rows have zero second jet.
The force-to-field estimate for $C_a$ now yields
$\|C_aB\|_{\mathcal A_\alpha\to\mathcal A_\alpha}
\le C_\alpha\varepsilon_s$.
A fixed reduction of the current ceiling makes this at most
$1/2$. The same calculation uses the smaller exponential rate
because every quotient reads its complete local carrier.

Apply $C_a$ to the first saddle equation and use
\eqref{mb:edge-identities}. This gives the first formula in
\eqref{mb:seam-saddle-inverse}. Conversely that formula gives
$Ku=b$, and its first-row residual lies in
$\ker C_a=\operatorname{ran}K^{\mathsf T}$.
Since $KH_a=I$, applying $H_a^{\mathsf T}$ gives the second
formula and the unique multiplier.

The Neumann series gives the field bound. To estimate the
multiplier without taking a second pointwise derivative of
$u$, use the variational pairing: one curl difference falls
on $H_ae$ and one on $u$. For the gauge term use
$\langle D^{\mathsf T}H_ae,RD^{\mathsf T}u\rangle_f$.
The two first-difference amplitudes and physical cell volume
again cancel their scale powers. The remaining curvature
and constraint terms have their force bounds. The column
estimate controls $H_a^{\mathsf T}f$. The same proof with
one exponential weight on the changed factor proves the
weighted estimate.
\end{proof}


% Construction source: draft/finite-window.tex
\section{Restrictions and physical gauge null directions}\label{mb:finite-window}

\subsection{A marginal precision retains the gap}

Work first in real physical orthonormal coordinates on a finite
index set. Let $A=A^*$ satisfy $a_-I\le A\le a_+I$, where
$a_->0$, and assume that its entries have uniformly bounded
two-sided exponential sums at rate $\mu>0$. Split the indices
into $I$ and $O$. Empty blocks have their usual zero-dimensional
meaning. The covariance restricted to $I$ has precision
\begin{equation}\label{mb:marginal-schur}
 K_I=((A^{-1})_{II})^{-1}
       =A_{II}-A_{IO}A_{OO}^{-1}A_{OI}.
\end{equation}
This is different from the conditional precision $A_{II}$.

\begin{proposition}[Uniform restrictions and weighted perturbations]
\label{mb:window-schur}
The bounds $a_-I\le K_I\le a_+I$ hold for every subset $I$.
There are $s_0,C>0$, independent of that subset and the admitted
volume, with the following property. For every bounded real
one-Lipschitz function $\psi$ and $|s|\le s_0$,
\begin{equation}\label{mb:window-weighted}
 \|e^{s\psi}K_Ie^{-s\psi}-K_I\|\le C|s|.
\end{equation}
If a possibly complex perturbation $B_I$ satisfies
$\|e^{s\psi}B_Ie^{-s\psi}\|\le\varepsilon$, then
\begin{equation}\label{mb:window-perturbation}
 \operatorname{Re}\bigl(v^*
        e^{s\psi}(K_I-B_I)e^{-s\psi}v\bigr)
          \ge(a_--C|s|-\varepsilon)\|v\|^2.
\end{equation}
In particular, $C|s|+\varepsilon\le a_-/2$ gives an inverse
bound $2/a_-$ in that weighted space.
\end{proposition}
\begin{proof}
The positive block $A_{OO}$ is invertible. Solving the second
block equation in $A(u,w)=(f,0)$ gives
$w=-A_{OO}^{-1}A_{OI}u$ and proves
\eqref{mb:marginal-schur}. Equivalently,
\[
 \langle u,K_Iu\rangle
       =\inf_w\langle(u,w),A(u,w)\rangle.
\]
The lower gap follows by dropping $\|w\|^2$ from the lower
bound on $A$. Taking $w=0$ proves the upper bound.

Let $A_s=e^{s\psi}Ae^{-s\psi}$. The entrywise inequality
$|e^{s(\psi(x)-\psi(y))}-1|
 \le |s|d(x,y)e^{|s|d(x,y)}$
and the exponential Schur bound imply
$\|A_s-A\|\le C_0|s|$ for $|s|\le\mu/2$.
Compression preserves that estimate. Choose $s_0$ so that
$C_0s_0\le a_-/2$. The resolvent identity then gives
\[
 \|(A_s)_{OO}^{-1}\|\le2/a_-,\qquad
 \|(A_s)_{OO}^{-1}-A_{OO}^{-1}\|
                       \le2C_0|s|/a_-^2.
\]
Conjugation commutes with the block decomposition. Insert these
bounds and the bounds for the two off-diagonal blocks in
\eqref{mb:marginal-schur}. Expanding the difference of the
three-factor products proves \eqref{mb:window-weighted}.
Taking the real part and using the norm bound for $B_I$ proves
\eqref{mb:window-perturbation}.
\end{proof}

Taking $\psi(x)=\min\{d(x,y),N\}$ in the same argument gives
pointwise exponential bounds for the restricted inverse, uniformly
in $N$. Sending $N$ to infinity removes the truncation. When a
two-sided cell sum is required, the metric growth must also be
paid for as in \eqref{mb:rate-reserve}. A small unweighted norm
of $B_I$ alone gives no uniform bound for its conjugates.

\subsection{The complete Hessian on an actual fibre}

Let $\mathfrak q(a)$ be a finite analytic raw constraint, with
$\mathfrak q(0)=0$, surjective derivative $P$, and pivot right
inverse $h$, so $Ph=I$. Let $C$ identify the free coordinates
with $\ker P$. The actual pivot chart has the form
$\Phi(u)=Cu+hz(u)$, where $z(0)=Dz(0)=0$ and
$\mathfrak q(\Phi(u))=0$.
Let $\mathcal A$ be an invariant action that is constrained
critical at the reference. Its gradient is $P^{\mathsf T}\omega$.
The reduced action $\mathcal S(u)=\mathcal A(\Phi(u))$ therefore
has Hessian
\begin{equation}\label{mb:fibre-hessian}
 K_A[v,w]=D^2\mathcal A(0)[Cv,Cw]
               -\langle\omega,D^2\mathfrak q(0)[Cv,Cw]\rangle.
\end{equation}
Indeed twice differentiating the constraint gives
$D^2z(0)[v,w]=-D^2\mathfrak q(0)[Cv,Cw]$. Pairing the
second chart derivative with the ambient action gradient proves
\eqref{mb:fibre-hessian}.

For the gauge splitting use the complete one-step outer
fibre chart. Its residual site gauges fix exactly the block
centres, and its contour coordinates are indexed by the
remaining sites. This residual action preserves the exact
constraint fibre. Let $\mathfrak h$ be its site gauge space,
$T:\mathfrak h\to E$ its tangent in the free chart, and
$\Gamma:E\to\mathfrak h$ the derivative of its contour
coordinate at a block-axial real reference. The contour's
centre endpoint is fixed and its other endpoint contributes
the inverse gauge factor. Differentiating this endpoint law gives
\begin{equation}\label{mb:gauge-tangent-readout}
                              \Gamma T=-I.
\end{equation}
This identity concerns the contour mean on its actual branch.
It does not require an inverse Hessian.

\mbneedspace{14\baselineskip}
\begin{lemma}[Gauge null splitting]\label{mb:fibre-gauge-null}
For the complete outer fibre and its action-preserving residual gauge action just specified, the complete fibre Hessian satisfies $K_AT=0$. Moreover,
\begin{equation}\label{mb:fibre-splitting}
 E=\ker\Gamma\oplus\operatorname{ran}T.
\end{equation}
For a positive gauge weight $a$, the form
$K_G=\Gamma^{\mathsf T}a\Gamma$ vanishes on $\ker\Gamma$.
The two summands in \eqref{mb:fibre-splitting} have zero mixed
terms for both $K_A$ and $K_G$.
\end{lemma}
\begin{proof}
The residual gauge acts on the exact constraint fibre. In its
free chart, differentiate the action-invariance identity once
in the gauge parameter and once in a free direction at zero.
The derivative of the gauge tangent is multiplied by
$D\mathcal S(0)=0$. The remaining term is $K_A[v,T\lambda]$,
which is therefore zero. The chart-curvature term in
\eqref{mb:fibre-hessian} is included in this differentiation.

Every $u$ decomposes as $(u+T\Gamma u)-T\Gamma u$.
By \eqref{mb:gauge-tangent-readout}, the first term lies in
$\ker\Gamma$. If $T\lambda\in\ker\Gamma$, then
$-\lambda=\Gamma T\lambda=0$, proving that the sum is direct.
The remaining assertions follow from the Gram form of $K_G$.
\end{proof}

The splitting is a decomposition for the quadratic forms; it
need not be orthogonal in the original free-coordinate Euclidean
metric. It proves no physical gap on $\ker\Gamma$ by itself.
That lower bound is supplied by the contour and outer-fibre
coercivity estimates. At the block-axial reference the contour
gauge coordinate vanishes, so the gauge Hessian is exactly
$\Gamma^{\mathsf T}a\Gamma$ without a missing second-chart term.


% Construction source: audit/gauge-determinant.tex
\mbneedspace{16\baselineskip}
\subsection{The two Gaussian powers of the action weight}

\begin{corollary}[Physical gauge splitting and the determinant]
\label{mb:cor-gauge-determinant}
Use the actual fibre and gauge decomposition of \cref{mb:fibre-gauge-null}. Let \(d=\dim E\), \(r=\dim\operatorname{ran}T\), and assume the complete physical precision \(K=K_A+K_G\) is positive. Retain the reference gauge weight and vary only the action weight, putting \(Q_\tau=\tau K_A+K_G\) for \(\operatorname{Re}\tau>0\). With branches continued from \(\tau=1\),
\begin{align}
 \det Q_\tau &= \tau^{d-r}\det K,\label{mb:eq-physical-determinant-power}\\
 \int_E e^{-u^TQ_\tau u/2}\,du
    &=\tau^{-(d-r)/2}\int_E e^{-u^TKu/2}\,du,\label{mb:eq-physical-gaussian-power}\\
 \int_E e^{(\tau-1)u^TK_Gu/2}\,d\mu_{(\tau K)^{-1}}(u)
    &=\tau^{r/2}.\label{mb:eq-physical-gauge-compensation}
\end{align}
Here \(\mu_{(\tau K)^{-1}}\) is the normalized Gaussian on the full free space, before insertion of fluctuation tests. The two logarithmic normalizations are \(-\frac d2\log\tau\) and \(\frac r2\log\tau\), with sum \(-\frac{d-r}{2}\log\tau\).
\end{corollary}

\begin{proof}
Choose a basis of \(\ker\Gamma\) followed by one of \(\operatorname{ran}T\), and denote the combined coordinate map by \(B\). In these coordinates \(B^T K_A B=\operatorname{diag}(A,0)\) and \(B^T K_G B=\operatorname{diag}(0,G)\). The complete positivity gives \(A>0\) and \(G>0\) on their respective blocks. Consequently \(B^TQ_\tau B=\operatorname{diag}(\tau A,G)\). The factor \((\det B)^2\) is the same in both determinants, proving \eqref{mb:eq-physical-determinant-power} even when the chosen splitting is oblique in the original coordinates.

The real part of \(Q_\tau\) is at least \(\min(1,\operatorname{Re}\tau)K\) as a real quadratic form. The Gaussian integrals are therefore holomorphic on the right half-plane by dominated integration on its compact subsets. Real Gaussian integration in the two sectors proves \eqref{mb:eq-physical-gaussian-power} for \(\tau>0\), and hence for the whole half-plane. Finally \(-\tau u^TKu/2+(\tau-1)u^TK_Gu/2=-u^TQ_\tau u/2\). The normalization of the precision \(\tau K\) contributes \(\tau^{d/2}\) against the reference integral. Inserting \eqref{mb:eq-physical-gaussian-power} proves \eqref{mb:eq-physical-gauge-compensation}. Empty blocks use the zero-dimensional integral and determinant equal to one.
\end{proof}

On a full periodic one-step fibre with $n$ coarse sites,
there are $4L^4n$ fine positive-bond occurrences and $4n$
pivots. Each carries three colour coordinates, so
\[
 d=12(L^4-1)n,\qquad r=3(L^4-1)n,\qquad d-r=9(L^4-1)n.
\]
The value of $r$ counts the noncentral sites, with injectivity
of their gauge tangents given by \eqref{mb:gauge-tangent-readout}.
Consequently the compensating expectation depends only on
these dimensions and is constant across the admitted full
physical fibres. On another domain the determinant formula
requires its actual splitting in \cref{mb:fibre-gauge-null};
the full-period counts do not apply to an arbitrary mask.

\clearpage
\begin{thebibliography}{9}
\raggedright
\bibitem{Averaging} \emph{Block averaging and gauge geometry for four-dimensional lattice Yang--Mills theory}, companion manuscript, 2026.
\bibitem{Bal83} T. Ba\l aban, \emph{Regularity and decay of lattice Green's functions}, Comm. Math. Phys. \textbf{89} (1983), 571--597. \href{https://doi.org/10.1007/BF01214744}{doi:10.1007/BF01214744}.
\bibitem{Bal84a} T. Ba\l aban, \emph{Propagators and renormalization transformations for lattice gauge theories. I}, Comm. Math. Phys. \textbf{95} (1984), 17--40. \href{https://doi.org/10.1007/BF01215753}{doi:10.1007/BF01215753}.
\bibitem{Bal84b} T. Ba\l aban, \emph{Propagators and renormalization transformations for lattice gauge theories. II}, Comm. Math. Phys. \textbf{96} (1984), 223--250. \href{https://doi.org/10.1007/BF01240221}{doi:10.1007/BF01240221}.
\bibitem{Bal85a} T. Ba\l aban, \emph{Averaging operations for lattice gauge theories}, Comm. Math. Phys. \textbf{98} (1985), 17--51. \href{https://doi.org/10.1007/BF01211042}{doi:10.1007/BF01211042}.
\bibitem{Bal85p} T. Ba\l aban, \emph{Propagators for lattice gauge theories in a background field}, Comm. Math. Phys. \textbf{99} (1985), 389--434. \href{https://doi.org/10.1007/BF01240355}{doi:10.1007/BF01240355}.
\end{thebibliography}
\end{document}
